\documentclass{ueap-sdm}

%------------------------------------------------
% PACOTES EXTRAS DO DOCUMENTO
%------------------------------------------------
\usepackage{float}
\usepackage{csquotes}
\usepackage{booktabs}
\usepackage{array}
\usepackage{graphicx}

%------------------------------------------------
% REFERÊNCIAS BIBLIOGRÁFICAS
%------------------------------------------------
\usepackage[
backend=biber,
style=authoryear,
natbib=true,
doi=true,
url=true,
isbn=false,
eprint=true,
hyperref=true,
maxcitenames=2,
maxbibnames=99
]{biblatex}

\DeclareFieldFormat{doi}{
	\href{https://doi.org/#1}
	{\textcolor{UEAPBlue}{doi:#1}}
}

\DeclareFieldFormat{url}{
	\href{#1}
	{\textcolor{UEAPBlue}{#1}}
}

\addbibresource{bibliografia/referencias_sdm.bib}

%------------------------------------------------
% LINKS E CORES
%------------------------------------------------
\hypersetup{
	colorlinks=true,
	linkcolor=UEAPDarkBlue,
	citecolor=UEAPDarkBlue,
	urlcolor=UEAPBlue
}

%------------------------------------------------
% METADADOS DO DOCUMENTO
%------------------------------------------------
\title{Modelagem de Distribuição de Espécies em R}
\subtitulo{Fundamentos Ecológicos, Aplicações em Conservação e Mudanças Climáticas}

\instituicao{Universidade do Estado do Amapá}
\curso{Colegiado de Engenharia Florestal}
\programa{Programa de Pós-graduação em Recursos Naturais Amazônicos -- RENAmazon}
\disciplina{Modelagem de Distribuição de Espécies em R}

\autor{Robson Borges de Lima}
\emailautor{robson.lima.ef@gmail.com}
\githubautor{https://github.com/robson-lima2023}
\orcidautor{0000-0001-5915-4045}

\logoueap{Brasao_UEAP.png}
\localano{Macapá, Amapá -- 2026}
\versao{Versão preliminar 1.0}

%------------------------------------------------
% INÍCIO DO DOCUMENTO
%------------------------------------------------
\begin{document}
	
	\begin{titlepage}
		\thispagestyle{empty}
		\newgeometry{margin=0cm}
		\noindent
		\includegraphics[
		width=\paperwidth,
		height=\paperheight,
		keepaspectratio=false
		]{figuras/capa/capa_frente_sdm.png}
		\restoregeometry
	\end{titlepage}
	
	\clearpage
	
	\folhaderosto
	\copyrightpage
	\sobreautor
	\paginainstitucional
	\prefacio
	
	\frontmatter
	\tableofcontents
	
	\cleardoublepage
	%================================================
% FLUXOGRAMA MESTRE DO LIVRO
%================================================
\cleardoublepage

\chapter*{Fluxograma Mestre}
\addcontentsline{toc}{chapter}{Fluxograma Mestre da Modelagem de Distribuição de Espécies}

\section*{Visão geral}

A Modelagem de Distribuição de Espécies (\textit{Species Distribution Modeling} -- SDM)
constitui uma estrutura integrada de inferência ecológica que combina
biogeografia, ecologia, climatologia, ciência de dados espaciais e modelagem
estatística para compreender e projetar a distribuição das espécies no espaço
e no tempo.

Embora frequentemente apresentada como uma sequência de algoritmos
computacionais, a SDM é fundamentalmente uma abordagem ecológica baseada na
relação entre organismos e ambiente. Cada etapa do processo influencia
diretamente a qualidade dos modelos, a interpretação dos resultados e a
robustez das inferências produzidas.

O fluxograma apresentado na Figura~\ref{fig:fluxograma-mestre-sdm}
resume o fluxo de trabalho adotado neste livro, desde a teoria ecológica
do nicho até aplicações em conservação da biodiversidade sob cenários de
mudanças climáticas.

\begin{figure}[p]
	\centering
	\includegraphics[
	width=0.95\textwidth,
	height=0.88\textheight,
	keepaspectratio
	]{figuras/fluxograma_mestre_sdm.png}
	
	\caption{
		Fluxograma mestre da Modelagem de Distribuição de Espécies (SDM)
		utilizado neste livro. O fluxo integra fundamentos ecológicos,
		dados de ocorrência, limpeza e rarefação espacial, definição da área acessível
		($M$), variáveis ambientais, multicolinearidade, background e pseudo-ausências,
		ajuste e avaliação de modelos, modelagem ensemble, projeções climáticas futuras,
		paleoclima, extrapolação, quantificação de incertezas e aplicações em conservação
		da biodiversidade. As cores representam grandes componentes conceituais do
		processo: fundamentos ecológicos, dados biológicos, preditores ambientais,
		modelagem, projeções climáticas, incertezas e aplicações em conservação.
	}
	\label{fig:fluxograma-mestre-sdm}
\end{figure}

\cleardoublepage

\section*{Como interpretar o fluxograma}

O processo inicia-se com a teoria ecológica do nicho, que fornece a base
conceitual para compreender a distribuição das espécies em função de gradientes
ambientais. Em seguida, registros de ocorrência são obtidos, validados e
submetidos a procedimentos de limpeza e rarefação espacial para reduzir erros
e vieses amostrais.

Posteriormente, define-se a área acessível ($M$) da espécie, conforme o
\textit{BAM framework}, e selecionam-se variáveis ambientais potencialmente
relevantes para explicar sua distribuição. Essas variáveis devem ser avaliadas
quanto à multicolinearidade por meio de correlação, VIF e outras abordagens de
seleção.

Após a construção da base de modelagem, são gerados pontos de background ou
pseudo-ausências e ajustados diferentes algoritmos, incluindo GLM, GAM,
Random Forest, Boosted Regression Trees e Maxent. Os modelos resultantes são
avaliados por métricas de desempenho e posteriormente integrados em modelos
ensemble.

Os modelos calibrados podem então ser projetados para cenários climáticos
futuros, períodos paleoclimáticos ou novas regiões geográficas. Nessas
situações, torna-se essencial avaliar a extrapolação ambiental utilizando
ferramentas como MESS e MOP.

Finalmente, diferentes fontes de incerteza são quantificadas e incorporadas
às análises, permitindo aplicações robustas em conservação, planejamento
ambiental, identificação de refúgios climáticos e priorização de áreas para
proteção da biodiversidade.

\begin{center}
	\fbox{
		\begin{minipage}{0.90\textwidth}
			
			\textbf{Síntese metodológica.}
			
			A qualidade de um modelo de distribuição não depende exclusivamente do
			algoritmo utilizado. Em muitos estudos, decisões tomadas durante a preparação
			dos dados, seleção de variáveis, definição da área acessível ($M$) e avaliação
			da extrapolação exercem impacto maior sobre os resultados finais do que a
			escolha do algoritmo propriamente dita.
			
		\end{minipage}
	}
\end{center}

\section*{Relação entre o fluxograma e as unidades do livro}

\begin{table}[H]
	\centering
	\caption{Correspondência entre as etapas do fluxo de trabalho e as unidades do livro.}
	\begin{tabular}{ll}
		\toprule
		Etapa do fluxograma & Unidade correspondente \\
		\midrule
		Teoria ecológica do nicho & Unidade 1 \\
		Ocorrências, limpeza e rarefação espacial & Unidade 2 \\
		Variáveis ambientais & Unidade 3 \\
		Multicolinearidade & Unidade 4 \\
		Área acessível ($M$) & Unidade 5 \\
		GLM & Unidade 6 \\
		GAM & Unidade 7 \\
		Random Forest & Unidade 8 \\
		BRT & Unidade 9 \\
		Maxent & Unidade 10 \\
		Ensemble & Unidade 11 \\
		Avaliação de modelos & Unidade 12 \\
		Projeções climáticas futuras & Unidade 13 \\
		Transferência e extrapolação & Unidade 14 \\
		Paleoclima & Unidade 15 \\
		Incertezas em SDM & Unidade 16 \\
		Aplicações à conservação & Unidade 17 \\
		Considerações finais & Unidade 18 \\
		\bottomrule
	\end{tabular}
\end{table}

\clearpage

	
	\mainmatter

%================================================
\chapter{Fundamentos ecológicos do nicho}
%================================================

\begin{objetivos}
Ao final desta unidade, o estudante deverá ser capaz de:
\begin{itemize}
  \item compreender a base ecológica da Modelagem de Distribuição de Espécies (SDM) e da Modelagem de Nicho Ecológico (ENM);
  \item diferenciar os conceitos de nicho de Grinnell, Elton e Hutchinson;
  \item distinguir nicho fundamental, nicho realizado, distribuição observada e distribuição potencial;
  \item reconhecer o papel dos fatores abióticos, bióticos e históricos na distribuição das espécies;
  \item interpretar criticamente mapas de adequabilidade ambiental e projeções sob mudanças climáticas;
  \item compreender por que a conservação de nicho é uma hipótese central em projeções espaciais e temporais.
\end{itemize}
\end{objetivos}

\section{Introdução}

A Modelagem de Distribuição de Espécies, conhecida internacionalmente como \textit{Species Distribution Modeling} (SDM), \textit{Habitat Suitability Modeling} (HSM) ou \textit{Ecological Niche Modeling} (ENM), constitui uma das abordagens mais importantes da ecologia quantitativa contemporânea. Seu objetivo central é relacionar registros de ocorrência de uma espécie com condições ambientais observadas no espaço geográfico, permitindo estimar áreas ambientalmente adequadas à sua presença atual, passada ou futura \parencite{guisan2000,guisan2005predicting,guisan2017habitat,elith2009,franklin2010}.

Apesar de frequentemente ser apresentada como uma técnica computacional baseada em mapas, algoritmos e camadas ambientais, a SDM é, antes de tudo, uma aplicação ecológica do conceito de nicho. A qualidade de um modelo de distribuição depende não apenas da escolha do algoritmo, mas da coerência entre pergunta ecológica, dados de ocorrência, variáveis ambientais, escala espacial, área acessível, hipóteses biogeográficas e interpretação da adequabilidade ambiental \parencite{peterson2011,guisan2017habitat,barve2011crucial}.

Nesta unidade são apresentados os fundamentos ecológicos necessários para compreender o que os modelos de distribuição realmente estimam. Serão discutidos os conceitos clássicos de nicho de Grinnell, Elton e Hutchinson; a distinção entre nicho fundamental e nicho realizado; a diferença entre distribuição potencial e distribuição observada; e o papel da conservação de nicho em estudos de mudanças climáticas, invasões biológicas e conservação da biodiversidade.

\begin{conceito}
A SDM não deve ser entendida apenas como uma técnica de produção de mapas. Ela é uma estrutura de inferência ecológica que conecta ocorrências biológicas, espaço ambiental, espaço geográfico, hipóteses de nicho e decisões de conservação.
\end{conceito}

\section{Por que começar pelo conceito de nicho?}

Um erro comum em aplicações de SDM é iniciar a modelagem diretamente pelo software, pelos pacotes em R ou pelos algoritmos, sem definir previamente qual dimensão ecológica está sendo representada. No entanto, todo modelo de distribuição carrega uma hipótese sobre por que uma espécie ocorre em determinados locais e está ausente de outros.

De forma geral, a distribuição geográfica de uma espécie resulta da interação entre pelo menos três componentes principais: condições abióticas adequadas, interações bióticas favoráveis ou toleráveis e possibilidade histórica ou atual de dispersão até a área considerada. Essa lógica está no centro do chamado \textit{BAM framework}, no qual $B$ representa fatores bióticos, $A$ representa fatores abióticos e $M$ representa a área acessível por dispersão \parencite{soberon2005,soberon2009,barve2011crucial}.

Assim, uma espécie somente será observada em um local quando as condições ambientais estiverem dentro de seus limites de tolerância, quando as interações ecológicas não impedirem sua persistência e quando a espécie tiver conseguido alcançar aquele local ao longo de sua história evolutiva, biogeográfica ou demográfica.

\begin{equationbox}
\[
G_0 \approx A \cap B \cap M
\]
Em termos conceituais, a distribuição ocupada observável ($G_0$) tende a ocorrer na interseção entre condições abióticas adequadas ($A$), condições bióticas compatíveis ($B$) e área acessível à espécie ($M$).
\end{equationbox}

\begin{atencao}
Um modelo calibrado fora da área acessível ecologicamente plausível pode confundir ausência de dispersão com inadequação ambiental. Por isso, a delimitação da área de calibração é uma decisão ecológica, não apenas cartográfica.
\end{atencao}

\section{O nicho de Grinnell}

O conceito grinnelliano de nicho está associado principalmente aos requisitos ambientais que permitem a existência de uma espécie em determinada região. Nesse sentido, o nicho é entendido como o conjunto de condições físicas e climáticas necessárias para que uma população sobreviva e se mantenha ao longo do tempo \parencite{grinnell1917,soberon2007}.

Esse conceito é particularmente importante para SDM porque muitos modelos utilizam variáveis climáticas, topográficas e edáficas como preditoras da distribuição das espécies. Quando usamos temperatura média anual, precipitação, sazonalidade climática, altitude, declividade, textura do solo ou déficit hídrico climático, estamos trabalhando predominantemente com uma interpretação grinnelliana do nicho.

Em escalas amplas, como continentes, biomas ou grandes regiões hidroclimáticas, fatores climáticos tendem a exercer forte controle sobre os limites de distribuição das espécies. Por isso, modelos baseados em variáveis bioclimáticas são frequentemente aplicados para prever distribuição potencial, mudanças futuras de área adequada e deslocamentos de faixa climática sob cenários de mudanças climáticas \parencite{guisan2005predicting,elith2009,guisan2017habitat}.

\begin{dica}[title={\faLightbulb\quad dica}]
Em SDM, variáveis grinnellianas são especialmente úteis quando a pergunta envolve macroecologia, biogeografia, mudanças climáticas, invasões biológicas ou conservação em grandes escalas espaciais.
\end{dica}

\section{O nicho de Elton}

O conceito eltoniano de nicho enfatiza o papel funcional da espécie na comunidade ecológica. Diferentemente da visão grinnelliana, centrada nos requisitos ambientais, o nicho eltoniano destaca o que a espécie faz no ecossistema: suas interações tróficas, seus recursos alimentares, seus predadores, seus competidores, seus mutualistas e seus efeitos sobre outras espécies \parencite{elton1927,soberon2007}.

Em termos práticos, o nicho de Elton está relacionado à posição da espécie na rede de interações ecológicas. Para plantas, isso pode incluir relações com polinizadores, dispersores, fungos micorrízicos, herbívoros, patógenos, competidores por luz, água e nutrientes. Para animais, pode envolver dieta, comportamento de forrageamento, predação, parasitismo e relações de dependência com habitats ou espécies-chave.

Embora essencial para a compreensão ecológica completa da distribuição das espécies, o nicho eltoniano é mais difícil de incorporar em SDM tradicionais. Isso ocorre porque dados espacialmente explícitos sobre interações bióticas raramente estão disponíveis na mesma extensão, resolução e qualidade dos dados climáticos. Por isso, muitos modelos acabam representando principalmente o componente abiótico da distribuição, enquanto as interações bióticas permanecem implícitas, ausentes ou apenas parcialmente representadas.

\begin{atencao}[title={\faExclamationTriangle\quad atenção}]
A ausência de variáveis bióticas em um SDM não significa que interações ecológicas sejam irrelevantes. Significa apenas que, em muitos casos, elas não foram medidas ou não estão disponíveis em resolução compatível com a análise.
\end{atencao}

\section{O nicho de Hutchinson}

O conceito de nicho de Hutchinson representa uma das formulações mais influentes da ecologia moderna. Hutchinson definiu o nicho como um hipervolume $n$-dimensional, no qual cada dimensão representa uma variável ambiental ou ecológica relevante para a persistência da espécie. Dentro desse hipervolume, as condições permitem crescimento populacional positivo; fora dele, a população não consegue se manter ao longo do tempo \parencite{hutchinson1957,colwell2009}.

Essa definição é fundamental para SDM porque permite separar o espaço geográfico do espaço ambiental. O espaço geográfico corresponde às coordenadas, mapas e regiões onde a espécie pode ocorrer. O espaço ambiental corresponde aos valores das variáveis que caracterizam os locais, como temperatura, precipitação, altitude, solo, sazonalidade e disponibilidade hídrica.

A modelagem de distribuição opera justamente nessa relação entre espaço ambiental e espaço geográfico. Primeiro, os registros de ocorrência são associados às condições ambientais dos locais onde a espécie foi observada. Em seguida, o modelo estima a resposta da espécie nesse espaço ambiental. Por fim, essa resposta é projetada de volta para o espaço geográfico, gerando mapas de adequabilidade ambiental \parencite{guisan2017habitat,raes2018modeling}.

\begin{conceito}
O espaço geográfico é o mapa: coordenadas, pixels, polígonos, bacias, biomas e unidades de conservação. O espaço ambiental é o conjunto de condições associadas a cada local: temperatura, precipitação, solo, altitude, sazonalidade, déficit hídrico e outras variáveis. A SDM aprende padrões no espaço ambiental e projeta esses padrões no espaço geográfico.
\end{conceito}

\section{Nicho fundamental e nicho realizado}

O nicho fundamental corresponde ao conjunto de condições ambientais sob as quais uma espécie poderia manter populações viáveis na ausência de restrições bióticas negativas e barreiras de dispersão. Em outras palavras, representa o espaço de tolerância fisiológica e ecológica potencial da espécie \parencite{hutchinson1957,pulliam2000,soberon2005}.

O nicho realizado é a porção do nicho fundamental efetivamente ocupada pela espécie sob condições naturais. Ele resulta da combinação entre tolerâncias ambientais, interações bióticas, dispersão, história evolutiva, filtros regionais e disponibilidade real de ambientes \parencite{pulliam2000,soberon2009}.

Na prática, o nicho fundamental raramente é estimado diretamente por SDM correlativos baseados em ocorrências naturais. Isso ocorre porque os registros de campo já refletem restrições de dispersão, história biogeográfica, interações ecológicas, filtros ambientais atuais e limitações amostrais. Assim, modelos baseados apenas em ocorrências geralmente estimam uma aproximação do nicho realizado ou de parte dele, não o nicho fundamental completo.

\begin{atencao}
Modelos correlativos inferem associações entre ocorrência e ambiente. Eles não medem diretamente taxa de crescimento populacional, sobrevivência, fecundidade, mortalidade, condutância estomática, vulnerabilidade hidráulica ou tolerância térmica. Para aproximar o nicho fundamental, são necessários dados experimentais, fisiológicos, demográficos ou modelos mecanísticos.
\end{atencao}

\section{Distribuição observada, real e potencial}

A distribuição observada corresponde aos registros conhecidos da espécie. Esses registros podem vir de herbários, museus, inventários florestais, bases de dados de biodiversidade, literatura científica, coleções biológicas, levantamentos de campo ou monitoramentos sistemáticos. Embora sejam a base dos modelos, os dados de ocorrência raramente representam perfeitamente a distribuição real da espécie \parencite{hijmans2021,guisan2017habitat}.

A distribuição real corresponde ao conjunto de locais efetivamente ocupados pela espécie em determinado período, independentemente de terem sido amostrados. Já a distribuição potencial representa o conjunto de áreas que apresentam condições ambientais adequadas para a espécie, de acordo com o modelo ajustado.

\begin{conceito}
A distribuição observada é uma amostra dos registros conhecidos; a distribuição real é a ocupação efetiva da espécie na paisagem; e a distribuição potencial é a projeção espacial das condições estimadas como adequadas pelo modelo.
\end{conceito}

Em espécies amazônicas, esse problema é especialmente relevante, pois grandes regiões permanecem subamostradas, o acesso é difícil e muitos registros estão concentrados próximos a rios, estradas, cidades, unidades de pesquisa e áreas historicamente mais inventariadas. Assim, ausência de registros não deve ser automaticamente interpretada como ausência ecológica.

\begin{dica}
Prefira interpretar mapas contínuos de adequabilidade como gradientes de suporte ambiental ao modelo, e não como mapas absolutos de presença e ausência. A transformação em mapas binários pode ser útil para algumas aplicações, mas sempre depende da escolha de limiares, que introduzem novas decisões metodológicas e incertezas.
\end{dica}

\section{Conservação de nicho}

A conservação de nicho refere-se à tendência de uma espécie ou linhagem manter, ao longo do tempo evolutivo ou entre regiões, requisitos ambientais semelhantes aos de seus ancestrais ou populações de origem. Esse conceito é central em estudos de mudanças climáticas, invasões biológicas, biogeografia histórica e projeções futuras de distribuição \parencite{peterson2011,guisan2017habitat,raes2018modeling}.

Quando projetamos um modelo calibrado no presente para o futuro, assumimos que a relação entre a espécie e o ambiente será suficientemente conservada para permitir inferência. Da mesma forma, quando usamos ocorrências da área nativa para prever o potencial invasivo em outra região, assumimos que a espécie manterá parte importante de seu nicho climático.

No entanto, a conservação de nicho não deve ser tratada como regra absoluta. Nichos podem se expandir, contrair ou mudar entre regiões e períodos. Mudanças aparentes podem resultar de evolução adaptativa, plasticidade fenotípica, liberação de inimigos naturais, novas interações ecológicas, ausência de competidores, mudanças na disponibilidade ambiental ou diferenças na amostragem.

\begin{atencao}
Projetar não é prever com certeza. Projeções futuras de SDM representam cenários condicionais: se a relação estimada entre ocorrência e ambiente permanecer válida, se as variáveis escolhidas forem relevantes, se o cenário climático ocorrer, e se dispersão e interações não impedirem a ocupação, então determinadas áreas poderão se tornar mais ou menos adequadas.
\end{atencao}

\section{Relação entre nicho e mudanças climáticas}

As mudanças climáticas modificam a distribuição espacial das condições ambientais que compõem o nicho das espécies. À medida que temperatura, precipitação, sazonalidade, déficit hídrico e frequência de extremos se alteram, áreas atualmente adequadas podem se tornar inadequadas, enquanto novas áreas podem se tornar ambientalmente compatíveis.

Em regiões tropicais, especialmente na Amazônia, mudanças no regime de chuvas, aumento da temperatura, intensificação do déficit de pressão de vapor, secas extremas e eventos associados ao El Niño podem alterar profundamente os filtros ambientais que estruturam a distribuição das espécies arbóreas. Para árvores gigantes, essas mudanças podem ser ainda mais críticas, pois indivíduos de grande porte dependem de estabilidade hidráulica, solos profundos, acesso contínuo à água e resistência mecânica a eventos extremos.

\begin{aplicacao}
Em estudos sobre árvores gigantes da Amazônia, a distribuição atual de indivíduos adultos pode refletir décadas ou séculos de condições ambientais favoráveis. Portanto, modelos baseados apenas em presença de adultos podem capturar a adequabilidade histórica para sobrevivência, mas não necessariamente a adequabilidade atual para regeneração. Uma abordagem robusta deve considerar, sempre que possível, diferentes estágios ontogenéticos, estrutura florestal, clima extremo, solo, hidrologia e estabilidade da paisagem.
\end{aplicacao}

\section{Síntese conceitual da unidade}

A modelagem de distribuição de espécies é uma ferramenta poderosa, mas sua interpretação depende de fundamentos ecológicos sólidos. O nicho de Grinnell orienta a relação entre espécie e ambiente físico; o nicho de Elton destaca o papel das interações ecológicas; e o nicho de Hutchinson fornece a base formal para compreender o nicho como um hipervolume multidimensional.

O nicho fundamental representa o conjunto de condições potencialmente toleradas pela espécie, enquanto o nicho realizado corresponde à porção efetivamente ocupada sob restrições bióticas, históricas, ambientais e de dispersão. A distribuição observada é uma amostra geográfica imperfeita dos locais conhecidos de ocorrência, enquanto a distribuição potencial representa a projeção espacial das condições estimadas como adequadas pelo modelo.

\begin{leitura}
Para aprofundamento conceitual, recomenda-se iniciar por \textcite{guisan2005predicting}, \textcite{soberon2007}, \textcite{soberon2009}, \textcite{peterson2011} e \textcite{guisan2017habitat}. Esses trabalhos formam a base conceitual moderna para interpretar nichos ecológicos, áreas de distribuição, adequabilidade ambiental e projeções espaço-temporais.
\end{leitura}

\subsection{Estudo de caso real: nicho ecológico climático de \textit{Dinizia excelsa}}

Nesta etapa, a simulação didática inicial é substituída por um estudo de caso real com registros de ocorrência de \textit{Dinizia excelsa} no bioma Amazônia. A espécie é uma árvore amazônica de grande porte, com importância ecológica, estrutural e conservacionista, sendo adequada para demonstrar como dados reais de ocorrência podem ser integrados a variáveis ambientais em uma análise inicial de nicho ecológico.

O objetivo desta unidade não é ainda ajustar um modelo preditivo completo, mas representar o espaço ambiental ocupado pela espécie em relação ao conjunto de condições climáticas disponíveis no bioma Amazônia. Para isso, os registros corrigidos de ocorrência foram filtrados espacialmente pelo limite do bioma e associados a variáveis bioclimáticas atuais derivadas do WorldClim.

Foram utilizadas quatro variáveis ambientais selecionadas: temperatura média anual, sazonalidade da temperatura, precipitação anual e sazonalidade da precipitação. Essas variáveis representam gradientes climáticos amplos associados ao nicho grinnelliano da espécie, isto é, ao conjunto de condições ambientais que podem favorecer, limitar ou restringir sua ocorrência geográfica.

\rscriptfile{scripts/unidade01/01_simulacao_nicho_ecologico.R}

\begin{figure}[H]
	\centering
	\includegraphics[width=0.96\textwidth]{figuras/unidade01/unidade01_mapas_variaveis_patchwork_dinizia.png}
	\caption{Variáveis ambientais selecionadas no bioma Amazônia com registros corrigidos de \textit{Dinizia excelsa}. Cada painel possui escala própria de cores, permitindo visualizar adequadamente a distribuição espacial de temperatura média anual, sazonalidade da temperatura, precipitação anual e sazonalidade da precipitação.}
	\label{fig:unidade01_variaveis_patchwork_dinizia}
\end{figure}

\begin{figure}[H]
	\centering
	\includegraphics[width=0.82\textwidth]{figuras/unidade01/unidade01_pca_nicho_climatico_dinizia.png}
	\caption{Nicho climático multivariado de \textit{Dinizia excelsa} representado por PCA das variáveis bioclimáticas selecionadas. Os pontos cinza representam o background ambiental do bioma Amazônia e os pontos vermelhos representam as ocorrências da espécie.}
	\label{fig:unidade01_pca_dinizia}
\end{figure}

\begin{figure}[H]
	\centering
	\includegraphics[width=0.82\textwidth]{figuras/unidade01/unidade01_nicho_temperatura_precipitacao_dinizia.png}
	\caption{Espaço climático ocupado por \textit{Dinizia excelsa} em relação ao background ambiental do bioma Amazônia, considerando temperatura média anual e precipitação anual.}
	\label{fig:unidade01_nicho_tp_dinizia}
\end{figure}

\begin{aplicacao}
	As Figuras~\ref{fig:unidade01_pca_dinizia} e~\ref{fig:unidade01_nicho_tp_dinizia} demonstram que a distribuição geográfica de \textit{Dinizia excelsa} está associada a um subconjunto específico das condições climáticas disponíveis no bioma Amazônia. O gráfico bivariado evidencia a ocupação do espaço climático definido pela temperatura média anual e pela precipitação anual, enquanto a PCA representa a posição da espécie em um espaço ambiental multidimensional construído a partir de múltiplas variáveis bioclimáticas. Observa-se que as ocorrências concentram-se em regiões ambientais relativamente mais restritas que o conjunto total de condições disponíveis no bioma, sugerindo processos de filtragem ecológica relacionados às exigências climáticas da espécie. Em estudos de Modelagem de Distribuição de Espécies, essas relações constituem a base para a estimativa do nicho ecológico e para a projeção de áreas potencialmente adequadas sob condições ambientais atuais, passadas ou futuras.
\end{aplicacao}

\begin{exercicios}
\begin{enumerate}
  \item Explique a diferença entre nicho de Grinnell, nicho de Elton e nicho de Hutchinson.
  \item Por que modelos de distribuição baseados em dados de ocorrência geralmente estimam uma aproximação do nicho realizado, e não do nicho fundamental?
  \item Diferencie distribuição observada, distribuição real e distribuição potencial.
  \item Explique por que ausência de registro em bases de biodiversidade não deve ser automaticamente interpretada como ausência ecológica da espécie.
  \item Discuta a importância da conservação de nicho para projeções climáticas futuras.
  \item Em espécies arbóreas amazônicas de grande porte, quais fatores podem fazer com que a distribuição de adultos não represente adequadamente a distribuição potencial de regenerantes?
\end{enumerate}
\end{exercicios}

% ============================================================
% UNIDADE 2
% ============================================================

\chapter[Dados de ocorrência]{Dados de ocorrência, limpeza espacial e preparação de registros}

\begin{objetivos}
	Ao final desta unidade, o estudante deverá ser capaz de compreender as principais fontes de dados de ocorrência utilizadas em Modelagem de Distribuição de Espécies; reconhecer limitações associadas a dados de herbários, inventários e repositórios globais; padronizar registros taxonômicos e geográficos; filtrar ocorrências pelo limite do bioma Amazônia; remover coordenadas problemáticas, duplicatas e registros espacialmente redundantes; aplicar rarefação espacial; extrair variáveis ambientais associadas às ocorrências; e preparar uma base final de registros para as unidades posteriores de modelagem.
\end{objetivos}

\section{Introdução}

Dados de ocorrência constituem a base empírica da Modelagem de Distribuição de Espécies. Em modelos correlativos, os registros geográficos da espécie são associados às condições ambientais dos locais onde ela foi observada, permitindo estimar relações entre ocorrência e ambiente. Entretanto, a qualidade dos modelos depende diretamente da qualidade desses registros.

Em aplicações reais, os dados de ocorrência podem ser obtidos de inventários florestais, herbários, coleções biológicas, plataformas digitais, literatura científica e repositórios globais como o GBIF. Esses dados frequentemente combinam diferentes níveis de precisão taxonômica, geográfica e temporal. Por isso, antes da modelagem, é necessário realizar uma etapa cuidadosa de padronização, limpeza e filtragem espacial \parencite{guisan2017habitat,hijmans2021species,zizka2019coordinatecleaner}.

\begin{conceito}[title={\faBook\quad Registro de ocorrência}]
	Um registro de ocorrência é uma evidência espacial da presença de uma espécie em determinado local. Em SDM, esse registro normalmente é representado por coordenadas geográficas associadas ao nome da espécie, fonte do dado e, quando disponível, data, instituição, coletor e precisão espacial.
\end{conceito}

\section{Fontes de dados de ocorrência}

Os dados usados em SDM podem vir de diferentes fontes. Inventários florestais geralmente apresentam maior controle taxonômico e espacial, mas tendem a cobrir áreas limitadas. Herbários e coleções biológicas fornecem registros históricos valiosos, porém podem conter coordenadas imprecisas ou ausentes. Plataformas globais, como GBIF, ampliam a cobertura espacial, mas exigem filtros rigorosos de qualidade.

No estudo de referência usado nesta apostila, os registros de \textit{Dinizia excelsa} foram integrados a dados de inventário e GBIF, e posteriormente submetidos a filtros espaciais, rarefação e extração de variáveis bioclimáticas para análise de nicho climático na escala amazônica \parencite{delima2026warmer}. Esse fluxo é semelhante ao que será implementado nesta unidade.

\begin{dica}[title={\faLightbulb\quad Integração de fontes}]
	Uma boa base de ocorrência pode combinar inventários, herbários e GBIF, desde que todos os registros passem por padronização, limpeza e controle de duplicatas.
\end{dica}

\section{Problemas comuns em dados de ocorrência}

Os problemas mais frequentes incluem nomes científicos inconsistentes, coordenadas inválidas, registros no oceano, pontos em capitais, coordenadas de instituições, duplicatas, inversão de latitude e longitude, erros de georreferenciamento e forte viés espacial de amostragem.

Em espécies amazônicas, é comum haver concentração de registros próximos a rios, estradas, cidades, herbários, áreas de inventário e regiões de maior acesso logístico. Esse padrão não representa necessariamente maior adequabilidade ambiental, mas sim maior esforço de coleta.

\begin{atencao}[title={\faExclamationTriangle\quad Ocorrência não é ausência}]
	Registros de presença indicam onde a espécie foi observada. A ausência de registros em uma região não significa necessariamente ausência da espécie; pode refletir ausência de amostragem.
\end{atencao}

\section{Fluxo de preparação dos dados}

O fluxo usado nesta unidade segue a lógica geral de estudos de SDM em escala amazônica: compilação dos registros, padronização de colunas, filtragem por coordenadas válidas, correção espacial, recorte pelo bioma, remoção de duplicatas, rarefação espacial, extração ambiental e exportação da base final.

\rscriptfile{scripts/unidade02/01_compilar_ocorrencias_dinizia.R}

\section{Limpeza geográfica e recorte pelo bioma Amazônia}

Após a compilação, os registros são convertidos para objetos espaciais, avaliados quanto à validade das coordenadas e filtrados pelo limite do bioma Amazônia. Essa etapa garante que a base de trabalho represente a escala geográfica definida para o estudo.

\rscriptfile{scripts/unidade02/02_limpeza_espacial_dinizia.R}

\begin{figure}[H]
	\centering
	\includegraphics[width=0.92\textwidth]{figuras/unidade02/unidade02_mapa_bruto_limpo_dinizia.png}
	\caption{Comparação entre registros brutos e registros filtrados de \textit{Dinizia excelsa} no bioma Amazônia.}
	\label{fig:unidade02_bruto_limpo_dinizia}
\end{figure}

\section{Extração ambiental e duplicatas por célula raster}

Mesmo após remover coordenadas idênticas, diferentes registros podem cair no mesmo pixel climático. Como os modelos usam valores ambientais extraídos de rasters, múltiplos pontos no mesmo pixel podem representar pseudorreplicação ambiental. Por isso, é recomendável remover duplicatas por célula raster, mantendo apenas um registro por pixel ambiental.

\rscriptfile{scripts/unidade02/03_extrair_ambiente_dinizia.R}

\section{Rarefação espacial}

A rarefação espacial reduz a concentração excessiva de registros em áreas muito amostradas. Nesta unidade, aplicamos uma distância mínima entre ocorrências, permitindo reduzir autocorrelação espacial e viés de amostragem. Esse procedimento é amplamente usado em estudos de nicho ecológico e pode ser implementado com o pacote \texttt{spThin} \parencite{aiellolammens2015spthin}.

\rscriptfile{scripts/unidade02/04_rarefacao_espacial_dinizia.R}

\begin{figure}[H]
	\centering
	\includegraphics[width=0.92\textwidth]{figuras/unidade02/unidade02_rarefacao_dinizia.png}
	\caption{Efeito da rarefação espacial nos registros de \textit{Dinizia excelsa}. O procedimento reduz a redundância espacial e a concentração de pontos próximos.}
	\label{fig:unidade02_rarefacao_dinizia}
\end{figure}

\section{Diagnóstico ambiental dos registros}

Após a limpeza e rarefação, as ocorrências finais são associadas às variáveis ambientais selecionadas. Essa etapa permite verificar a amplitude climática ocupada pela espécie, identificar valores extremos e preparar a matriz final para modelagem.

\rscriptfile{scripts/unidade02/05_diagnostico_ambiental_dinizia.R}

\begin{figure}[H]
	\centering
	\includegraphics[width=0.92\textwidth]{figuras/unidade02/unidade02_histogramas_variaveis_dinizia.png}
	\caption{Distribuição das variáveis bioclimáticas extraídas nos registros finais de \textit{Dinizia excelsa}.}
	\label{fig:unidade02_hist_variaveis_dinizia}
\end{figure}

\begin{figure}[H]
	\centering
	\includegraphics[width=0.86\textwidth]{figuras/unidade02/unidade02_espaco_climatico_dinizia.png}
	\caption{Espaço climático ocupado por \textit{Dinizia excelsa} em relação ao background ambiental do bioma Amazônia.}
	\label{fig:unidade02_espaco_climatico_dinizia}
\end{figure}

\section{Fluxograma do processamento}

O fluxograma resume a sequência de decisões adotadas para transformar registros brutos em uma base final de modelagem.

\rscriptfile{scripts/unidade02/06_fluxograma_preparacao_dinizia.R}

\begin{figure}[H]
	\centering
	\includegraphics[width=0.92\textwidth]{figuras/unidade02/unidade02_fluxograma_preparacao_dinizia.png}
	\caption{Fluxograma de preparação de dados de ocorrência para SDM usando \textit{Dinizia excelsa} como estudo de caso.}
	\label{fig:unidade02_fluxograma_dinizia}
\end{figure}

\section{Pipeline completo da Unidade 2}

\rscriptfile{scripts/unidade02/07_pipeline_unidade02_dinizia.R}

\section{Síntese da unidade}

A preparação dos dados de ocorrência é uma das etapas mais importantes da SDM. Modelos sofisticados não corrigem dados ruins. A qualidade final da modelagem depende da coerência taxonômica, validade espacial, controle de duplicatas, redução de viés amostral e compatibilidade entre registros e variáveis ambientais.

\begin{conceito}[title={\faBook\quad Mensagem central}]
	Antes de ajustar qualquer algoritmo, é necessário transformar registros brutos em uma base de ocorrência limpa, espacialmente coerente, ambientalmente consistente e compatível com a escala do estudo.
\end{conceito}

\section{Exercícios de fixação}

\begin{enumerate}
	\item Explique por que registros de GBIF e herbários precisam ser limpos antes da modelagem.
	\item Diferencie duplicata geográfica e duplicata por célula raster.
	\item Por que registros fora do bioma de estudo devem ser removidos?
	\item Explique a finalidade da rarefação espacial.
	\item O que é viés de amostragem em SDM?
	\item Por que a ausência de registros não deve ser interpretada automaticamente como ausência da espécie?
	\item Proponha um fluxo de limpeza para uma espécie arbórea amazônica rara.
\end{enumerate}

% ============================================================
% UNIDADE 3
% ============================================================

\chapter[Variáveis ambientais]{Variáveis Ambientais em Modelagem de Distribuição de Espécies}

\begin{objetivos}
	Ao final desta unidade, o estudante deverá ser capaz de compreender o papel das variáveis ambientais em SDM; reconhecer diferenças entre variáveis climáticas, topográficas, edáficas e índices derivados; acessar bases como WorldClim, CHELSA, EarthEnv, SoilGrids e ENVIREM; compreender resolução espacial, extensão geográfica, CRS, reamostragem, recorte e harmonização de camadas; construir pilhas ambientais reprodutíveis em R; e aplicar esses conceitos a um estudo de caso real com \textit{Dinizia excelsa} no bioma Amazônia.
\end{objetivos}

\section{Introdução às variáveis ambientais em SDM}

As variáveis ambientais constituem a base explicativa da Modelagem de Distribuição de Espécies. Elas representam gradientes climáticos, edáficos, topográficos, hidrológicos ou estruturais usados para caracterizar os locais onde uma espécie ocorre e os locais disponíveis para calibração e projeção do modelo. Assim, a qualidade ecológica de um SDM depende diretamente da coerência entre registros de ocorrência, área de calibração, escala espacial e conjunto de variáveis selecionadas \parencite{guisan2005predicting,elith2009,guisan2017habitat}.

Em modelos correlativos, as variáveis ambientais funcionam como preditoras. Elas são usadas para estimar relações estatísticas entre ocorrência da espécie e condições ambientais observadas. No entanto, essas relações não devem ser interpretadas automaticamente como relações causais. Uma variável pode ser importante porque representa diretamente um processo ecológico, mas também pode atuar apenas como substituta de outro gradiente não medido \parencite{franklin2010,guisan2017habitat}.

\begin{conceito}[title={\faBook\quad Variável ambiental não é necessariamente mecanismo ecológico}]
	Uma variável ambiental representa um gradiente associado à distribuição da espécie. O fato de uma variável apresentar alta importância estatística em um modelo não implica necessariamente relação causal direta. A interpretação ecológica deve considerar a história natural da espécie, a escala espacial, a resolução dos dados, o conjunto de variáveis disponíveis e o conhecimento prévio sobre os processos que limitam sua distribuição.
\end{conceito}

\subsection{Espaço ambiental, resolução e escala}

O espaço ambiental corresponde ao conjunto de condições representadas pelas variáveis ambientais usadas no modelo. Enquanto o espaço geográfico é formado por coordenadas, pixels, polígonos e mapas, o espaço ambiental é formado pelos valores das variáveis associadas a cada local. Dois pontos distantes no mapa podem ser ambientalmente semelhantes, enquanto pontos próximos podem ser muito diferentes ao longo de gradientes de altitude, solo, umidade ou cobertura vegetal.

A resolução espacial define o tamanho da unidade mínima de análise. Em dados raster, corresponde ao tamanho do pixel. A extensão geográfica define a área total considerada. A escala ecológica envolve simultaneamente resolução, extensão, tempo e nível de organização biológica. Em ampla escala, o clima tende a atuar como filtro dominante; em escala regional ou local, solo, relevo, hidrologia, cobertura da terra e microclima podem ser determinantes.

\begin{atencao}[title={\faExclamationTriangle\quad Compatibilidade entre ocorrência e ambiente}]
	A resolução das camadas ambientais precisa ser compatível com a precisão dos registros de ocorrência. Um registro com incerteza espacial de vários quilômetros não deve ser interpretado como representação precisa de um pixel ambiental de poucos metros.
\end{atencao}

\section{WorldClim}

O WorldClim é uma das bases climáticas mais utilizadas em SDM. Ele disponibiliza camadas globais de temperatura, precipitação e variáveis bioclimáticas derivadas, com diferentes resoluções espaciais. As variáveis BIOCLIM sintetizam aspectos médios, extremos e sazonais do clima, sendo amplamente usadas em estudos de nicho ecológico, biogeografia e mudanças climáticas \parencite{fick2017worldclim}.

\begin{conceito}[title={\faBook\quad Variáveis BIOCLIM}]
	As variáveis BIOCLIM são descritores climáticos derivados de temperatura e precipitação mensal. Elas incluem temperatura média anual, sazonalidade térmica, temperatura dos trimestres mais quente e mais frio, precipitação anual, sazonalidade da precipitação e precipitação dos trimestres mais úmido, seco, quente e frio.
\end{conceito}

\section{CHELSA}

O CHELSA é uma base climática global de alta resolução desenvolvida para representar padrões climáticos considerando processos orográficos e atmosféricos. Em regiões montanhosas ou com forte heterogeneidade topográfica, o CHELSA pode representar gradientes de precipitação e temperatura de forma mais adequada que bases interpoladas mais simples \parencite{karger2017chelsa}.

\begin{dica}[title={\faLightbulb\quad WorldClim ou CHELSA?}]
	WorldClim é amplamente usado e de fácil acesso. CHELSA pode ser particularmente útil em regiões montanhosas, áreas com forte variação topográfica ou estudos sensíveis à precipitação orográfica. A escolha deve considerar região, escala, objetivo e consistência com projeções futuras.
\end{dica}

\section{EarthEnv}

O EarthEnv reúne variáveis ambientais globais voltadas à caracterização da heterogeneidade ambiental, incluindo produtos derivados de topografia, clima, cobertura e habitat. Em SDM, variáveis de heterogeneidade podem auxiliar em estudos de endemismo, diversidade beta, refúgios ambientais e diferenciação de habitats \parencite{amatulli2018data}.

\section{ENVIREM}

O ENVIREM amplia o conjunto clássico de variáveis bioclimáticas ao incluir índices derivados com maior interpretação ecofisiológica, como variáveis relacionadas à evapotranspiração potencial, umidade climática, continentalidade e balanço água-energia. Essas variáveis podem melhorar a representação ambiental em modelos de nicho ecológico \parencite{title2018envirem}.

\section{Variáveis topográficas}

Variáveis topográficas derivadas de modelos digitais de elevação podem representar gradientes importantes em SDM. Elevação, declividade, orientação de vertente, rugosidade, índice de posição topográfica e heterogeneidade do relevo influenciam microclima, drenagem, exposição solar, erosão, acúmulo de água e distribuição de habitats \parencite{farr2007srtm}.

Em ampla escala, como no bioma Amazônia, a topografia pode parecer menos variável do que em regiões montanhosas, mas ainda pode representar gradientes associados a drenagem, altitude relativa, áreas de terra firme, várzeas, escarpas, interflúvios e condições hidrológicas locais.

\section{Variáveis edáficas e SoilGrids}

Variáveis edáficas são fundamentais para representar filtros ambientais associados à fertilidade, textura, acidez, capacidade de troca catiônica, retenção de água e disponibilidade de nutrientes. O SoilGrids, desenvolvido pelo ISRIC, fornece predições globais de propriedades do solo em múltiplas profundidades, sendo útil para aplicações regionais e globais em SDM \parencite{poggio2021soilgrids}.

\begin{conceito}[title={\faBook\quad Solo como filtro ecológico}]
	O solo representa um componente relativamente estável da paisagem e pode explicar padrões de distribuição que não são capturados apenas por clima. Em plantas, variáveis como pH, textura, carbono orgânico, argila, areia e capacidade de troca catiônica podem estar fortemente associadas à ocorrência.
\end{conceito}

\section{Harmonização de camadas ambientais}

Bases ambientais frequentemente possuem resoluções, extensões, sistemas de coordenadas e unidades diferentes. Antes de usar essas camadas em SDM, é necessário harmonizá-las. Esse processo envolve recorte, reprojeção, reamostragem, padronização de nomes, remoção de camadas redundantes e construção de uma pilha ambiental consistente.

\begin{atencao}[title={\faExclamationTriangle\quad Harmonização antes da modelagem}]
	Nunca combine diretamente camadas ambientais com resoluções, extensões ou projeções diferentes. A inconsistência espacial entre rasters pode gerar extrações erradas de valores ambientais e comprometer todo o modelo.
\end{atencao}

\section{Estudo de caso real: variáveis ambientais para \textit{Dinizia excelsa}}

Nesta seção, os conceitos apresentados anteriormente são aplicados a um estudo de caso real com \textit{Dinizia excelsa} no bioma Amazônia. Os registros de ocorrência foram previamente compilados, limpos, filtrados espacialmente e rarefeitos na Unidade 2. Agora, esses registros são associados a variáveis ambientais reais, recortadas e harmonizadas para a escala do bioma Amazônia.

O objetivo não é apenas gerar mapas ambientais, mas preparar uma base espacial coerente para as etapas seguintes da apostila: controle de multicolinearidade, seleção de variáveis, ajuste de modelos, avaliação, ensemble, projeções climáticas futuras e reconstruções paleoclimáticas.

\begin{aplicacao}
	No estudo de caso com \textit{Dinizia excelsa}, as variáveis ambientais são tratadas como hipóteses ecológicas sobre os gradientes que podem influenciar a distribuição da espécie. Variáveis climáticas representam filtros energéticos e hídricos; variáveis topográficas representam heterogeneidade do terreno; e variáveis edáficas, quando disponíveis, representam filtros associados ao substrato e à fertilidade.
\end{aplicacao}

\subsection{Variáveis bioclimáticas selecionadas}

Para o estudo de caso, foram preparadas variáveis bioclimáticas derivadas do WorldClim e recortadas para o bioma Amazônia. As variáveis selecionadas representam gradientes de temperatura, precipitação e sazonalidade:

\begin{itemize}
	\item BIO3 -- Isotermalidade;
	\item BIO4 -- Sazonalidade da temperatura;
	\item BIO5 -- Temperatura máxima do mês mais quente;
	\item BIO8 -- Temperatura média do trimestre mais úmido;
	\item BIO12 -- Precipitação anual;
	\item BIO13 -- Precipitação do mês mais úmido;
	\item BIO14 -- Precipitação do mês mais seco;
	\item BIO18 -- Precipitação do trimestre mais quente;
	\item BIO19 -- Precipitação do trimestre mais frio.
\end{itemize}

Essas variáveis são úteis para representar componentes climáticos associados à disponibilidade hídrica, estresse térmico, sazonalidade e extremos ambientais. Em espécies arbóreas tropicais, tais gradientes podem estar associados à sobrevivência, crescimento, regeneração, recrutamento e persistência populacional.

\rscriptfile{scripts/unidade03/01_variaveis_ambientais_dinizia.R}

\begin{figure}[H]
	\centering
	\includegraphics[width=0.98\textwidth]{figuras/unidade03/unidade03_variaveis_bioclimaticas_dinizia.png}
	\caption{Variáveis bioclimáticas selecionadas para o bioma Amazônia com registros finais de \textit{Dinizia excelsa}. Cada painel possui escala própria de cores, permitindo visualizar adequadamente gradientes de temperatura, precipitação e sazonalidade.}
	\label{fig:unidade03_variaveis_bioclimaticas_dinizia}
\end{figure}

\subsection{Variáveis topográficas}

Além das variáveis climáticas, o script prepara variáveis topográficas derivadas de um modelo digital de elevação. Nesta unidade, foram calculadas elevação, declividade e orientação de vertente. Essas camadas são harmonizadas com as variáveis bioclimáticas, garantindo mesma resolução, extensão, máscara espacial e sistema de referência.

\begin{figure}[H]
	\centering
	\includegraphics[width=0.98\textwidth]{figuras/unidade03/unidade03_variaveis_topograficas_dinizia.png}
	\caption{Variáveis topográficas preparadas para o bioma Amazônia e associadas ao estudo de caso com \textit{Dinizia excelsa}. Os mapas representam elevação, declividade e orientação de vertente.}
	\label{fig:unidade03_topografia_dinizia}
\end{figure}

\subsection{Extração ambiental nas ocorrências}

Após a preparação das camadas ambientais, os valores de cada variável são extraídos nos pontos finais de ocorrência. Essa matriz ambiental é essencial para as próximas etapas da modelagem, pois permite analisar a amplitude ambiental ocupada pela espécie, avaliar correlações entre variáveis, selecionar preditores e construir modelos de adequabilidade.

\begin{figure}[H]
	\centering
	\includegraphics[width=0.94\textwidth]{figuras/unidade03/unidade03_histogramas_ambientais_dinizia.png}
	\caption{Distribuição das variáveis ambientais extraídas nos registros finais de \textit{Dinizia excelsa}. Os histogramas permitem avaliar a amplitude ambiental ocupada pela espécie em relação às variáveis selecionadas.}
	\label{fig:unidade03_histogramas_dinizia}
\end{figure}

\section{Produtos gerados pela Unidade 3}

O script desta unidade gera os seguintes produtos principais:

\begin{itemize}
	\item raster multivariado com variáveis bioclimáticas selecionadas para o bioma Amazônia;
	\item raster com variáveis topográficas harmonizadas;
	\item pilha ambiental completa para uso nas próximas unidades;
	\item tabela com valores ambientais extraídos nas ocorrências de \textit{Dinizia excelsa};
	\item background ambiental do bioma Amazônia;
	\item metadados das camadas ambientais;
	\item mapas bioclimáticos, mapas topográficos e histogramas ambientais.
\end{itemize}

\begin{dica}[title={\faLightbulb\quad Organização dos arquivos}]
	Os arquivos gerados nesta unidade devem ser mantidos como referência para as unidades seguintes. A pilha ambiental salva em \texttt{dados/unidade03/processados/variaveis\_ambientais\_amazonia\_unidade03.tif} será utilizada em análises de multicolinearidade, seleção de variáveis, modelagem e projeções.
\end{dica}

\section{Síntese da unidade}

As variáveis ambientais representam a ponte entre o nicho ecológico e a modelagem espacial. A seleção adequada dessas variáveis exige conhecimento ecológico, coerência com a escala da pergunta, compatibilidade com os registros de ocorrência e cuidado técnico no processamento espacial. Bases como WorldClim, CHELSA, EarthEnv, ENVIREM, SoilGrids e SRTM são extremamente úteis, mas devem ser usadas criticamente.

No estudo de caso com \textit{Dinizia excelsa}, a preparação das camadas ambientais para o bioma Amazônia permitiu construir uma base espacial reprodutível para as próximas etapas da apostila. A partir deste ponto, os dados já estão organizados para avaliar multicolinearidade, selecionar variáveis, ajustar modelos e produzir mapas de adequabilidade atual, passada e futura.

\begin{conceito}[title={\faBook\quad Mensagem central}]
	Em SDM, variáveis ambientais não são apenas camadas de mapa. Elas representam hipóteses ecológicas sobre os fatores que limitam, favorecem ou estruturam a distribuição das espécies.
\end{conceito}

\section{Exercícios de fixação}

\begin{enumerate}
	\item Explique a diferença entre espaço geográfico e espaço ambiental.
	\item Por que uma variável ambiental importante no modelo não representa necessariamente um mecanismo ecológico causal?
	\item Diferencie resolução espacial, extensão geográfica e escala ecológica.
	\item Explique por que variáveis climáticas são amplamente usadas em SDM de ampla escala.
	\item Cite exemplos de variáveis topográficas importantes para distribuição de espécies.
	\item Explique a importância de variáveis edáficas em modelos de distribuição de plantas.
	\item Por que é necessário harmonizar camadas ambientais antes da modelagem?
	\item No estudo de caso com \textit{Dinizia excelsa}, quais variáveis ambientais poderiam representar disponibilidade hídrica?
	\item Por que os rasters precisam ter a mesma resolução, extensão e CRS antes da modelagem?
\end{enumerate}

% ============================================================
% UNIDADE 4
% ============================================================

\chapter[Multicolinearidade]{Multicolinearidade em Modelagem de Distribuição de Espécies}

\begin{objetivos}
	Ao final desta unidade, o estudante deverá ser capaz de compreender o conceito de multicolinearidade em SDM; identificar variáveis ambientais redundantes por correlação; calcular e interpretar o Fator de Inflação da Variância (VIF); compreender os efeitos da colinearidade sobre interpretação ecológica, estabilidade dos parâmetros e transferibilidade dos modelos; aplicar PCA como estratégia de redução dimensional; e construir um fluxo reprodutível em R para seleção e transformação de variáveis ambientais antes da modelagem.
\end{objetivos}

\section{Introdução}

A multicolinearidade ocorre quando duas ou mais variáveis preditoras carregam informação semelhante. Em Modelagem de Distribuição de Espécies, esse problema é comum porque variáveis climáticas, topográficas, edáficas e hidrológicas frequentemente covariam no espaço. Por exemplo, temperatura média anual pode estar fortemente associada à altitude; precipitação anual pode estar correlacionada com precipitação do trimestre mais úmido; e variáveis topográficas derivadas de um mesmo modelo digital de elevação podem apresentar forte redundância.

A presença de variáveis altamente correlacionadas pode gerar modelos instáveis, dificultar a interpretação da importância das variáveis, inflar incertezas em modelos paramétricos e reduzir a capacidade de transferência espacial ou temporal. Por isso, a seleção de preditores ambientais deve combinar fundamentação ecológica, análise exploratória e critérios estatísticos \parencite{guisan2017habitat,naimi2016sdm,raes2018modeling}.

\begin{atencao}[title={\faExclamationTriangle\quad Mais variáveis não significam melhor modelo}]
	Adicionar muitas variáveis ambientais não garante melhor desempenho. Variáveis redundantes podem aumentar complexidade, gerar instabilidade, dificultar interpretação ecológica e reduzir a transferibilidade do modelo.
\end{atencao}

\section{O problema da multicolinearidade}

Variáveis ambientais frequentemente apresentam forte correlação entre si. Em regiões tropicais, por exemplo, áreas com elevada precipitação anual frequentemente apresentam baixa sazonalidade hídrica, enquanto gradientes térmicos podem estar associados a gradientes altitudinais.

Quando duas ou mais variáveis carregam informação semelhante, ocorre multicolinearidade. Embora a capacidade preditiva do modelo possa permanecer elevada, a interpretação ecológica dos coeficientes torna-se instável e a transferência espacial ou temporal pode ser comprometida \parencite{dormann2013collinearity}.

\section{Correlação entre variáveis ambientais}

A forma mais simples de diagnosticar redundância entre preditores é calcular correlações par a par. Em SDM, os coeficientes de Pearson e Spearman são os mais utilizados. Pearson mede associação linear entre variáveis contínuas, enquanto Spearman avalia associação monotônica baseada em postos, sendo útil quando as distribuições não são normais ou apresentam assimetria.

Na prática, muitos estudos usam limiares como \(|r| > 0{,}7\) ou \(|r| > 0{,}8\) para indicar alta correlação. Esses limiares não devem ser tratados como regras absolutas, mas como critérios operacionais. Quando duas variáveis estão fortemente correlacionadas, recomenda-se manter aquela com maior relevância ecológica, maior qualidade espacial, maior interpretabilidade ou melhor adequação ao objetivo do estudo \parencite{dormann2013collinearity,raes2018modeling}.

\begin{conceito}[title={\faBook\quad Correlação não é causalidade}]
	Duas variáveis podem ser altamente correlacionadas sem que uma cause a outra. Em SDM, a correlação indica redundância informacional, mas a decisão sobre qual variável manter deve considerar a ecologia da espécie e a hipótese do estudo.
\end{conceito}

\begin{conceito}[title={\faBook\quad Correlação não implica redundância ecológica}]
	Duas variáveis podem apresentar alta correlação estatística e, ainda assim, representar processos ecológicos distintos.
	
	Por exemplo, precipitação anual e precipitação do trimestre mais seco podem ser altamente correlacionadas, mas representam mecanismos ecológicos diferentes relacionados à disponibilidade hídrica.
\end{conceito}

\rscriptfile{scripts/unidade04/01_estrutura_pastas.R}

\rscriptfile{scripts/unidade04/02_preparar_variaveis_colineares.R}

\rscriptfile{scripts/unidade04/03_correlacao_variaveis.R}

\begin{figure}[H]
	\centering
	\includegraphics[width=\textwidth]
	{figuras/unidade04/unidade04_matriz_correlacao.png}
	\caption{
		Matriz de correlação de Pearson entre as variáveis ambientais
		utilizadas para modelagem de \textit{Dinizia excelsa}
		no bioma Amazônia.
		Os coeficientes indicam a intensidade da associação linear
		entre pares de variáveis ambientais.
	}
	\label{fig:unidade04_matriz_correlacao}
\end{figure}

\section{Fator de Inflação da Variância}

O Fator de Inflação da Variância, conhecido como VIF, avalia o quanto a variância estimada de um coeficiente é inflada pela dependência linear entre uma variável preditora e todas as demais. Diferentemente da correlação par a par, o VIF considera a relação de cada variável com o conjunto completo dos preditores.

\begin{equationbox}
	\[
	VIF_j = \frac{1}{1 - R_j^2}
	\]
	
	Em que \(VIF_j\) é o Fator de Inflação da Variância da variável \(j\), e \(R_j^2\) é o coeficiente de determinação obtido ao ajustar uma regressão da variável \(j\) contra todas as demais variáveis preditoras.
\end{equationbox}

Valores elevados de VIF indicam que uma variável pode ser explicada por uma combinação das outras. Como regra prática, VIF acima de 10 é frequentemente interpretado como evidência de colinearidade severa, embora limiares mais conservadores, como 5, também sejam usados em estudos ecológicos \parencite{obrien2007vif,dormann2013collinearity,naimi2016sdm}.

\begin{dica}[title={\faLightbulb\quad Correlação e VIF são complementares}]
	A correlação identifica pares de variáveis redundantes. O VIF identifica variáveis redundantes em relação ao conjunto completo de preditores. Em estudos robustos, os dois diagnósticos devem ser avaliados conjuntamente.
\end{dica}

\rscriptfile{scripts/unidade04/04_vif_variaveis.R}

\begin{figure}[H]
	\centering
	\includegraphics[width=0.95\textwidth]
	{figuras/unidade04/unidade04_corrplot_correlacao.png}
	\caption{
		Representação gráfica da correlação entre variáveis ambientais.
		Cores mais intensas representam correlações mais fortes,
		enquanto correlações próximas de zero indicam maior independência
		entre variáveis.
	}
	\label{fig:unidade04_corrplot}
\end{figure}

\begin{figure}[H]
	\centering
	\includegraphics[width=0.85\textwidth]
	{figuras/unidade04/unidade04_vif_variaveis.png}
	\caption{
		Valores do Fator de Inflação da Variância (VIF)
		para as variáveis ambientais avaliadas.
		A linha tracejada representa o limiar adotado
		para identificação de multicolinearidade excessiva.
	}
	\label{fig:unidade04_vif}
\end{figure}

\begin{aplicacao}
	As Figuras~\ref{fig:unidade04_matriz_correlacao}
	e~\ref{fig:unidade04_corrplot}
	permitem identificar conjuntos de variáveis que carregam
	informação semelhante.
	
	Em SDM, manter simultaneamente variáveis altamente correlacionadas
	pode dificultar a interpretação ecológica dos modelos e aumentar
	a instabilidade das projeções espaciais e temporais.
\end{aplicacao}

\begin{aplicacao}
	O VIF quantifica o quanto a variância de uma variável
	é inflacionada pela presença de correlação com outras variáveis.
	
	Valores elevados indicam redundância estatística
	e sugerem que a variável deve ser removida ou substituída
	por outra ecologicamente mais representativa.
\end{aplicacao}

\section{Seleção de variáveis ambientais}

A seleção de variáveis deve equilibrar três dimensões: relevância ecológica, qualidade dos dados e redundância estatística. Remover automaticamente variáveis apenas com base em limiares pode gerar modelos pobres do ponto de vista ecológico. Por outro lado, manter variáveis redundantes pode gerar modelos difíceis de interpretar.

Em SDM, recomenda-se iniciar com um conjunto amplo de variáveis candidatas, justificar ecologicamente cada grupo de preditores e, em seguida, aplicar filtros de correlação e VIF. Quando duas variáveis são correlacionadas, a decisão deve priorizar aquela com relação mais direta com processos ecológicos, maior estabilidade temporal, melhor resolução ou menor erro conhecido.

\begin{atencao}[title={\faExclamationTriangle\quad Não selecione variáveis apenas pelo algoritmo}]
	A seleção automática pode remover variáveis ecologicamente importantes e manter variáveis estatisticamente convenientes, mas biologicamente frágeis. A seleção de preditores deve ser documentada e justificada.
\end{atencao}

\rscriptfile{scripts/unidade04/05_selecao_variaveis.R}

\begin{figure}[H]
	\centering
	\includegraphics[width=0.85\textwidth]
	{figuras/unidade04/unidade04_variaveis_selecionadas.png}
	\caption{
		Conjunto final de variáveis ambientais selecionadas
		após análise de correlação, VIF e interpretação ecológica.
		Essas variáveis serão utilizadas nas unidades subsequentes
		para ajuste dos modelos de distribuição de espécies.
	}
	\label{fig:unidade04_variaveis_final}
\end{figure}

\section{PCA aplicada a SDM}

A Análise de Componentes Principais (PCA) é uma estratégia de redução dimensional que transforma variáveis correlacionadas em novos eixos ortogonais. Esses eixos, chamados componentes principais, são combinações lineares das variáveis originais e não são correlacionados entre si.

Em SDM, a PCA pode ser útil quando existe forte colinearidade entre variáveis ambientais e o objetivo é representar o espaço ambiental de forma compacta. Entretanto, a PCA reduz a interpretabilidade ecológica direta, pois os componentes representam combinações de variáveis, não gradientes simples como temperatura ou precipitação.

\begin{conceito}[title={\faBook\quad PCA como ortogonalização}]
	A PCA não escolhe variáveis. Ela transforma um conjunto de variáveis correlacionadas em eixos independentes. Isso reduz colinearidade, mas exige cuidado na interpretação ecológica dos componentes.
\end{conceito}

\rscriptfile{scripts/unidade04/06_pca_variaveis.R}

\begin{figure}[H]
	\centering
	\includegraphics[width=0.90\textwidth]
	{figuras/unidade04/unidade04_pca_biplot_variaveis.png}
	\caption{
		Análise de Componentes Principais (PCA)
		das variáveis ambientais utilizadas no estudo de caso.
		Os vetores representam a contribuição de cada variável
		para os principais gradientes ambientais do bioma Amazônia.
	}
	\label{fig:unidade04_pca_biplot}
\end{figure}

\begin{figure}[H]
	\centering
	\includegraphics[width=0.80\textwidth]
	{figuras/unidade04/unidade04_pca_variancia.png}
	\caption{
		Proporção da variância explicada pelos componentes principais.
		As barras representam a contribuição individual de cada eixo,
		enquanto a linha representa a variância acumulada.
	}
	\label{fig:unidade04_pca_variancia}
\end{figure}

\begin{aplicacao}
	As Figuras~\ref{fig:unidade04_pca_biplot}
	e~\ref{fig:unidade04_pca_variancia}
	mostram que grande parte da estrutura ambiental
	do bioma Amazônia pode ser representada por poucos
	eixos ambientais principais.
	
	A PCA não substitui a seleção ecológica de variáveis,
	mas auxilia na identificação de redundâncias e na interpretação
	dos gradientes ambientais que estruturam a distribuição
	de \textit{Dinizia excelsa}.
\end{aplicacao}

\section{Pipeline integrador}

O script final reúne as etapas da Unidade 4: simulação de variáveis ambientais, análise de correlação, diagnóstico por VIF, seleção de variáveis e PCA.

\rscriptfile{scripts/unidade04/08_pipeline_multicolinearidade.R}

\section{Síntese da unidade}

A multicolinearidade é um problema central em SDM porque afeta interpretação, estabilidade e transferibilidade dos modelos. Correlação, VIF e PCA são ferramentas complementares. A correlação identifica redundância par a par; o VIF avalia redundância multivariada; e a PCA transforma variáveis correlacionadas em eixos independentes. A escolha entre remover variáveis ou usar PCA depende da pergunta científica, do algoritmo, da necessidade de interpretação ecológica e do objetivo de projeção espacial ou temporal.

\begin{conceito}[title={\faBook\quad Mensagem central da Unidade 4}]
	A seleção de variáveis não é apenas um procedimento estatístico.
	O objetivo é construir um conjunto de preditores que represente
	os principais filtros ecológicos associados à distribuição da espécie,
	evitando redundâncias excessivas e aumentando a interpretabilidade
	dos modelos.
	
	As variáveis selecionadas nesta unidade serão utilizadas
	nas Unidades 5 a 17 para modelagem, avaliação,
	projeções climáticas e aplicações à conservação.
\end{conceito}

\begin{aplicacao}
	As variáveis selecionadas nesta unidade serão utilizadas nas etapas subsequentes de modelagem. A partir deste ponto, o conjunto ambiental passa a representar uma hipótese ecológica refinada sobre os fatores que influenciam a distribuição de \textit{Dinizia excelsa} na Amazônia.
	
	Na próxima unidade, essas variáveis serão integradas à definição da área acessível (M) dentro do framework BAM.
\end{aplicacao}

\section{Exercícios de fixação}

\begin{enumerate}
	\item Explique o que é multicolinearidade e por que ela é comum em variáveis ambientais.
	\item Diferencie correlação de Pearson e correlação de Spearman.
	\item Por que duas variáveis altamente correlacionadas não devem ser mantidas automaticamente no mesmo modelo?
	\item Explique o significado do VIF.
	\item Qual é a diferença entre diagnóstico por correlação e diagnóstico por VIF?
	\item Em que situações a PCA pode ser útil em SDM?
	\item Discuta uma desvantagem ecológica de usar componentes principais como preditores.
\end{enumerate}

% ============================================================
% UNIDADE 5
% ============================================================

\chapter[Área acessível M]{Área Acessível (M), BAM Framework e Seleção de Background}

\begin{objetivos}
	Ao final desta unidade, o estudante deverá compreender o conceito de área acessível (M); interpretar o framework BAM; diferenciar área de estudo, área de calibração e área de projeção; delimitar uma área M biologicamente defensável; gerar pontos de background; comparar background no bioma inteiro versus background restrito à área M; e preparar dados de presença-background para modelagem de \textit{Dinizia excelsa}.
\end{objetivos}

\section{Introdução}

A definição da área acessível, frequentemente representada pela letra \(M\), é uma das decisões mais importantes em Modelagem de Distribuição de Espécies. A área \(M\) representa a porção do espaço geográfico que esteve historicamente acessível à espécie por dispersão, história biogeográfica, conectividade e ausência de barreiras intransponíveis \parencite{soberon2005,soberon2007,soberon2009,barve2011crucial}.

Em SDM, a área \(M\) influencia diretamente a seleção de background, a disponibilidade ambiental, os contrastes entre presença e ausência/background e, consequentemente, os parâmetros estimados pelos modelos. Uma área \(M\) muito ampla pode incluir ambientes nunca acessíveis à espécie; uma área muito restrita pode excluir gradientes relevantes para estimar seu nicho.

\begin{conceito}[title={\faBook\quad Área acessível (M)}]
	A área acessível (M) é o conjunto de regiões que a espécie poderia ter alcançado ao longo de sua história recente, considerando dispersão, barreiras, conectividade da paisagem e história biogeográfica. Ela não é necessariamente igual ao bioma inteiro nem à distribuição observada da espécie.
\end{conceito}

\section{BAM framework}

O framework BAM organiza a distribuição das espécies como resultado da interação entre três componentes principais: fatores bióticos (B), fatores abióticos (A) e área acessível por dispersão (M). Em termos simplificados, uma espécie ocorre onde as condições ambientais são adequadas, as interações bióticas permitem sua persistência e a região é acessível historicamente.

\begin{equationbox}
	\[
	G_O = A \cap B \cap M
	\]
	
	em que \(G_O\) representa a distribuição observada; \(A\), as condições abióticas adequadas; \(B\), as interações bióticas permissivas; e \(M\), a área acessível.
\end{equationbox}

\begin{atencao}[title={\faExclamationTriangle\quad Área M não é detalhe técnico}]
	A definição de \(M\) altera o conjunto de condições ambientais disponíveis para o modelo. Portanto, ela deve ser justificada biologicamente e não escolhida apenas por conveniência computacional.
\end{atencao}

\section{Delimitação da área M para \textit{Dinizia excelsa}}

Nesta unidade, a área \(M\) será delimitada de forma operacional a partir dos registros finais de \textit{Dinizia excelsa}, usando um buffer espacial ao redor das ocorrências e restringindo o resultado ao limite do bioma Amazônia. Essa abordagem é didática e reprodutível, adequada para construir a base das próximas unidades de modelagem.

\rscriptfile{scripts/unidade05/01_estrutura_pastas.R}

\rscriptfile{scripts/unidade05/02_delimitar_area_m_dinizia.R}

\begin{figure}[H]
	\centering
	\includegraphics[width=0.90\textwidth]{figuras/unidade05/unidade05_area_m_dinizia.png}
	\caption{Área acessível (M) delimitada para \textit{Dinizia excelsa} no bioma Amazônia. A área foi definida por buffer espacial ao redor das ocorrências e recortada pelo limite do bioma.}
	\label{fig:unidade05_area_m_dinizia}
\end{figure}

\section{Seleção de background}

Pontos de background representam o ambiente disponível para calibração do modelo. Em modelos presença-background, eles não são ausências verdadeiras. Eles caracterizam as condições ambientais disponíveis dentro da área considerada acessível.

\rscriptfile{scripts/unidade05/03_gerar_background_area_m.R}

\begin{figure}[H]
	\centering
	\includegraphics[width=0.90\textwidth]{figuras/unidade05/unidade05_background_area_m.png}
	\caption{Pontos de background gerados dentro da área acessível (M) para \textit{Dinizia excelsa}.}
	\label{fig:unidade05_background_m}
\end{figure}

\section{Background no bioma versus background na área M}

A seleção do background pode alterar substancialmente a distribuição ambiental disponível para o modelo. Por isso, é útil comparar o background amostrado no bioma inteiro com o background restrito à área \(M\).

\rscriptfile{scripts/unidade05/04_comparar_background_bioma_m.R}

\begin{figure}[H]
	\centering
	\includegraphics[width=0.90\textwidth]{figuras/unidade05/unidade05_comparacao_background.png}
	\caption{Comparação entre background amostrado no bioma Amazônia e background restrito à área acessível (M).}
	\label{fig:unidade05_comparacao_background}
\end{figure}

\section{Extração ambiental e base presença-background}

A última etapa consiste em extrair os valores das variáveis ambientais selecionadas nas presenças e nos pontos de background da área \(M\), construindo a base presença-background que será usada nas unidades seguintes.

\rscriptfile{scripts/unidade05/05_preparar_presenca_background.R}

\begin{figure}[H]
	\centering
	\includegraphics[width=0.86\textwidth]{figuras/unidade05/unidade05_espaco_ambiental_presenca_background.png}
	\caption{Espaço ambiental comparando presenças de \textit{Dinizia excelsa} e background dentro da área acessível (M).}
	\label{fig:unidade05_espaco_presenca_background}
\end{figure}

\section{Pipeline completo}

\rscriptfile{scripts/unidade05/06_pipeline_area_m_background.R}

\section{Síntese da unidade}

A definição da área acessível (M) é uma decisão ecológica, biogeográfica e metodológica. Ela influencia o background, a calibração, a avaliação e as projeções dos modelos. Para \textit{Dinizia excelsa}, a delimitação operacional de \(M\) no bioma Amazônia fornece uma base coerente para os modelos das próximas unidades.

\begin{conceito}[title={\faBook\quad Mensagem central}]
	Em SDM, o background deve representar o ambiente disponível para a espécie dentro da área acessível. Modelos calibrados com background inadequado podem produzir mapas biologicamente frágeis, mesmo quando apresentam boas métricas estatísticas.
\end{conceito}

% ============================================================
% UNIDADE 6
% ============================================================

\chapter[GLM]{Modelos Lineares Generalizados}

\begin{objetivos}
	Ao final desta unidade, o estudante deverá compreender os fundamentos dos Modelos Lineares Generalizados (GLM) aplicados à Modelagem de Distribuição de Espécies; ajustar modelos presença-background com família binomial; interpretar coeficientes, termos lineares e quadráticos; avaliar desempenho preditivo por AUC e TSS; gerar curvas de resposta parcial; projetar adequabilidade ambiental no espaço geográfico; e produzir mapas de adequabilidade atual para \textit{Dinizia excelsa} no bioma Amazônia.
\end{objetivos}

\section{Introdução}

Os Modelos Lineares Generalizados, conhecidos como GLM, representam uma das abordagens estatísticas mais clássicas e interpretáveis em Modelagem de Distribuição de Espécies. Em SDM, eles são frequentemente usados para relacionar registros de presença e ausência, ou presença e background, com variáveis ambientais preditoras \parencite{guisan2000,guisan2005predicting,guisan2017habitat}.

Em modelos presença-background, a resposta é binária: presença recebe valor 1 e background recebe valor 0. Embora pontos de background não sejam ausências verdadeiras, eles representam o conjunto de condições ambientais disponíveis dentro da área de calibração. Assim, o GLM estima a probabilidade relativa de ocorrência ou adequabilidade ambiental em função dos gradientes ambientais.

\begin{conceito}[title={\faBook\quad GLM em SDM}]
	Em SDM, um GLM binomial estima a relação entre a ocorrência da espécie e variáveis ambientais, usando uma função de ligação logit para transformar uma combinação linear de preditores em valores entre 0 e 1.
\end{conceito}

\section{Formulação matemática}

No caso binomial, seja \(Y_i\) uma variável resposta indicando presença ou background:

\[
Y_i =
\begin{cases}
	1, & \text{presença}\\
	0, & \text{background}
\end{cases}
\]

O GLM com ligação logit é definido por:

\begin{equationbox}
	\[
	\text{logit}(p_i)=
	\log\left(\frac{p_i}{1-p_i}\right)
	=
	\beta_0+\beta_1X_{1i}+\beta_2X_{2i}+...+\beta_kX_{ki}
	\]
	\[
	p_i=
	\frac{1}{1+\exp(-\eta_i)}
	\]
\end{equationbox}

em que \(p_i\) representa a adequabilidade estimada no pixel \(i\), \(X_k\) representa as variáveis ambientais e \(\beta_k\) representa os coeficientes estimados.

\begin{dica}[title={\faLightbulb\quad Termos quadráticos}]
	Em SDM, termos quadráticos são úteis porque permitem respostas unimodais, isto é, situações nas quais a espécie apresenta maior adequabilidade em valores intermediários de uma variável ambiental.
\end{dica}

\section{Preparação dos dados para GLM}

Nesta unidade, utilizamos a base presença-background construída na Unidade 5. As presenças correspondem aos registros finais de \textit{Dinizia excelsa}, enquanto os pontos de background foram amostrados dentro da área acessível \(M\). Essa base representa o conjunto de calibração dos modelos.

\rscriptfile{scripts/unidade06/01_estrutura_pastas.R}

\rscriptfile{scripts/unidade06/02_preparar_dados_glm.R}

\section{Ajuste do GLM}

O GLM é ajustado com família binomial e ligação logit. Para aumentar a flexibilidade ecológica, o modelo inclui termos lineares e quadráticos das variáveis selecionadas. Antes do ajuste, as variáveis ambientais são padronizadas para média zero e desvio-padrão um, reduzindo problemas numéricos e facilitando a comparação entre coeficientes.

\rscriptfile{scripts/unidade06/03_ajustar_glm_dinizia.R}

\begin{figure}[H]
	\centering
	\includegraphics[width=0.88\textwidth]{figuras/unidade06/unidade06_coeficientes_glm.png}
	\caption{Coeficientes estimados pelo GLM para \textit{Dinizia excelsa}. Valores positivos indicam associação positiva com a adequabilidade relativa; valores negativos indicam associação negativa, considerando as demais variáveis do modelo.}
	\label{fig:unidade06_coeficientes_glm}
\end{figure}

\section{Avaliação do modelo}

A avaliação do GLM é feita usando AUC, TSS, sensibilidade, especificidade e matriz de confusão. AUC resume a capacidade de discriminação do modelo ao longo de diferentes limiares; TSS combina sensibilidade e especificidade em um único índice. Como trabalhamos com presença-background, essas métricas devem ser interpretadas como indicadores comparativos, não como validação absoluta da distribuição real da espécie.

\rscriptfile{scripts/unidade06/04_avaliar_glm_dinizia.R}

\begin{figure}[H]
	\centering
	\includegraphics[width=0.82\textwidth]{figuras/unidade06/unidade06_avaliacao_glm_patchwork.png}
	\caption{Curva ROC do GLM ajustado para \textit{Dinizia excelsa}.}
	\label{fig:unidade06_roc_glm}
\end{figure}

\section{Curvas de resposta parcial}

As curvas de resposta parcial mostram como a adequabilidade estimada varia ao longo de cada variável ambiental, mantendo as demais constantes. Elas são essenciais para interpretar ecologicamente o comportamento do modelo.

\rscriptfile{scripts/unidade06/06_curvas_resposta_glm.R}

\begin{figure}[H]
	\centering
	\includegraphics[width=0.96\textwidth]{figuras/unidade06/unidade06_curvas_resposta_glm.png}
	\caption{Curvas de resposta parcial do GLM para as variáveis ambientais selecionadas.}
	\label{fig:unidade06_curvas_resposta}
\end{figure}

\section{Projeção espacial da adequabilidade atual}

Após o ajuste e avaliação, o GLM é projetado para a pilha ambiental da área de estudo. O resultado é um mapa contínuo de adequabilidade ambiental atual para \textit{Dinizia excelsa}.

\rscriptfile{scripts/unidade06/05_predicao_espacial_glm.R}

\begin{figure}[H]
	\centering
	\includegraphics[width=0.90\textwidth]{figuras/unidade06/unidade06_adequabilidade_glm.png}
	\caption{Mapa de adequabilidade ambiental atual de \textit{Dinizia excelsa} estimado por GLM.}
	\label{fig:unidade06_adequabilidade_glm}
\end{figure}

\section{Pipeline completo}

\rscriptfile{scripts/unidade06/07_pipeline_glm.R}

\section{Síntese da unidade}

O GLM é um modelo simples, interpretável e robusto para iniciar análises de SDM. Embora menos flexível que GAM, Random Forest, BRT ou Maxent, ele permite compreender explicitamente a relação entre ocorrência e gradientes ambientais. Por isso, é uma ferramenta fundamental em materiais didáticos, análises exploratórias e estudos que priorizam transparência e interpretação ecológica.

\begin{conceito}[title={\faBook\quad Mensagem central}]
	O GLM é uma ponte entre estatística clássica e SDM moderna: simples o suficiente para interpretação ecológica, mas poderoso o suficiente para produzir mapas de adequabilidade ambiental quando bem calibrado.
\end{conceito}

\section{Exercícios de fixação}

\begin{enumerate}
	\item Explique por que a família binomial é apropriada para modelos presença-background.
	\item O que representa a função logit em um GLM?
	\item Por que termos quadráticos são úteis em SDM?
	\item Diferencie coeficiente positivo e coeficiente negativo em um GLM.
	\item Explique a diferença entre AUC e TSS.
	\item Por que curvas de resposta parcial são importantes para interpretação ecológica?
	\item Quais limitações devem ser consideradas ao usar GLM em projeções espaciais?
\end{enumerate}

% ============================================================
% UNIDADE 7
% ============================================================

\chapter[GAM]{Modelos Aditivos Generalizados}

\begin{objetivos}
	Ao final desta unidade, o estudante deverá compreender a lógica dos Modelos Aditivos Generalizados (GAM) em Modelagem de Distribuição de Espécies; ajustar modelos presença-background com suavizadores; interpretar funções suaves; avaliar desempenho preditivo por AUC, TSS, sensibilidade e especificidade; gerar curvas de resposta parcial; produzir mapas de adequabilidade ambiental para \textit{Dinizia excelsa}; e comparar a flexibilidade do GAM com o GLM da unidade anterior.
\end{objetivos}

\section{Introdução}

Os Modelos Aditivos Generalizados, ou GAMs, ampliam a estrutura dos Modelos Lineares Generalizados ao substituir parte dos termos lineares por funções suaves estimadas a partir dos dados. Em SDM, isso é particularmente útil porque as respostas das espécies aos gradientes ambientais raramente são estritamente lineares \parencite{guisan2002,guisan2017habitat,wood2017generalized}.

Enquanto o GLM exige que a forma da resposta seja especificada explicitamente, por exemplo com termos lineares e quadráticos, o GAM permite que a relação entre ocorrência e ambiente seja estimada de forma mais flexível. Assim, ele pode representar respostas unimodais, limiares, platôs e transições ambientais complexas.

\begin{equationbox}
	\[
	\text{logit}(p_i) =
	\beta_0 + f_1(x_{1i}) + f_2(x_{2i}) + \cdots + f_k(x_{ki})
	\]
\end{equationbox}

em que \(p_i\) representa a probabilidade relativa de presença no local \(i\), \(\beta_0\) é o intercepto e \(f_k(x_{ki})\) são funções suaves estimadas para cada variável ambiental.

\begin{conceito}[title={\faBook\quad GAM em SDM}]
	Em SDM, o GAM estima relações não lineares entre presença-background e variáveis ambientais por meio de funções suaves. Sua principal vantagem é a flexibilidade; sua principal cautela é evitar suavizadores excessivamente complexos que representem ruído em vez de sinal ecológico.
\end{conceito}

\section{Suavizadores e complexidade}

A flexibilidade de um GAM depende da estrutura dos suavizadores. No pacote \texttt{mgcv}, o parâmetro \(k\) define a dimensão máxima da base usada para estimar a suavização. Valores muito baixos podem produzir respostas excessivamente simples; valores muito altos podem permitir sobreajuste.

\begin{atencao}[title={\faExclamationTriangle\quad Flexibilidade não é sinônimo de qualidade}]
	Um GAM muito flexível pode capturar padrões locais dos dados de calibração, mas perder capacidade de transferência espacial ou temporal. Em SDM, a forma da curva deve ser estatisticamente adequada e ecologicamente plausível.
\end{atencao}

\section{Preparação dos dados}

Nesta unidade, utilizamos a mesma base presença-background preparada na Unidade 5 e as variáveis selecionadas após diagnóstico de multicolinearidade na Unidade 4. Isso garante continuidade metodológica entre GLM, GAM e os demais algoritmos.

\rscriptfile{scripts/unidade07/01_estrutura_pastas.R}

\rscriptfile{scripts/unidade07/02_preparar_dados_gam.R}

\section{Ajuste do GAM}

O modelo é ajustado com distribuição binomial e função de ligação logit. Cada variável ambiental selecionada entra como função suave, permitindo respostas ecológicas não lineares.

\rscriptfile{scripts/unidade07/03_ajustar_gam_dinizia.R}

\begin{figure}[H]
	\centering
	\includegraphics[width=0.96\textwidth]{figuras/unidade07/unidade07_suavizadores_gam.png}
	\caption{Funções suaves estimadas pelo GAM para as variáveis ambientais selecionadas. Cada painel representa a contribuição parcial de uma variável para a adequabilidade ambiental de \textit{Dinizia excelsa}.}
	\label{fig:unidade07_suavizadores_gam}
\end{figure}

\section{Avaliação do modelo}

O desempenho preditivo do GAM é avaliado com AUC, TSS, sensibilidade, especificidade e matriz de confusão. Essa avaliação permite comparar o GAM com o GLM ajustado na Unidade 6.

\rscriptfile{scripts/unidade07/04_avaliar_gam_dinizia.R}

\begin{figure}[H]
	\centering
	\includegraphics[width=0.78\textwidth]{figuras/unidade07/unidade07_avaliacao_gam_patchwork.png}
	\caption{Curva ROC do GAM ajustado para \textit{Dinizia excelsa}.}
	\label{fig:unidade07_roc_gam}
\end{figure}

\section{Mapa de adequabilidade atual}

O GAM é projetado espacialmente para a área acessível \(M\), gerando um mapa contínuo de adequabilidade ambiental atual.

\rscriptfile{scripts/unidade07/05_predicao_espacial_gam.R}

\begin{figure}[H]
	\centering
	\includegraphics[width=0.88\textwidth]{figuras/unidade07/unidade07_adequabilidade_gam.png}
	\caption{Mapa de adequabilidade ambiental atual de \textit{Dinizia excelsa} estimado por GAM dentro da área acessível \(M\).}
	\label{fig:unidade07_adequabilidade_gam}
\end{figure}

\section{Curvas de resposta parcial}

As curvas de resposta parcial do GAM mostram respostas flexíveis ao longo dos gradientes ambientais, permitindo identificar faixas de maior adequabilidade, limiares e potenciais restrições climáticas.

\rscriptfile{scripts/unidade07/06_curvas_resposta_gam.R}

\begin{figure}[H]
	\centering
	\includegraphics[width=0.96\textwidth]{figuras/unidade07/unidade07_curvas_resposta_gam.png}
	\caption{Curvas de resposta parcial do GAM para as variáveis ambientais selecionadas.}
	\label{fig:unidade07_curvas_resposta_gam}
\end{figure}

\section{Comparação entre GLM e GAM}

A comparação entre GLM e GAM permite avaliar como a maior flexibilidade do GAM altera a distribuição espacial prevista e as curvas de resposta ambiental.

\rscriptfile{scripts/unidade07/07_comparar_glm_gam.R}

\begin{figure}[H]
	\centering
	\includegraphics[width=0.96\textwidth]{figuras/unidade07/unidade07_comparacao_glm_gam.png}
	\caption{Comparação entre mapas de adequabilidade estimados por GLM e GAM para \textit{Dinizia excelsa}.}
	\label{fig:unidade07_comparacao_glm_gam}
\end{figure}

\section{Pipeline completo}

\rscriptfile{scripts/unidade07/08_pipeline_gam.R}

\section{Síntese da unidade}

O GAM representa uma etapa intermediária entre modelos paramétricos simples e algoritmos de aprendizado de máquina. Ele mantém boa interpretabilidade, mas permite capturar relações ambientais não lineares. Para \textit{Dinizia excelsa}, o GAM permite visualizar respostas ecológicas flexíveis aos gradientes ambientais selecionados.

\begin{conceito}[title={\faBook\quad Mensagem central}]
	O GAM é especialmente útil quando se espera que a espécie responda de forma não linear aos gradientes ambientais. Em SDM, ele oferece um equilíbrio entre flexibilidade, interpretação ecológica e desempenho preditivo.
\end{conceito}

\section{Exercícios de fixação}

\begin{enumerate}
	\item Explique a diferença entre GLM e GAM.
	\item O que representa uma função suave em um GAM?
	\item Por que o parâmetro \(k\) deve ser escolhido com cautela?
	\item Explique por que GAMs podem capturar respostas ecológicas não lineares.
	\item O que significa sobreajuste em SDM?
	\item Como curvas de resposta parcial ajudam na interpretação ecológica?
	\item Compare vantagens e limitações de GLM e GAM em SDM.
\end{enumerate}

% ============================================================
% UNIDADE 8
% ============================================================

\chapter[Random Forest]{Random Forest aplicado}

\begin{objetivos}
	Ao final desta unidade, o estudante deverá compreender os fundamentos do algoritmo Random Forest; ajustar modelos presença-background para \textit{Dinizia excelsa}; controlar parâmetros para reduzir sobreajuste; avaliar desempenho com dados independentes de teste; comparar desempenho OOB e teste; interpretar importância das variáveis; gerar curvas de resposta parcial; produzir mapas de adequabilidade ambiental; e compreender limitações de algoritmos de aprendizado de máquina em SDM.
\end{objetivos}

\section{Introdução}

Random Forest é um algoritmo de aprendizado de máquina baseado em múltiplas árvores de decisão. Cada árvore é ajustada a partir de uma amostra dos dados e de um subconjunto aleatório das variáveis preditoras. A predição final é obtida pela média das árvores, no caso de regressão ou probabilidade, ou pelo voto majoritário, no caso de classificação \parencite{breiman2001random,cutler2007random,guisan2017habitat}.

Em Modelagem de Distribuição de Espécies, Random Forest é frequentemente utilizado por sua capacidade de representar relações não lineares, interações entre variáveis e respostas ambientais complexas. Entretanto, sua flexibilidade exige cuidado. Modelos muito complexos podem apresentar desempenho aparentemente excelente nos dados de calibração, mas baixa capacidade de generalização espacial ou temporal.

\begin{conceito}[title={\faBook\quad Random Forest em SDM}]
	Random Forest combina muitas árvores de decisão para estimar relações complexas entre ocorrência da espécie e variáveis ambientais. Sua força está na flexibilidade; seu risco está no sobreajuste quando a validação não é adequadamente controlada.
\end{conceito}

\section{Sobreajuste em Random Forest}

Um problema comum em SDM com algoritmos flexíveis é obter métricas muito elevadas, como AUC próxima de 1. Isso pode indicar excelente discriminação, mas também pode ser sintoma de sobreajuste, especialmente quando há autocorrelação espacial, muitos pontos de background, variáveis redundantes ou validação não independente.

Nesta unidade, adotamos quatro estratégias para reduzir esse risco:

\begin{itemize}
	\item separação treino/teste estratificada;
	\item uso de variáveis selecionadas após correlação e VIF;
	\item controle de complexidade por \texttt{mtry}, \texttt{min.node.size} e \texttt{sample.fraction};
	\item comparação entre desempenho OOB e desempenho no conjunto teste.
\end{itemize}

\begin{atencao}[title={\faExclamationTriangle\quad AUC muito alto exige cautela}]
	Em SDM, AUC próximo de 1 nem sempre significa modelo ecologicamente excelente. Pode indicar sobreajuste, viés espacial, background muito fácil de distinguir ou validação não independente.
\end{atencao}

\section{Preparação dos dados}

A base utilizada nesta unidade é a mesma preparada na Unidade 5, com presenças reais de \textit{Dinizia excelsa}, background amostrado dentro da área acessível \(M\) e variáveis ambientais selecionadas na Unidade 4.

\rscriptfile{scripts/unidade08/01_estrutura_pastas.R}

\rscriptfile{scripts/unidade08/02_preparar_dados_rf.R}

\section{Ajuste controlado do Random Forest}

O modelo é ajustado com o pacote \texttt{ranger}, usando probabilidade de presença como saída. Para reduzir sobreajuste, utilizamos valores moderados de \texttt{mtry}, \texttt{min.node.size} maior que o padrão e \texttt{sample.fraction} menor que 1.

\rscriptfile{scripts/unidade08/03_ajustar_rf_dinizia.R}

\begin{figure}[H]
	\centering
	\includegraphics[width=0.86\textwidth]{figuras/unidade08/unidade08_importancia_variaveis_rf.png}
	\caption{Importância das variáveis ambientais no modelo Random Forest ajustado para \textit{Dinizia excelsa}.}
	\label{fig:unidade08_importancia_rf}
\end{figure}

\section{Avaliação do modelo}

O desempenho do Random Forest é avaliado no conjunto teste, não nos dados de treino. Também comparamos AUC OOB e AUC teste. Diferenças muito grandes entre esses valores podem indicar instabilidade ou sobreajuste.

\rscriptfile{scripts/unidade08/04_avaliar_rf_dinizia.R}

\begin{figure}[H]
	\centering
	\includegraphics[width=0.78\textwidth]{figuras/unidade08/unidade08_avaliacao_rf_patchwork.png}
	\caption{Curva ROC do Random Forest avaliado no conjunto teste.}
	\label{fig:unidade08_roc_rf}
\end{figure}

\section{Mapa de adequabilidade atual}

Após a avaliação, o modelo é projetado para a área acessível \(M\), produzindo um mapa contínuo de adequabilidade ambiental atual.

\rscriptfile{scripts/unidade08/05_predicao_espacial_rf.R}

\begin{figure}[H]
	\centering
	\includegraphics[width=0.88\textwidth]{figuras/unidade08/unidade08_adequabilidade_rf.png}
	\caption{Mapa de adequabilidade ambiental atual de \textit{Dinizia excelsa} estimado por Random Forest dentro da área acessível \(M\).}
	\label{fig:unidade08_adequabilidade_rf}
\end{figure}

\begin{figure}[H]
	\centering
	\includegraphics[width=0.78\textwidth]{figuras/unidade08/unidade08_incerteza_sd_glm_gam_rf.png}
	\caption{Mapa de incerteza espacial obtido a partir do desvio-padrão das predições dos modelos GLM, GAM e Random Forest para \textit{Dinizia excelsa}. Valores elevados indicam regiões onde os modelos apresentam maior divergência nas estimativas de adequabilidade ambiental, enquanto valores reduzidos indicam maior consenso entre os algoritmos.}
	\label{fig:unidade08_incerteza_modelos}
\end{figure}

\section{Curvas de resposta parcial}

As curvas de resposta parcial auxiliam na interpretação ecológica do modelo. Em Random Forest, essas curvas não são coeficientes diretos, mas representam a tendência média da predição ao variar uma variável, mantendo as demais constantes.

\rscriptfile{scripts/unidade08/06_curvas_resposta_rf.R}

\begin{figure}[H]
	\centering
	\includegraphics[width=0.96\textwidth]{figuras/unidade08/unidade08_curvas_resposta_rf.png}
	\caption{Curvas de resposta parcial do Random Forest para as variáveis ambientais selecionadas.}
	\label{fig:unidade08_curvas_resposta_rf}
\end{figure}

\section{Comparação com GLM e GAM}

Por fim, comparamos o mapa produzido pelo Random Forest com os mapas gerados por GLM e GAM. Essa comparação mostra como modelos com diferentes níveis de flexibilidade podem produzir padrões espaciais distintos.

\rscriptfile{scripts/unidade08/07_comparar_glm_gam_rf.R}

\begin{figure}[H]
	\centering
	\includegraphics[width=0.96\textwidth]{figuras/unidade08/unidade08_comparacao_glm_gam_rf.png}
	\caption{Comparação entre mapas de adequabilidade estimados por GLM, GAM e Random Forest.}
	\label{fig:unidade08_comparacao_modelos}
\end{figure}

\section{Pipeline completo}

\rscriptfile{scripts/unidade08/08_pipeline_rf.R}

\section{Síntese da unidade}

Random Forest é um algoritmo poderoso para SDM, especialmente quando há relações não lineares e interações entre variáveis. Entretanto, seu uso exige validação cuidadosa. Métricas elevadas devem ser interpretadas em conjunto com diagnóstico de sobreajuste, coerência ecológica, curvas de resposta, distribuição espacial das predições e comparação com modelos mais simples.

\begin{conceito}[title={\faBook\quad Mensagem central}]
	Random Forest pode melhorar o desempenho preditivo, mas não substitui raciocínio ecológico. Um modelo complexo só é útil quando generaliza bem, produz respostas plausíveis e é avaliado de forma independente.
\end{conceito}

\section{Exercícios de fixação}

\begin{enumerate}
	\item Explique como o Random Forest combina múltiplas árvores de decisão.
	\item Por que Random Forest pode capturar relações não lineares?
	\item O que é sobreajuste em SDM?
	\item Por que AUC muito alto pode ser preocupante?
	\item Explique a diferença entre desempenho OOB e desempenho em teste.
	\item Qual o papel de \texttt{mtry}, \texttt{min.node.size} e \texttt{sample.fraction}?
	\item Como curvas de resposta parcial ajudam na interpretação de Random Forest?
\end{enumerate}

% ============================================================
% UNIDADE 9
% ============================================================

\chapter[BRT]{Boosted Regression Trees}

\begin{objetivos}
	Ao final desta unidade, o estudante deverá compreender os fundamentos das Boosted Regression Trees; ajustar modelos BRT para dados presença-background; controlar hiperparâmetros para reduzir sobreajuste; selecionar o número ótimo de árvores por validação cruzada; avaliar o modelo com AUC, TSS, sensibilidade e especificidade; interpretar importância relativa das variáveis; gerar curvas de resposta parcial; produzir mapas de adequabilidade ambiental de \textit{Dinizia excelsa}; e comparar BRT com GLM, GAM e Random Forest.
\end{objetivos}

\section{Introdução}

Boosted Regression Trees, ou BRT, combinam árvores de decisão com boosting. Diferentemente do Random Forest, que ajusta muitas árvores independentes e combina suas predições, o boosting ajusta árvores sequencialmente, de modo que cada nova árvore tenta melhorar os erros deixados pelas árvores anteriores. Em SDM, BRTs são amplamente usadas porque capturam respostas não lineares, interações entre variáveis e padrões ambientais complexos \parencite{elith2008working,elith2009,guisan2017habitat}.

\begin{conceito}[title={\faBook\quad BRT em SDM}]
	BRT é um algoritmo flexível que combina árvores de decisão e aprendizado sequencial. Ele pode representar relações ecológicas complexas, mas exige controle rigoroso de hiperparâmetros para evitar sobreajuste.
\end{conceito}

\section{Hiperparâmetros principais}

Os principais hiperparâmetros em BRT são:

\begin{itemize}
	\item \texttt{n.trees}: número máximo de árvores;
	\item \texttt{interaction.depth}: complexidade das interações;
	\item \texttt{shrinkage}: taxa de aprendizado;
	\item \texttt{bag.fraction}: fração dos dados usada em cada iteração;
	\item \texttt{n.minobsinnode}: número mínimo de observações nos nós terminais.
\end{itemize}

\begin{atencao}[title={\faExclamationTriangle\quad BRT e sobreajuste}]
	BRT pode produzir modelos excessivamente ajustados se o número de árvores for alto, a taxa de aprendizado for inadequada ou a profundidade das árvores for muito grande. Por isso, o número ótimo de árvores deve ser definido por validação cruzada.
\end{atencao}

\section{Preparação dos dados}

\rscriptfile{scripts/unidade09/01_estrutura_pastas.R}

\rscriptfile{scripts/unidade09/02_preparar_dados_brt.R}

\section{Ajuste do BRT}

O modelo BRT é ajustado usando distribuição Bernoulli, adequada para dados binários presença-background. A seleção do número ótimo de árvores é feita por validação cruzada interna.

\rscriptfile{scripts/unidade09/03_ajustar_brt_dinizia.R}

\begin{figure}[H]
	\centering
	\includegraphics[width=0.86\textwidth]{figuras/unidade09/unidade09_importancia_variaveis_brt.png}
	\caption{Importância relativa das variáveis ambientais no modelo BRT ajustado para \textit{Dinizia excelsa}.}
	\label{fig:unidade09_importancia_brt}
\end{figure}

\section{Avaliação do modelo}

\rscriptfile{scripts/unidade09/04_avaliar_brt_dinizia.R}

\begin{figure}[H]
	\centering
	\includegraphics[width=0.78\textwidth]{figuras/unidade09/unidade09_avaliacao_brt_patchwork.png}
	\caption{Curva ROC do BRT avaliado no conjunto teste.}
	\label{fig:unidade09_roc_brt}
\end{figure}

\begin{figure}[H]
	\centering
	\includegraphics[width=0.78\textwidth]{figuras/unidade09/unidade09_incerteza_sd_glm_gam_rf_brt.png}
	\caption{Incerteza associada às predições de distribuição potencial de \textit{Dinizia excelsa}, calculada como o desvio-padrão entre os modelos GLM, GAM, Random Forest e BRT. Áreas com maiores valores indicam menor consenso entre algoritmos e, consequentemente, maior incerteza preditiva, enquanto áreas com baixos valores representam regiões de elevada concordância entre modelos.}
	\label{fig:unidade09_incerteza_brt}
\end{figure}

\section{Mapa de adequabilidade atual}

\rscriptfile{scripts/unidade09/05_predicao_espacial_brt.R}

\begin{figure}[H]
	\centering
	\includegraphics[width=0.88\textwidth]{figuras/unidade09/unidade09_adequabilidade_brt.png}
	\caption{Mapa de adequabilidade ambiental atual de \textit{Dinizia excelsa} estimado por BRT dentro da área acessível \(M\).}
	\label{fig:unidade09_adequabilidade_brt}
\end{figure}

\section{Curvas de resposta parcial}

\rscriptfile{scripts/unidade09/06_curvas_resposta_brt.R}

\begin{figure}[H]
	\centering
	\includegraphics[width=0.96\textwidth]{figuras/unidade09/unidade09_curvas_resposta_brt.png}
	\caption{Curvas de resposta parcial do BRT para as variáveis ambientais selecionadas.}
	\label{fig:unidade09_curvas_resposta_brt}
\end{figure}

\section{Comparação com GLM, GAM e Random Forest}

\rscriptfile{scripts/unidade09/07_comparar_modelos_brt.R}

\begin{figure}[H]
	\centering
	\includegraphics[width=0.96\textwidth]{figuras/unidade09/unidade09_comparacao_glm_gam_rf_brt.png}
	\caption{Comparação entre mapas de adequabilidade estimados por GLM, GAM, Random Forest e BRT.}
	\label{fig:unidade09_comparacao_modelos}
\end{figure}

\section{Pipeline completo}

\rscriptfile{scripts/unidade09/08_pipeline_brt.R}

\section{Síntese da unidade}

BRT é um algoritmo poderoso para SDM por combinar flexibilidade, capacidade de modelar interações e boa performance preditiva. Entretanto, seu uso exige cautela metodológica, especialmente no controle de hiperparâmetros e na avaliação em dados independentes.

\begin{conceito}[title={\faBook\quad Mensagem central}]
	BRT pode revelar respostas ambientais complexas, mas sua credibilidade depende de validação adequada, controle de sobreajuste e interpretação ecológica das respostas estimadas.
\end{conceito}

\section{Exercícios de fixação}

\begin{enumerate}
	\item Explique a diferença entre Random Forest e BRT.
	\item O que significa boosting?
	\item Por que o parâmetro \texttt{shrinkage} é importante?
	\item Qual o papel de \texttt{interaction.depth}?
	\item Por que o número ótimo de árvores deve ser selecionado por validação cruzada?
	\item Como interpretar a importância relativa das variáveis em BRT?
	\item Por que curvas de resposta parcial devem ser avaliadas criticamente?
\end{enumerate}

% ============================================================
% UNIDADE 10
% ============================================================

\chapter[Maxent]{Máxima Entropia}

\begin{objetivos}
	Ao final desta unidade, o estudante deverá compreender os fundamentos do Maxent em SDM; interpretar o princípio da máxima entropia; diferenciar features lineares, quadráticas, hinge e produto; compreender o papel da regularização; ajustar modelos Maxent com dados presença-background de \textit{Dinizia excelsa}; realizar tuning de feature classes e regularization multiplier; avaliar o modelo com AUC, TSS, sensibilidade e especificidade; gerar curvas de resposta; produzir mapa de adequabilidade ambiental atual; e comparar Maxent com GLM, GAM, Random Forest e BRT.
\end{objetivos}

\section{Introdução}

O Maxent é um dos métodos mais utilizados em Modelagem de Distribuição de Espécies, especialmente quando os dados disponíveis são registros de presença e pontos de background. O método estima uma distribuição de máxima entropia, isto é, a distribuição mais uniforme possível, sujeita às restrições impostas pelas condições ambientais observadas nos locais de presença \parencite{phillips2006maximum,phillips2008modeling,elith2011statistical,merow2013practical,guisan2017habitat}.

Em termos práticos, o Maxent compara os valores ambientais nas presenças com os valores ambientais disponíveis no background. A partir dessa comparação, estima uma função de adequabilidade relativa que indica quais combinações ambientais são mais semelhantes às condições onde a espécie foi observada.

\begin{conceito}[title={\faBook\quad Maxent em SDM}]
	Maxent é um método presença-background baseado no princípio da máxima entropia. Ele estima uma distribuição espacial de adequabilidade ambiental que satisfaz as restrições ambientais observadas nas presenças, evitando assumir mais estrutura do que a informação disponível permite.
\end{conceito}

\section{Máxima entropia e adequabilidade ambiental}

O princípio da máxima entropia busca a distribuição menos enviesada possível, dado um conjunto de restrições. Em SDM, essas restrições são derivadas das variáveis ambientais nos pontos de presença. Assim, o modelo busca uma distribuição que seja o mais uniforme possível, mas que respeite as características ambientais associadas à ocorrência da espécie.

\begin{equationbox}
	\[
	\hat{p}(x) =
	\frac{\exp\left(\sum_j \lambda_j f_j(x)\right)}
	{Z}
	\]
\end{equationbox}

em que \(\hat{p}(x)\) representa a adequabilidade relativa em um local \(x\), \(f_j(x)\) são funções das variáveis ambientais, \(\lambda_j\) são pesos estimados pelo modelo e \(Z\) é uma constante de normalização.

\section{Feature classes}

No Maxent, as relações entre espécie e ambiente são representadas por diferentes tipos de funções, chamadas \textit{feature classes}. As principais são:

\begin{itemize}
	\item linear (L): respostas lineares;
	\item quadratic (Q): respostas unimodais ou em forma de parábola;
	\item product (P): interações entre variáveis;
	\item hinge (H): respostas com limiares e mudanças abruptas;
	\item threshold (T): respostas por cortes discretos.
\end{itemize}

\begin{atencao}[title={\faExclamationTriangle\quad Complexidade do Maxent}]
	Modelos Maxent com muitas features e baixa regularização podem apresentar sobreajuste. Por isso, a escolha de feature classes e do parâmetro de regularização deve ser feita por tuning e validação independente.
\end{atencao}

\section{Regularização}

A regularização controla a complexidade do modelo. Valores baixos permitem respostas mais flexíveis, mas aumentam o risco de sobreajuste. Valores altos tornam o modelo mais simples e suave, podendo reduzir a capacidade de capturar padrões ecológicos reais. O parâmetro geralmente avaliado é o \textit{regularization multiplier}.

\begin{conceito}[title={\faBook\quad Regularização}]
	A regularização penaliza modelos excessivamente complexos. Em SDM, ela é fundamental para aumentar a capacidade de transferência espacial e temporal, especialmente em projeções futuras ou paleoclimáticas.
\end{conceito}

\section{Preparação dos dados}

Nesta unidade utilizamos a base presença-background construída na Unidade 5 e as variáveis selecionadas na Unidade 4. O ajuste será feito com o pacote \texttt{maxnet}, uma implementação em R do Maxent baseada em modelos penalizados, sem necessidade do arquivo Java \texttt{maxent.jar}.

\rscriptfile{scripts/unidade10/01_estrutura_pastas.R}

\rscriptfile{scripts/unidade10/02_preparar_dados_maxent.R}

\section{Tuning de feature classes e regularização}

O tuning avalia diferentes combinações de feature classes e multiplicadores de regularização. O objetivo é escolher um modelo com bom desempenho preditivo, mas sem complexidade excessiva.

\rscriptfile{scripts/unidade10/03_tuning_maxent_dinizia.R}

\begin{figure}[H]
	\centering
	\includegraphics[width=0.88\textwidth]{figuras/unidade10/unidade10_tuning_maxent.png}
	\caption{Resultados do tuning do Maxent para \textit{Dinizia excelsa}, considerando diferentes combinações de feature classes e regularization multiplier.}
	\label{fig:unidade10_tuning_maxent}
\end{figure}

\section{Ajuste do modelo Maxent final}

Após o tuning, o melhor conjunto de hiperparâmetros é usado para ajustar o modelo final.

\rscriptfile{scripts/unidade10/04_ajustar_maxent_final.R}

\begin{figure}[H]
	\centering
	\includegraphics[width=0.86\textwidth]{figuras/unidade10/unidade10_importancia_variaveis_maxent.png}
	\caption{Importância relativa das variáveis ambientais no modelo Maxent final ajustado para \textit{Dinizia excelsa}.}
	\label{fig:unidade10_importancia_maxent}
\end{figure}

\section{Avaliação do modelo}

O modelo Maxent é avaliado com dados de teste independentes usando AUC, TSS, sensibilidade, especificidade e matriz de confusão.

\rscriptfile{scripts/unidade10/05_avaliar_maxent_dinizia.R}

\begin{figure}[H]
	\centering
	\includegraphics[width=0.78\textwidth]{figuras/unidade10/unidade10_avaliacao_maxent_patchwork.png}
	\caption{Curva ROC do Maxent avaliado no conjunto teste para \textit{Dinizia excelsa}.}
	\label{fig:unidade10_roc_maxent}
\end{figure}

\section{Mapa de adequabilidade atual}

O modelo final é projetado espacialmente para a área acessível \(M\), produzindo um mapa contínuo de adequabilidade ambiental atual.

\rscriptfile{scripts/unidade10/06_predicao_espacial_maxent.R}

\begin{figure}[H]
	\centering
	\includegraphics[width=0.88\textwidth]{figuras/unidade10/unidade10_adequabilidade_maxent.png}
	\caption{Mapa de adequabilidade ambiental atual de \textit{Dinizia excelsa} estimado por Maxent dentro da área acessível \(M\).}
	\label{fig:unidade10_adequabilidade_maxent}
\end{figure}

\section{Curvas de resposta parcial}

As curvas de resposta parcial permitem visualizar como a adequabilidade estimada pelo Maxent varia ao longo dos gradientes ambientais selecionados.

\rscriptfile{scripts/unidade10/07_curvas_resposta_maxent.R}

\begin{figure}[H]
	\centering
	\includegraphics[width=0.96\textwidth]{figuras/unidade10/unidade10_curvas_resposta_maxent.png}
	\caption{Curvas de resposta parcial do Maxent para as variáveis ambientais selecionadas.}
	\label{fig:unidade10_curvas_resposta_maxent}
\end{figure}

\section{Comparação com os modelos anteriores}

\rscriptfile{scripts/unidade10/08_comparar_modelos_maxent.R}

\begin{figure}[H]
	\centering
	\includegraphics[width=0.98\textwidth]{figuras/unidade10/unidade10_comparacao_modelos_maxent.png}
	\caption{Comparação entre mapas de adequabilidade estimados por GLM, GAM, Random Forest, BRT e Maxent.}
	\label{fig:unidade10_comparacao_modelos}
\end{figure}

\section{Pipeline completo}

\rscriptfile{scripts/unidade10/09_pipeline_maxent.R}

\section{Síntese da unidade}

O Maxent é um método robusto, amplamente utilizado e especialmente adequado para dados de presença-background. Entretanto, sua aplicação exige decisões cuidadosas sobre background, feature classes, regularização, tuning e validação. Modelos excessivamente complexos podem produzir mapas aparentemente detalhados, mas pouco transferíveis para outros períodos climáticos ou regiões.

\begin{figure}[H]
	\centering
	\includegraphics[width=\textwidth]{figuras/unidade10/unidade10_sintese_consenso_incerteza_glm_gam_rf_brt_maxent.png}
	\caption{
		Síntese espacial multi-modelo para \textit{Dinizia excelsa} no bioma Amazônia. 
		O painel esquerdo apresenta o consenso médio de adequabilidade ambiental obtido a partir da média das predições dos modelos GLM, GAM, Random Forest, BRT e MaxEnt. 
		O painel direito apresenta a incerteza espacial associada às predições, estimada pelo desvio-padrão entre os modelos. 
		Áreas com elevados valores de consenso indicam maior concordância quanto à adequabilidade ambiental da espécie, enquanto áreas com maior desvio-padrão representam regiões de maior divergência entre algoritmos e, consequentemente, maior incerteza preditiva.
	}
	\label{fig:unidade10_consenso_incerteza}
\end{figure}

\begin{conceito}[title={\faBook\quad Mensagem central}]
	Maxent não deve ser usado como uma caixa-preta. O desempenho do modelo depende da definição adequada da área acessível, da escolha do background, da seleção de variáveis, da regularização e da avaliação independente.
\end{conceito}

%=========================================================
\subsection*{Síntese comparativa dos algoritmos}

Após a apresentação dos principais algoritmos utilizados em SDM, torna-se
importante compreender suas diferenças conceituais, vantagens operacionais,
limitações metodológicas e potencial de interpretação ecológica. A
Tabela~\ref{tab:comparacao_algoritmos_sdm} apresenta uma visão integrada dos
métodos abordados nesta apostila, auxiliando na seleção do algoritmo mais
adequado para diferentes objetivos ecológicos e biogeográficos.

\begin{table}[H]
	\centering
	\caption{Comparação dos principais algoritmos utilizados em Modelagem de Distribuição de Espécies (SDM).}
	\label{tab:comparacao_algoritmos_sdm}
	
	\scriptsize
	
	\begin{tabular}{
			p{0.12\textwidth}
			p{0.16\textwidth}
			p{0.14\textwidth}
			p{0.21\textwidth}
			p{0.22\textwidth}
			p{0.11\textwidth}
		}
		
		\toprule
		\textbf{Método} &
		\textbf{Tipo de dado} &
		\textbf{Relação ecológica} &
		\textbf{Vantagens} &
		\textbf{Limitações} &
		\textbf{Interpretação} \\
		\midrule
		
		GLM &
		Presença--Ausência &
		Linear &
		Simples, robusto e amplamente utilizado &
		Baixa flexibilidade para relações complexas &
		Excelente \\
		
		GAM &
		Presença--Ausência &
		Não linear &
		Flexível e interpretável &
		Mais complexo que o GLM &
		Muito alta \\
		
		Random Forest &
		Presença--Ausência &
		Não linear &
		Elevado desempenho preditivo &
		Menor interpretabilidade ecológica &
		Média \\
		
		BRT &
		Presença--Ausência &
		Não linear &
		Captura interações e relações não lineares &
		Ajuste de hiperparâmetros delicado &
		Média \\
		
		Maxent &
		Presença--Background &
		Não linear &
		Excelente para dados de presença apenas &
		Sensível ao background e ao viés amostral &
		Boa \\
		
		Ensemble &
		Múltiplos &
		Integrado &
		Reduz a incerteza algorítmica &
		Maior complexidade de interpretação &
		Variável \\
		
		\bottomrule
		
	\end{tabular}
	
\end{table}

%=========================================================
\subsection*{Guia prático de seleção de algoritmos}

\begin{table}[H]
	\centering
	\caption{Guia prático para seleção de algoritmos em diferentes aplicações de SDM.}
	\label{tab:guia_escolha_algoritmos}
	
	\small
	
	\begin{tabular}{
			p{0.48\textwidth}
			p{0.36\textwidth}
		}
		
		\toprule
		\textbf{Situação} &
		\textbf{Método recomendado} \\
		\midrule
		
		Poucos registros de ocorrência &
		Maxent \\
		
		Interpretação ecológica dos efeitos ambientais &
		GLM ou GAM \\
		
		Predição com máximo desempenho estatístico &
		Random Forest ou BRT \\
		
		Modelagem para conservação da biodiversidade &
		Ensemble \\
		
		Projeções sob mudanças climáticas &
		Ensemble \\
		
		Análise exploratória inicial &
		GLM \\
		
		Curvas de resposta ecológica &
		GAM \\
		
		Avaliação da importância relativa das variáveis &
		Random Forest ou BRT \\
		
		Dados de presença apenas &
		Maxent \\
		
		Redução da incerteza algorítmica &
		Ensemble \\
		
		\bottomrule
		
	\end{tabular}
	
\end{table}

%=========================================================
\begin{center}
	\fbox{
		\begin{minipage}{0.90\textwidth}
			
			\textbf{Síntese metodológica.}
			
			Não existe um algoritmo universalmente superior para todos os problemas
			ecológicos. O desempenho de um modelo depende da qualidade dos dados,
			da definição da área acessível ($M$), da seleção de variáveis ambientais,
			da estratégia de validação e da pergunta ecológica proposta.
			
			Em aplicações modernas de conservação e mudanças climáticas, abordagens
			\textit{ensemble} são frequentemente recomendadas por reduzirem a dependência
			de algoritmos individuais e fornecerem estimativas mais robustas da
			adequabilidade ambiental.
			
		\end{minipage}
	}
\end{center}
%=========================================================

\section{Exercícios de fixação}

\begin{enumerate}
	\item Explique o princípio da máxima entropia.
	\item Diferencie dados de presença, ausência e background.
	\item O que são feature classes no Maxent?
	\item Qual o papel da regularização?
	\item Por que modelos Maxent muito complexos podem ser problemáticos?
	\item Explique por que tuning é necessário em Maxent.
	\item Compare Maxent com GLM, GAM, Random Forest e BRT.
\end{enumerate}

% ============================================================
% UNIDADE 11
% ============================================================

\chapter[Modelagem Ensemble]{Modelagem Ensemble}

\begin{objetivos}
	Ao final desta unidade, o estudante deverá compreender o conceito de modelagem ensemble em SDM; integrar predições de GLM, GAM, Random Forest, BRT e Maxent; comparar ensemble por média simples, média ponderada e consenso binário; ponderar modelos por desempenho preditivo; gerar mapas de adequabilidade ensemble para \textit{Dinizia excelsa}; avaliar a concordância entre modelos; e preparar produtos para projeções futuras, avaliação de incertezas e aplicações à conservação.
\end{objetivos}

\section{Introdução}

A Modelagem Ensemble é uma abordagem amplamente utilizada em Modelagem de Distribuição de Espécies para combinar predições de diferentes algoritmos. Em vez de depender de um único modelo, o ensemble integra múltiplas representações estatísticas da relação entre ocorrência e ambiente, reduzindo a dependência de escolhas algorítmicas específicas \parencite{araujo2007ensemble,thuiller2009biomod,guisan2017habitat}.

Em SDM, diferentes algoritmos podem produzir mapas distintos porque representam respostas ambientais de maneiras diferentes. GLM e GAM são mais interpretáveis; Random Forest e BRT capturam relações complexas; Maxent é especialmente adequado para dados presença-background. O ensemble sintetiza essas diferenças em um produto integrado.

\begin{conceito}[title={\faBook\quad Ensemble em SDM}]
	Um ensemble combina predições de múltiplos modelos para representar uma estimativa integrada da adequabilidade ambiental. Ele não elimina incertezas, mas permite quantificar concordância e divergência entre algoritmos.
\end{conceito}

\section{Tipos de ensemble}

Nesta unidade serão considerados três tipos principais de ensemble:

\begin{itemize}
	\item média simples das predições;
	\item média ponderada por desempenho preditivo;
	\item consenso binário baseado em limiares de presença-ausência.
\end{itemize}

A média simples considera todos os modelos com o mesmo peso. A média ponderada atribui maior influência aos modelos com melhor desempenho. O consenso binário representa a proporção de modelos que classificam uma célula como adequada.

\begin{atencao}[title={\faExclamationTriangle\quad Ensemble não corrige modelo ruim}]
	O ensemble deve ser construído a partir de modelos minimamente aceitáveis. Incluir modelos com baixa qualidade, forte sobreajuste ou respostas ecologicamente incoerentes pode degradar o resultado final.
\end{atencao}

\section{Preparação das predições}

Nesta unidade, integramos os mapas de adequabilidade gerados nas Unidades 6 a 10. Todos os rasters devem possuir a mesma extensão, resolução, projeção e máscara espacial.

\rscriptfile{scripts/unidade11/01_estrutura_pastas.R}

\rscriptfile{scripts/unidade11/02_preparar_predicoes_ensemble.R}

\section{Ensemble por média simples}

A média simples representa a adequabilidade média prevista pelos modelos. É uma abordagem transparente e útil quando todos os modelos têm desempenho semelhante.

\rscriptfile{scripts/unidade11/03_ensemble_media_simples.R}

\begin{figure}[H]
	\centering
	\includegraphics[width=0.88\textwidth]{figuras/unidade11/unidade11_ensemble_media_simples.png}
	\caption{Mapa ensemble por média simples das predições de GLM, GAM, Random Forest, BRT e Maxent para \textit{Dinizia excelsa}.}
	\label{fig:unidade11_ensemble_media}
\end{figure}

\section{Ensemble ponderado por desempenho}

O ensemble ponderado utiliza métricas de avaliação para atribuir maior peso aos modelos com melhor desempenho. Nesta unidade, usamos AUC de teste como critério principal, mas o mesmo procedimento pode ser adaptado para TSS, Boyce Index ou métricas compostas.

\rscriptfile{scripts/unidade11/04_ensemble_ponderado.R}

\begin{figure}[H]
	\centering
	\includegraphics[width=0.88\textwidth]{figuras/unidade11/unidade11_ensemble_ponderado.png}
	\caption{Mapa ensemble ponderado por desempenho preditivo para \textit{Dinizia excelsa}.}
	\label{fig:unidade11_ensemble_ponderado}
\end{figure}

\section{Consenso binário}

O consenso binário converte cada mapa contínuo em presença-ausência usando limiares definidos por TSS. O mapa final representa a proporção de modelos que classificam cada célula como adequada.

\rscriptfile{scripts/unidade11/05_consenso_binario.R}

\begin{figure}[H]
	\centering
	\includegraphics[width=0.88\textwidth]{figuras/unidade11/unidade11_consenso_binario.png}
	\caption{Consenso binário entre os modelos. Valores próximos de 1 indicam alta concordância entre modelos quanto à adequabilidade ambiental.}
	\label{fig:unidade11_consenso_binario}
\end{figure}

\section{Incerteza entre algoritmos}

A variação entre modelos pode ser usada como uma medida de incerteza algorítmica. Regiões com alta adequabilidade média e baixa incerteza são mais robustas; regiões com alta divergência entre modelos exigem interpretação cautelosa.

\rscriptfile{scripts/unidade11/06_incerteza_algoritmica.R}

\begin{figure}[H]
	\centering
	\includegraphics[width=0.88\textwidth]{figuras/unidade11/unidade11_incerteza_algoritmica.png}
	\caption{Incerteza algorítmica estimada pelo desvio-padrão das predições entre modelos.}
	\label{fig:unidade11_incerteza_algoritmica}
\end{figure}

\section{Comparação integrada}

\rscriptfile{scripts/unidade11/07_comparacao_ensemble_modelos.R}

\begin{figure}[H]
	\centering
	\includegraphics[width=0.98\textwidth]{figuras/unidade11/unidade11_comparacao_ensemble.png}
	\caption{Comparação entre modelos individuais, ensemble médio, ensemble ponderado, consenso binário e incerteza algorítmica.}
	\label{fig:unidade11_comparacao_ensemble}
\end{figure}

\section{Pipeline completo}

\rscriptfile{scripts/unidade11/08_pipeline_ensemble.R}

\section{Síntese da unidade}

A modelagem ensemble permite sintetizar diferentes algoritmos em um produto integrado de adequabilidade ambiental. Para \textit{Dinizia excelsa}, o ensemble combina modelos estatísticos interpretáveis e algoritmos flexíveis, oferecendo uma visão mais robusta da distribuição potencial atual dentro da área acessível \(M\).

\begin{conceito}[title={\faBook\quad Mensagem central}]
	O ensemble não substitui a avaliação crítica dos modelos individuais. Ele deve ser interpretado junto com métricas de desempenho, curvas de resposta, coerência ecológica e mapas de incerteza.
\end{conceito}

\section{Exercícios de fixação}

\begin{enumerate}
	\item Explique por que ensembles são usados em SDM.
	\item Diferencie média simples, média ponderada e consenso binário.
	\item Por que modelos ruins não devem ser incluídos no ensemble?
	\item Como a incerteza algorítmica pode ser mapeada?
	\item Por que o ensemble ponderado depende da escolha da métrica?
	\item Como interpretar áreas com alta adequabilidade e alta incerteza?
	\item Explique como o ensemble será útil nas projeções climáticas futuras.
\end{enumerate}

% ============================================================
% UNIDADE 12
% ============================================================

\chapter[Avaliação de modelos]{Avaliação de Modelos}

\begin{objetivos}
	Ao final desta unidade, o estudante deverá compreender os principais critérios de avaliação de modelos de distribuição de espécies; calcular e interpretar AUC, TSS, Kappa, sensibilidade, especificidade, acurácia e Boyce Index; comparar GLM, GAM, Random Forest, BRT, Maxent e Ensemble; interpretar curvas ROC; avaliar limiares de conversão presença-ausência; produzir mapas binários; e discutir limitações das métricas em estudos presença-background.
\end{objetivos}

\section{Introdução}

A avaliação de modelos é uma etapa essencial em Modelagem de Distribuição de Espécies. Um modelo pode produzir mapas visualmente plausíveis, mas ainda assim apresentar baixa capacidade de discriminação, baixa sensibilidade, baixa especificidade ou forte dependência do conjunto de calibração. Por isso, modelos de SDM devem ser avaliados com dados independentes sempre que possível \parencite{fielding1997review,allouche2006assessing,hirzel2006evaluating,li2013thresholds,guisan2017habitat}.

Nesta unidade, a avaliação será aplicada ao fluxo real da apostila, comparando os modelos ajustados para \textit{Dinizia excelsa}: GLM, GAM, Random Forest, BRT, Maxent e Ensemble. O objetivo não é escolher automaticamente o modelo com maior métrica, mas compreender o desempenho, as limitações e a coerência ecológica de cada abordagem.

\begin{conceito}[title={\faBook\quad Avaliação não é apenas ranqueamento}]
	A avaliação de modelos não deve ser reduzida à escolha do maior AUC. Um modelo útil em SDM deve apresentar bom desempenho preditivo, coerência ecológica, respostas interpretáveis e capacidade de generalização espacial ou temporal.
\end{conceito}

\section{AUC e curva ROC}

A curva ROC representa a relação entre sensibilidade e \(1 -\) especificidade ao longo de diferentes limiares. A área sob a curva, ou AUC, mede a capacidade do modelo de discriminar presenças de background ou ausências.

\begin{equationbox}
	\[
	AUC = P(\hat{y}_{presença} > \hat{y}_{background})
	\]
\end{equationbox}

Em termos intuitivos, AUC representa a probabilidade de uma presença receber valor de adequabilidade maior que um ponto de background escolhido aleatoriamente.

\begin{atencao}[title={\faExclamationTriangle\quad Limitações do AUC}]
	Em dados presença-background, AUC pode ser influenciado pela extensão geográfica, pelo background escolhido, pelo viés amostral e pela prevalência. Valores altos não garantem necessariamente bom modelo ecológico.
\end{atencao}

\section{Sensibilidade, especificidade e TSS}

A sensibilidade mede a proporção de presenças corretamente previstas como adequadas. A especificidade mede a proporção de backgrounds ou ausências corretamente classificadas como inadequadas.

\begin{equationbox}
	\[
	Sensibilidade = \frac{TP}{TP + FN}
	\]
	
	\[
	Especificidade = \frac{TN}{TN + FP}
	\]
	
	\[
	TSS = Sensibilidade + Especificidade - 1
	\]
\end{equationbox}

O TSS varia de \(-1\) a \(1\), sendo valores próximos de \(1\) indicativos de melhor discriminação binária.

\section{Kappa}

O índice Kappa mede o acordo entre predição binária e observação, corrigindo pelo acordo esperado ao acaso. Embora tenha sido amplamente utilizado em SDM, ele pode ser sensível à prevalência e à distribuição dos dados.

\begin{equationbox}
	\[
	Kappa = \frac{p_o - p_e}{1 - p_e}
	\]
\end{equationbox}

em que \(p_o\) é o acordo observado e \(p_e\) é o acordo esperado ao acaso.

\section{Boyce Index}

O Boyce Index é uma métrica especialmente útil para modelos presença-background ou presença-apenas. Ele avalia se as presenças tendem a ocorrer em áreas com maiores valores de adequabilidade previstos pelo modelo \parencite{hirzel2006evaluating}.

\begin{conceito}[title={\faBook\quad Boyce Index}]
	O Boyce Index mede a consistência entre adequabilidade prevista e frequência observada de presenças. Valores positivos indicam que as presenças se concentram em classes de maior adequabilidade.
\end{conceito}

\section{Preparação da avaliação integrada}

Nesta unidade, organizamos as predições de teste dos modelos individuais e do ensemble para construir uma avaliação comparativa.

\rscriptfile{scripts/unidade12/01_estrutura_pastas.R}

\rscriptfile{scripts/unidade12/02_preparar_predicoes_avaliacao.R}

\section{Cálculo das métricas}

\rscriptfile{scripts/unidade12/03_calcular_metricas_modelos.R}

\begin{figure}[H]
	\centering
	\includegraphics[width=0.95\textwidth]{figuras/unidade12/unidade12_comparacao_metricas.png}
	\caption{Comparação das métricas de avaliação dos modelos ajustados para \textit{Dinizia excelsa}.}
	\label{fig:unidade12_metricas}
\end{figure}

\section{Curvas ROC comparativas}

\rscriptfile{scripts/unidade12/04_curvas_roc_comparativas.R}

\begin{figure}[H]
	\centering
	\includegraphics[width=0.86\textwidth]{figuras/unidade12/unidade12_curvas_roc_comparativas.png}
	\caption{Curvas ROC comparativas para GLM, GAM, Random Forest, BRT, Maxent e Ensemble.}
	\label{fig:unidade12_roc_comparativas}
\end{figure}

\section{Matrizes de confusão}

\rscriptfile{scripts/unidade12/05_matrizes_confusao.R}

\begin{figure}[H]
	\centering
	\includegraphics[width=0.95\textwidth]{figuras/unidade12/unidade12_matrizes_confusao.png}
	\caption{Matrizes de confusão dos modelos avaliados com limiares ótimos definidos pelo TSS.}
	\label{fig:unidade12_matrizes_confusao}
\end{figure}

\section{Mapas binários}

\rscriptfile{scripts/unidade12/06_mapas_binarios.R}

\begin{figure}[H]
	\centering
	\includegraphics[width=0.98\textwidth]{figuras/unidade12/unidade12_mapas_binarios.png}
	\caption{Mapas binários de adequabilidade ambiental para \textit{Dinizia excelsa}, gerados a partir dos limiares ótimos por TSS.}
	\label{fig:unidade12_mapas_binarios}
\end{figure}

\section{Síntese comparativa}

\rscriptfile{scripts/unidade12/07_sintese_avaliacao.R}

\begin{figure}[H]
	\centering
	\includegraphics[width=0.95\textwidth]{figuras/unidade12/unidade12_ranking_modelos.png}
	\caption{Ranking comparativo dos modelos considerando AUC, TSS, Kappa, sensibilidade, especificidade e Boyce Index.}
	\label{fig:unidade12_ranking}
\end{figure}

\section{Pipeline completo}

\rscriptfile{scripts/unidade12/08_pipeline_avaliacao.R}

\section{Síntese da unidade}

A avaliação de modelos deve combinar múltiplas métricas, validação independente e interpretação ecológica. Para \textit{Dinizia excelsa}, as métricas ajudam a comparar modelos com diferentes estruturas, mas a decisão final sobre quais modelos utilizar em projeções futuras deve considerar também coerência espacial, curvas de resposta, incerteza e plausibilidade ecológica.

\begin{conceito}[title={\faBook\quad Mensagem central}]
	Nenhuma métrica isolada é suficiente para validar um SDM. A avaliação robusta combina métricas contínuas, métricas binárias, interpretação ecológica e análise espacial das predições.
\end{conceito}

\section{Exercícios de fixação}

\begin{enumerate}
	\item Explique a diferença entre AUC e TSS.
	\item Por que AUC pode ser problemático em dados presença-background?
	\item O que significam sensibilidade e especificidade?
	\item Como o limiar ótimo por TSS é definido?
	\item Explique o significado ecológico do Boyce Index.
	\item Por que mapas binários devem ser interpretados com cautela?
	\item Como você escolheria modelos para projeções climáticas futuras?
\end{enumerate}

% ============================================================
% UNIDADE 13
% ============================================================

\chapter[Projeções climáticas futuras]{Projeções Climáticas Futuras em SDM}

\begin{objetivos}
	Ao final desta unidade, o estudante deverá compreender como modelos de distribuição de espécies podem ser projetados para cenários climáticos futuros; reconhecer os principais cenários SSP do CMIP6; preparar variáveis ambientais futuras compatíveis com as variáveis atuais; projetar modelos ajustados para \textit{Dinizia excelsa}; gerar mapas futuros de adequabilidade; calcular mudanças em adequabilidade; identificar áreas de estabilidade, perda e ganho; e discutir incertezas associadas a cenários, modelos climáticos e algoritmos de SDM.
\end{objetivos}

\section{Introdução}

As projeções climáticas futuras permitem avaliar como a adequabilidade ambiental potencial de uma espécie pode se alterar sob diferentes cenários de mudanças climáticas. Em Modelagem de Distribuição de Espécies, essa etapa consiste em ajustar modelos com dados atuais e projetá-los para camadas ambientais futuras derivadas de Modelos Climáticos Globais, períodos futuros e cenários socioeconômicos \parencite{guisan2017habitat,ipcc2021,araujo2006climate,thuiller2005climate}.

Nesta unidade, o fluxo será estruturado para \textit{Dinizia excelsa}, usando as variáveis ambientais selecionadas nas unidades anteriores, a área acessível \(M\), os modelos ajustados e o ensemble atual. O objetivo é construir uma base reprodutível para projeções sob cenários SSP126, SSP245, SSP370 e SSP585.

\begin{conceito}[title={\faBook\quad Projeção climática futura}]
	Uma projeção futura em SDM consiste em aplicar um modelo calibrado sob condições atuais a um conjunto de variáveis ambientais representando condições climáticas futuras. O resultado não é uma previsão determinística de presença, mas uma estimativa de adequabilidade ambiental sob determinado cenário.
\end{conceito}

\section{CMIP6 e cenários SSP}

O CMIP6 reúne simulações climáticas globais produzidas por diferentes modelos climáticos. Os cenários SSP, ou \textit{Shared Socioeconomic Pathways}, representam trajetórias alternativas de desenvolvimento socioeconômico e emissões de gases de efeito estufa.

\begin{itemize}
	\item SSP126: cenário de baixas emissões e forte mitigação climática;
	\item SSP245: cenário intermediário;
	\item SSP370: cenário de emissões elevadas;
	\item SSP585: cenário de emissões muito elevadas.
\end{itemize}

\begin{atencao}[title={\faExclamationTriangle\quad Cenário não é previsão única}]
	Cada SSP representa uma trajetória possível, não uma previsão certa. Por isso, projeções futuras devem ser interpretadas como cenários alternativos de adequabilidade ambiental.
\end{atencao}

\section{Compatibilidade entre variáveis atuais e futuras}

Para projetar modelos de SDM, as camadas futuras devem possuir os mesmos nomes, unidades, resolução espacial, sistema de coordenadas e extensão das variáveis usadas no ajuste atual. Qualquer inconsistência entre as variáveis atuais e futuras pode gerar erros ou projeções inválidas.

\rscriptfile{scripts/unidade13/01_estrutura_pastas.R}

\rscriptfile{scripts/unidade13/02_preparar_variaveis_futuras.R}

\section{Projeção dos modelos individuais}

Nesta etapa, os modelos GLM, GAM, Random Forest, BRT e Maxent são projetados para os cenários futuros. A projeção é feita usando as mesmas variáveis ambientais selecionadas e a mesma máscara espacial da área acessível \(M\).

\rscriptfile{scripts/unidade13/03_projetar_modelos_futuros.R}

\begin{figure}[H]
	\centering
	\includegraphics[width=0.98\textwidth]{figuras/unidade13/unidade13_modelos_futuros_ssp585.png}
	\caption{Projeções futuras dos modelos individuais para \textit{Dinizia excelsa} sob cenário SSP585.}
	\label{fig:unidade13_modelos_futuros}
\end{figure}

\section{Ensemble futuro}

As projeções individuais são combinadas para gerar mapas ensemble futuros por cenário. O ensemble reduz a dependência de um único algoritmo e permite comparar padrões médios de adequabilidade entre cenários.

\rscriptfile{scripts/unidade13/04_ensemble_futuro.R}

\begin{figure}[H]
	\centering
	\includegraphics[width=0.96\textwidth]{figuras/unidade13/unidade13_ensemble_futuro_cenarios.png}
	\caption{Mapas ensemble futuros de adequabilidade ambiental de \textit{Dinizia excelsa} sob diferentes cenários SSP.}
	\label{fig:unidade13_ensemble_futuro}
\end{figure}

\section{Mudança de adequabilidade}

A mudança de adequabilidade é calculada pela diferença entre a adequabilidade futura e a adequabilidade atual. Valores positivos indicam aumento relativo de adequabilidade; valores negativos indicam redução.

\begin{equationbox}
	\[
	\Delta S = S_{futuro} - S_{atual}
	\]
\end{equationbox}

em que \(S_{futuro}\) é a adequabilidade projetada para um cenário futuro e \(S_{atual}\) é a adequabilidade estimada no período atual.

\rscriptfile{scripts/unidade13/05_mudanca_adequabilidade.R}

\begin{figure}[H]
	\centering
	\includegraphics[width=0.96\textwidth]{figuras/unidade13/unidade13_mudanca_adequabilidade.png}
	\caption{Mudança de adequabilidade ambiental futura em relação ao cenário atual para \textit{Dinizia excelsa}.}
	\label{fig:unidade13_mudanca}
\end{figure}

\section{Áreas estáveis, perdas e ganhos}

A comparação entre mapas binários atuais e futuros permite classificar as células em áreas estáveis, áreas de perda potencial, áreas de ganho potencial e áreas persistentemente inadequadas.

\rscriptfile{scripts/unidade13/06_estabilidade_perda_ganho.R}

\begin{figure}[H]
	\centering
	\includegraphics[width=0.96\textwidth]{figuras/unidade13/unidade13_estabilidade_perda_ganho.png}
	\caption{Áreas de estabilidade, perda e ganho potencial de adequabilidade ambiental para \textit{Dinizia excelsa} sob cenários futuros.}
	\label{fig:unidade13_estabilidade}
\end{figure}

\section{Síntese tabular dos cenários}

\rscriptfile{scripts/unidade13/07_sintese_cenarios_futuros.R}

\begin{figure}[H]
	\centering
	\includegraphics[width=0.88\textwidth]{figuras/unidade13/unidade13_sintese_area_adequada.png}
	\caption{Síntese da área potencialmente adequada atual e futura para \textit{Dinizia excelsa} sob diferentes cenários SSP.}
	\label{fig:unidade13_sintese_area}
\end{figure}

\section{Pipeline completo}

\rscriptfile{scripts/unidade13/08_pipeline_projecoes_futuras.R}

\section{Síntese da unidade}

As projeções climáticas futuras permitem avaliar possíveis mudanças na adequabilidade ambiental de espécies sob diferentes trajetórias climáticas. Para \textit{Dinizia excelsa}, esse tipo de análise é especialmente relevante para discutir persistência climática, áreas prioritárias, estabilidade ambiental e riscos associados a cenários de aquecimento.

\begin{conceito}[title={\faBook\quad Mensagem central}]
	Projeções futuras em SDM devem ser interpretadas como cenários de adequabilidade ambiental, não como previsões determinísticas de ocorrência. A robustez da interpretação depende da qualidade dos dados atuais, da escolha dos cenários, da compatibilidade das variáveis e da avaliação das incertezas.
\end{conceito}

\section{Exercícios de fixação}

\begin{enumerate}
	\item Explique a diferença entre clima atual e clima futuro em SDM.
	\item O que são cenários SSP?
	\item Por que as variáveis futuras devem ter os mesmos nomes das variáveis usadas no ajuste atual?
	\item Como interpretar valores positivos e negativos de mudança de adequabilidade?
	\item O que são áreas de estabilidade climática?
	\item Por que projeções futuras não devem ser interpretadas como previsões absolutas?
	\item Como os resultados desta unidade podem apoiar decisões de conservação?
\end{enumerate}

% ============================================================
% UNIDADE 14
% ============================================================

\chapter[Transferência e extrapolação]{Transferência Espacial e Temporal, Extrapolação, MESS e MOP}

\begin{objetivos}
	Ao final desta unidade, o estudante deverá compreender os desafios da transferência espacial e temporal em SDM; diferenciar interpolação e extrapolação ambiental; calcular mapas MESS e MOP; identificar áreas futuras fora do domínio ambiental de calibração; avaliar riscos de extrapolação nas projeções climáticas de \textit{Dinizia excelsa}; e interpretar mapas futuros considerando incerteza ambiental.
\end{objetivos}

\section{Introdução}

Modelos de distribuição de espécies são frequentemente calibrados sob condições ambientais atuais e depois transferidos para outros espaços ou períodos, como cenários climáticos futuros, períodos paleoclimáticos ou regiões geográficas distintas. Essa transferência é uma das aplicações mais importantes de SDM, mas também uma das mais sensíveis a erros \parencite{peterson2011,araujo2012,yates2018transferability,zurell2020benchmarking}.

Quando um modelo é projetado para condições ambientais semelhantes às usadas na calibração, estamos diante de uma situação de interpolação ambiental. Por outro lado, quando o modelo é aplicado a condições ambientais fora do intervalo conhecido, ocorre extrapolação. Nesses casos, as predições podem ser estatisticamente possíveis, mas ecologicamente frágeis.

\begin{conceito}[title={\faBook\quad Interpolação versus extrapolação}]
	Interpolação ocorre quando o modelo é projetado para combinações ambientais semelhantes às observadas na calibração. Extrapolação ocorre quando o modelo é aplicado a valores ou combinações ambientais não representadas nos dados usados para ajustar o modelo.
\end{conceito}

\section{Transferência espacial e temporal}

A transferência espacial ocorre quando um modelo calibrado em uma região é aplicado a outra região. A transferência temporal ocorre quando um modelo calibrado no presente é aplicado ao passado ou futuro. Em ambos os casos, a confiabilidade depende da similaridade ambiental entre calibração e projeção.

\begin{atencao}[title={\faExclamationTriangle\quad Projeções futuras exigem diagnóstico de extrapolação}]
	Mapas futuros de adequabilidade não devem ser interpretados isoladamente. Eles precisam ser avaliados junto com mapas de similaridade ambiental, como MESS e MOP.
\end{atencao}

\section{MESS: Multivariate Environmental Similarity Surface}

O MESS avalia, para cada célula projetada, se seus valores ambientais estão dentro do intervalo observado na região de calibração. Valores positivos indicam similaridade ambiental; valores negativos indicam extrapolação univariada em pelo menos uma variável.

\rscriptfile{scripts/unidade14/01_estrutura_pastas.R}

\rscriptfile{scripts/unidade14/02_calcular_mess_futuro.R}

\begin{figure}[H]
	\centering
	\includegraphics[width=0.96\textwidth]{figuras/unidade14/unidade14_mess_cenarios.png}
	\caption{Mapas MESS para os cenários futuros de \textit{Dinizia excelsa}. Valores negativos indicam extrapolação ambiental em relação ao domínio de calibração.}
	\label{fig:unidade14_mess}
\end{figure}

\section{Variável limitante no MESS}

Além do valor de similaridade, é possível identificar qual variável ambiental é responsável pela maior extrapolação em cada célula. Essa análise ajuda a interpretar se a incerteza futura decorre principalmente de temperatura, precipitação ou sazonalidade.

\rscriptfile{scripts/unidade14/03_variavel_limitante_mess.R}

\begin{figure}[H]
	\centering
	\includegraphics[width=0.96\textwidth]{figuras/unidade14/unidade14_variavel_limitante_mess.png}
	\caption{Variável ambiental mais limitante no diagnóstico MESS para os cenários futuros.}
	\label{fig:unidade14_limitante}
\end{figure}

\section{MOP: Mobility-Oriented Parity}

O MOP avalia a similaridade multivariada entre ambientes de calibração e projeção, considerando a distância ambiental entre células. Diferentemente do MESS, que avalia limites univariados, o MOP considera combinações ambientais multivariadas e identifica áreas ambientalmente não análogas \parencite{owens2013constraints}.

\rscriptfile{scripts/unidade14/04_calcular_mop_futuro.R}

\begin{figure}[H]
	\centering
	\includegraphics[width=0.96\textwidth]{figuras/unidade14/unidade14_mop_cenarios.png}
	\caption{Mapas MOP simplificados para os cenários futuros. Valores mais altos indicam maior distância ambiental em relação ao domínio de calibração.}
	\label{fig:unidade14_mop}
\end{figure}

\section{Adequabilidade sob extrapolação}

Nesta etapa, os mapas de adequabilidade futura são combinados com os mapas MESS e MOP. O objetivo é identificar áreas projetadas como adequadas, mas sob condições extrapolativas.

\rscriptfile{scripts/unidade14/05_adequabilidade_extrapolacao.R}

\begin{figure}[H]
	\centering
	\includegraphics[width=0.96\textwidth]{figuras/unidade14/unidade14_adequabilidade_extrapolacao.png}
	\caption{Adequabilidade futura combinada com diagnóstico de extrapolação ambiental.}
	\label{fig:unidade14_adequabilidade_extrapolacao}
\end{figure}

\section{Síntese da extrapolação}

\rscriptfile{scripts/unidade14/06_sintese_extrapolacao.R}

\begin{figure}[H]
	\centering
	\includegraphics[width=0.88\textwidth]{figuras/unidade14/unidade14_sintese_extrapolacao.png}
	\caption{Proporção da área futura sob extrapolação ambiental em cada cenário SSP.}
	\label{fig:unidade14_sintese}
\end{figure}

\section{Pipeline completo}

\rscriptfile{scripts/unidade14/07_pipeline_transferencia_extrapolacao.R}

\section{Síntese da unidade}

A transferência espacial e temporal é essencial para aplicações de SDM em mudanças climáticas, mas exige diagnóstico rigoroso de extrapolação. Para \textit{Dinizia excelsa}, os mapas MESS e MOP permitem identificar regiões onde as projeções futuras devem ser interpretadas com maior cautela.

\begin{conceito}[title={\faBook\quad Mensagem central}]
	Projeções futuras só são ecologicamente defensáveis quando acompanhadas de diagnóstico de extrapolação. Áreas adequadas sob forte extrapolação devem ser interpretadas como hipóteses incertas, não como evidências robustas de persistência futura.
\end{conceito}

\section{Exercícios de fixação}

\begin{enumerate}
	\item Diferencie interpolação ambiental e extrapolação ambiental.
	\item Explique o que representa um valor negativo de MESS.
	\item Qual a diferença entre MESS e MOP?
	\item Por que modelos podem prever alta adequabilidade em áreas extrapolativas?
	\item Como mapas de extrapolação ajudam na interpretação de projeções futuras?
	\item Por que transferência temporal é especialmente sensível em cenários climáticos futuros?
	\item Como você usaria MESS e MOP para selecionar áreas prioritárias de conservação?
\end{enumerate}


% ============================================================
% UNIDADE 15
% ============================================================

\chapter[Paleoclima e nicho climático passado]{Paleoclima, Nicho Climático Passado e Reconstrução Histórica da Distribuição}

\begin{objetivos}
	Ao final desta unidade, o estudante deverá compreender como dados paleoclimáticos podem ser usados em SDM; reconhecer períodos como Holoceno Médio, Último Máximo Glacial e Último Interglacial; preparar camadas paleoclimáticas compatíveis com os modelos atuais; projetar modelos de \textit{Dinizia excelsa} para climas passados; avaliar estabilidade climática histórica; identificar possíveis áreas de refúgio; e interpretar o nicho climático passado como componente da história biogeográfica da espécie.
\end{objetivos}

\section{Introdução}

A distribuição atual de uma espécie arbórea amazônica é resultado de processos ecológicos contemporâneos e de uma história climática profunda. Mudanças passadas em temperatura, precipitação, sazonalidade, conectividade florestal e estabilidade ambiental podem ter influenciado a persistência, expansão, retração e estruturação espacial das populações.

Em Modelagem de Distribuição de Espécies, reconstruções paleoclimáticas permitem projetar modelos calibrados no presente para períodos passados, como o Holoceno Médio, o Último Máximo Glacial e o Último Interglacial. Essas projeções ajudam a formular hipóteses sobre estabilidade climática, refúgios, conectividade histórica e conservação de nicho.

\begin{conceito}[title={\faBook\quad Nicho climático passado}]
	O nicho climático passado representa o conjunto de condições ambientais históricas que poderiam ter favorecido a persistência, expansão ou retração de uma espécie ao longo do tempo geológico recente.
\end{conceito}

%------------------------------------------------
\section{Linha do tempo paleoclimática e cenários futuros}
%------------------------------------------------

A modelagem paleoclimática permite reconstruir a distribuição potencial das
espécies sob condições ambientais do passado, enquanto as projeções futuras
permitem avaliar possíveis respostas sob cenários de mudanças climáticas. Em
conjunto, essas abordagens conectam biogeografia histórica, estabilidade
climática, refúgios ecológicos e vulnerabilidade futura.

Nesta unidade, o nicho climático de \textit{Dinizia excelsa} é interpretado ao
longo de uma linha temporal que integra períodos paleoclimáticos clássicos,
condições climáticas atuais e cenários futuros derivados de modelos climáticos
globais. Essa perspectiva temporal amplia a interpretação dos modelos, pois
permite investigar se áreas atualmente adequadas também foram estáveis no
passado e se tendem a permanecer adequadas sob mudanças climáticas futuras.

\begin{center}
	\fbox{
		\begin{minipage}{0.88\textwidth}
			
			\centering
			\textbf{Linha do tempo climática utilizada na interpretação dos modelos}
			
			\vspace{0.4cm}
			
			\begin{tabular}{c c p{0.48\textwidth}}
				\toprule
				\textbf{Tempo} & \textbf{Período/Cenário} & \textbf{Interpretação ecológica} \\
				\midrule
				
				130 ka &
				Último Interglacial (LIG) &
				Período quente anterior ao último ciclo glacial, útil para avaliar respostas da espécie a condições mais quentes que as atuais. \\
				
				21 ka &
				Último Máximo Glacial (LGM) &
				Fase fria e seca do último ciclo glacial, associada a mudanças expressivas na distribuição de ambientes tropicais. \\
				
				6 ka &
				Holoceno Médio &
				Período pós-glacial com reorganização climática e ecológica, importante para inferir estabilidade recente de áreas adequadas. \\
				
				0 ka &
				Presente &
				Condição climática de referência usada para calibrar ou interpretar os modelos atuais. \\
				
				2050 &
				SSP126 &
				Cenário futuro de menor emissão, usado para avaliar respostas sob mitigação climática. \\
				
				2050 &
				SSP585 &
				Cenário futuro de alta emissão, representando maior pressão climática sobre a distribuição potencial. \\
				
				2070 &
				SSP126 &
				Projeção de médio-longo prazo sob cenário de menor emissão. \\
				
				2070 &
				SSP585 &
				Projeção de médio-longo prazo sob cenário de alta emissão e maior risco de perda de adequabilidade. \\
				
				\bottomrule
			\end{tabular}
			
		\end{minipage}
	}
\end{center}

Essa linha do tempo deve ser interpretada como uma estrutura conceitual para
comparar estabilidade, perda, ganho e deslocamento de áreas climaticamente
adequadas ao longo de diferentes janelas temporais. Áreas que permanecem
adequadas desde o passado até o futuro podem indicar maior estabilidade
climática e potencial importância como refúgios biogeográficos. Por outro lado,
áreas adequadas apenas no presente ou apenas em cenários futuros devem ser
interpretadas com cautela, especialmente quando associadas à extrapolação
ambiental.

\begin{center}
	\fbox{
		\begin{minipage}{0.90\textwidth}
			
			\textbf{Mensagem-chave.}
			
			A integração entre paleoclima, clima presente e cenários futuros transforma a
			SDM em uma ferramenta temporal. O objetivo deixa de ser apenas mapear onde a
			espécie pode ocorrer hoje e passa a incluir perguntas sobre estabilidade
			histórica, persistência climática, vulnerabilidade futura e identificação de
			possíveis refúgios para conservação.
			
		\end{minipage}
	}
\end{center}

\section{Períodos paleoclimáticos relevantes}

Nesta unidade, trabalharemos com três janelas temporais principais:

\begin{itemize}
	\item Holoceno Médio, aproximadamente 6 mil anos antes do presente;
	\item Último Máximo Glacial, aproximadamente 21 mil anos antes do presente;
	\item Último Interglacial, aproximadamente 120 mil anos antes do presente.
\end{itemize}

Esses períodos representam diferentes estados climáticos do sistema terrestre e permitem avaliar se áreas atualmente adequadas para \textit{Dinizia excelsa} também apresentavam adequabilidade no passado.

\begin{atencao}[title={\faExclamationTriangle\quad Projeções paleoclimáticas são hipóteses}]
	Mapas paleoclimáticos em SDM não demonstram ocorrência passada da espécie. Eles indicam áreas que poderiam ter sido climaticamente adequadas, assumindo conservação parcial do nicho e compatibilidade entre clima passado e modelo atual.
\end{atencao}

\section{Preparação das camadas paleoclimáticas}

As camadas paleoclimáticas devem ser harmonizadas com as variáveis atuais utilizadas no ajuste dos modelos. Isso inclui nomes, resolução, extensão, projeção e máscara espacial.

\rscriptfile{scripts/unidade15/01_estrutura_pastas.R}

\rscriptfile{scripts/unidade15/02_preparar_variaveis_paleoclimaticas.R}

\section{Projeção dos modelos para o passado}

Os modelos ajustados no presente são projetados para os cenários paleoclimáticos. Essa etapa permite comparar a adequabilidade atual com adequabilidade sob climas passados.

\rscriptfile{scripts/unidade15/03_projetar_modelos_passado.R}

\begin{figure}[H]
	\centering
	\includegraphics[width=0.98\textwidth]{figuras/unidade15/unidade15_modelos_passado_lgm.png}
	\caption{Projeções dos modelos individuais para \textit{Dinizia excelsa} no Último Máximo Glacial.}
	\label{fig:unidade15_modelos_lgm}
\end{figure}

\section{Ensemble paleoclimático}

As projeções individuais são combinadas em mapas ensemble para cada período paleoclimático.

\rscriptfile{scripts/unidade15/04_ensemble_passado.R}

\begin{figure}[H]
	\centering
	\includegraphics[width=0.96\textwidth]{figuras/unidade15/unidade15_ensemble_passado_periodos.png}
	\caption{Ensemble paleoclimático de adequabilidade ambiental para \textit{Dinizia excelsa}.}
	\label{fig:unidade15_ensemble_passado}
\end{figure}

\section{Trajetória do nicho climático}

A trajetória climática permite comparar o espaço ambiental ocupado no presente com o espaço climático disponível nos períodos passados. Nesta unidade, usamos PCA ambiental para representar mudanças no espaço climático.

\rscriptfile{scripts/unidade15/05_trajetoria_nicho_climatico.R}

\begin{figure}[H]
	\centering
	\includegraphics[width=0.88\textwidth]{figuras/unidade15/unidade15_pca_nicho_passado.png}
	\caption{Trajetória do espaço climático atual e paleoclimático para \textit{Dinizia excelsa}.}
	\label{fig:unidade15_pca_nicho_passado}
\end{figure}

\section{Estabilidade histórica e refúgios climáticos}

Áreas climaticamente adequadas no presente e também adequadas em períodos passados podem ser interpretadas como hipóteses de estabilidade climática histórica. Essas áreas não provam refúgios, mas indicam regiões candidatas à persistência de condições favoráveis.

\rscriptfile{scripts/unidade15/06_estabilidade_refugios_paleoclimaticos.R}

\begin{figure}[H]
	\centering
	\includegraphics[width=0.96\textwidth]{figuras/unidade15/unidade15_estabilidade_refugios.png}
	\caption{Áreas de estabilidade paleoclimática e possíveis refúgios climáticos para \textit{Dinizia excelsa}.}
	\label{fig:unidade15_refugios}
\end{figure}

\section{Síntese paleoclimática}

\rscriptfile{scripts/unidade15/07_sintese_paleoclimatica.R}

\begin{figure}[H]
	\centering
	\includegraphics[width=0.88\textwidth]{figuras/unidade15/unidade15_sintese_paleoclimatica.png}
	\caption{Síntese da área adequada estimada para \textit{Dinizia excelsa} em períodos paleoclimáticos.}
	\label{fig:unidade15_sintese}
\end{figure}

\section{Pipeline completo}

\rscriptfile{scripts/unidade15/08_pipeline_paleoclima.R}

\section{Síntese da unidade}

A reconstrução paleoclimática amplia a interpretação dos modelos de distribuição ao conectar o nicho atual da espécie com sua história ambiental. Para \textit{Dinizia excelsa}, esse tipo de análise permite discutir estabilidade climática, possíveis áreas de persistência histórica e hipóteses biogeográficas para a distribuição amazônica contemporânea.

\begin{conceito}[title={\faBook\quad Mensagem central}]
	Projeções paleoclimáticas não mostram onde a espécie ocorreu necessariamente no passado, mas indicam onde o clima passado poderia ter sido compatível com seu nicho climático atual.
\end{conceito}

\section{Exercícios de fixação}

\begin{enumerate}
	\item Explique a diferença entre clima atual, clima futuro e paleoclima em SDM.
	\item O que significa projetar um modelo atual para o Último Máximo Glacial?
	\item Por que áreas estáveis ao longo do tempo podem ser interpretadas como hipóteses de refúgio?
	\item Quais são as principais limitações das projeções paleoclimáticas?
	\item Como a PCA ajuda a interpretar mudanças no espaço climático passado?
	\item Por que conservação de nicho é uma hipótese importante nessa análise?
	\item Como resultados paleoclimáticos podem contribuir para conservação de espécies amazônicas?
\end{enumerate}

% ============================================================
% UNIDADE 16
% ============================================================

\chapter[Incertezas em SDM]{Incertezas em Modelagem de Distribuição de Espécies}

\begin{objetivos}
	Ao final desta unidade, o estudante deverá compreender as principais fontes de incerteza em SDM; diferenciar incerteza associada aos dados, algoritmos, cenários climáticos, extrapolação e escala espacial; quantificar incerteza algorítmica a partir de múltiplos modelos; avaliar incerteza entre cenários SSP; integrar incerteza de extrapolação ambiental; produzir mapas de incerteza para \textit{Dinizia excelsa}; e interpretar resultados de SDM de forma crítica, transparente e reprodutível.
\end{objetivos}

\section{Introdução}

Modelos de distribuição de espécies não produzem verdades absolutas sobre a distribuição das espécies. Eles produzem estimativas condicionadas aos dados de ocorrência, às variáveis ambientais, à área acessível, ao algoritmo, à parametrização, ao período climático e ao cenário de projeção. Portanto, toda aplicação robusta de SDM deve explicitar e quantificar incertezas \parencite{guisan2017habitat,franklin2023,zurell2020benchmarking,yates2018transferability}.

No caso de \textit{Dinizia excelsa}, as incertezas são particularmente importantes porque estamos trabalhando com uma espécie arbórea amazônica de grande porte, registros de ocorrência potencialmente enviesados, modelos climáticos, projeções futuras, reconstruções paleoclimáticas e extrapolação temporal.

\begin{conceito}[title={\faBook\quad Incerteza em SDM}]
	Incerteza é a variação ou limitação associada às estimativas do modelo. Em SDM, ela pode surgir dos dados de ocorrência, das variáveis ambientais, dos algoritmos, dos cenários climáticos, da extrapolação e da escala espacial.
\end{conceito}

\section{Fontes principais de incerteza}

Nesta unidade, organizamos as incertezas em cinco componentes principais:

\begin{itemize}
	\item incerteza dos dados de ocorrência;
	\item incerteza algorítmica;
	\item incerteza entre cenários climáticos;
	\item incerteza por extrapolação ambiental;
	\item incerteza espacial integrada.
\end{itemize}

\begin{atencao}[title={\faExclamationTriangle\quad Incerteza não é erro descartável}]
	Incerteza não deve ser escondida. Ela é parte fundamental da interpretação científica e ajuda a identificar onde as conclusões são robustas e onde devem ser tratadas com cautela.
\end{atencao}

\section{Incerteza algorítmica}

A incerteza algorítmica representa a divergência entre modelos ajustados com diferentes algoritmos. Nesta apostila, ela pode ser estimada pelo desvio-padrão das predições de GLM, GAM, Random Forest, BRT e Maxent.

\rscriptfile{scripts/unidade16/01_estrutura_pastas.R}

\rscriptfile{scripts/unidade16/02_incerteza_algoritmica_atual.R}

\begin{figure}[H]
	\centering
	\includegraphics[width=0.88\textwidth]{figuras/unidade16/unidade16_incerteza_algoritmica_atual.png}
	\caption{Incerteza algorítmica atual para \textit{Dinizia excelsa}, estimada pelo desvio-padrão das predições entre GLM, GAM, Random Forest, BRT e Maxent.}
	\label{fig:unidade16_incerteza_algoritmica}
\end{figure}

\section{Incerteza entre cenários futuros}

A incerteza entre cenários representa a variação nas projeções futuras associada às diferentes trajetórias SSP. Nesta unidade, comparamos os ensembles futuros sob SSP126, SSP245, SSP370 e SSP585.

\rscriptfile{scripts/unidade16/03_incerteza_cenarios_futuros.R}

\begin{figure}[H]
	\centering
	\includegraphics[width=0.88\textwidth]{figuras/unidade16/unidade16_incerteza_cenarios_futuros.png}
	\caption{Incerteza entre cenários futuros para \textit{Dinizia excelsa}, estimada pelo desvio-padrão entre ensembles SSP126, SSP245, SSP370 e SSP585.}
	\label{fig:unidade16_incerteza_cenarios}
\end{figure}

\section{Incerteza por extrapolação ambiental}

Mapas de extrapolação, como MESS e MOP, indicam áreas onde as condições futuras ou passadas estão fora do domínio ambiental de calibração. Essas regiões devem receber maior peso de incerteza interpretativa.

\rscriptfile{scripts/unidade16/04_incerteza_extrapolacao.R}

\begin{figure}[H]
	\centering
	\includegraphics[width=0.96\textwidth]{figuras/unidade16/unidade16_incerteza_extrapolacao.png}
	\caption{Incerteza associada à extrapolação ambiental futura para \textit{Dinizia excelsa}.}
	\label{fig:unidade16_incerteza_extrapolacao}
\end{figure}

\section{Incerteza paleoclimática}

As reconstruções paleoclimáticas possuem incertezas associadas aos modelos climáticos, à resolução espacial, à reconstrução das condições passadas e à hipótese de conservação de nicho. Nesta unidade, avaliamos a variação entre os ensembles paleoclimáticos.

\rscriptfile{scripts/unidade16/05_incerteza_paleoclimatica.R}

\begin{figure}[H]
	\centering
	\includegraphics[width=0.88\textwidth]{figuras/unidade16/unidade16_incerteza_paleoclimatica.png}
	\caption{Incerteza paleoclimática estimada pela variação entre períodos passados para \textit{Dinizia excelsa}.}
	\label{fig:unidade16_incerteza_paleo}
\end{figure}

\section{Incerteza integrada}

A incerteza integrada combina diferentes fontes em uma superfície única, permitindo identificar regiões onde as predições são mais robustas e regiões onde a interpretação deve ser mais cautelosa.

\rscriptfile{scripts/unidade16/06_incerteza_integrada.R}

\begin{figure}[H]
	\centering
	\includegraphics[width=0.88\textwidth]{figuras/unidade16/unidade16_incerteza_integrada.png}
	\caption{Mapa de incerteza integrada para \textit{Dinizia excelsa}, combinando incerteza algorítmica, incerteza entre cenários, extrapolação ambiental e incerteza paleoclimática.}
	\label{fig:unidade16_incerteza_integrada}
\end{figure}

\section{Adequabilidade versus incerteza}

Um resultado útil em conservação é identificar áreas com alta adequabilidade e baixa incerteza. Essas áreas representam regiões mais robustas para interpretação ecológica e priorização espacial.

\rscriptfile{scripts/unidade16/07_adequabilidade_incerteza.R}

\begin{figure}[H]
	\centering
	\includegraphics[width=0.96\textwidth]{figuras/unidade16/unidade16_adequabilidade_incerteza.png}
	\caption{Classificação combinada de adequabilidade e incerteza para \textit{Dinizia excelsa}.}
	\label{fig:unidade16_adequabilidade_incerteza}
\end{figure}

\section{Síntese das incertezas}

\rscriptfile{scripts/unidade16/08_sintese_incertezas.R}

\begin{figure}[H]
	\centering
	\includegraphics[width=0.88\textwidth]{figuras/unidade16/unidade16_sintese_incertezas.png}
	\caption{Síntese quantitativa das classes de adequabilidade e incerteza para \textit{Dinizia excelsa}.}
	\label{fig:unidade16_sintese}
\end{figure}

\section{Pipeline completo}

\rscriptfile{scripts/unidade16/09_pipeline_incertezas.R}

\section{Síntese da unidade}

A análise de incertezas transforma SDM de um produto cartográfico simples em uma ferramenta científica mais robusta. Para \textit{Dinizia excelsa}, a integração entre incerteza algorítmica, cenários futuros, extrapolação e paleoclima permite distinguir áreas de alta confiança de áreas onde os resultados devem ser tratados como hipóteses.

\begin{conceito}[title={\faBook\quad Mensagem central}]
	A interpretação de SDM deve sempre considerar duas perguntas simultâneas: onde o modelo indica alta adequabilidade e quão confiável é essa estimativa?
\end{conceito}

\section{Exercícios de fixação}

\begin{enumerate}
	\item Liste cinco fontes de incerteza em SDM.
	\item Explique o que é incerteza algorítmica.
	\item Por que cenários SSP diferentes produzem incerteza?
	\item Como MESS e MOP contribuem para interpretar incerteza?
	\item Por que projeções paleoclimáticas possuem incerteza própria?
	\item Como interpretar uma área com alta adequabilidade e alta incerteza?
	\item Como mapas de incerteza podem apoiar decisões de conservação?
\end{enumerate}

% ============================================================
% UNIDADE 17
% ============================================================

\chapter[Aplicações à conservação]{Aplicações de SDM à Conservação}

\begin{objetivos}
	Ao final desta unidade, o estudante deverá compreender como modelos de distribuição de espécies podem apoiar decisões de conservação; integrar adequabilidade atual, estabilidade futura, incerteza e áreas protegidas; identificar áreas prioritárias para \textit{Dinizia excelsa}; avaliar lacunas de proteção; produzir mapas de prioridade espacial; e discutir aplicações em espécies ameaçadas, árvores gigantes, restauração ecológica e planejamento sistemático da conservação.
\end{objetivos}

\section{Introdução}

Modelos de Distribuição de Espécies são amplamente utilizados para apoiar estratégias de conservação, especialmente em cenários de mudanças climáticas, perda de habitat, fragmentação florestal e expansão de pressões antrópicas. Em aplicações conservacionistas, o objetivo não é apenas mapear onde uma espécie pode ocorrer, mas identificar regiões onde a adequabilidade ambiental, a estabilidade climática, a baixa incerteza e a proteção territorial convergem \parencite{margules2000systematic,moilanen2009spatial,franklin2023,guisan2017habitat,zurell2020benchmarking}.

Nesta unidade, integramos os produtos construídos nas unidades anteriores para desenvolver uma análise aplicada à conservação de \textit{Dinizia excelsa}. O foco está na construção de mapas de prioridade espacial baseados em evidências ecológicas, climáticas e modelagem preditiva.

\begin{conceito}[title={\faBook\quad SDM aplicado à conservação}]
	Em conservação, SDM deve ser interpretado como uma camada de evidência espacial. Ele deve ser combinado com informações sobre proteção territorial, conectividade, ameaça, incerteza e viabilidade ecológica.
\end{conceito}

\section{Critérios de priorização}

Nesta unidade, usamos cinco critérios principais:

\begin{itemize}
	\item alta adequabilidade atual;
	\item estabilidade futura;
	\item estabilidade paleoclimática;
	\item baixa incerteza integrada;
	\item presença ou ausência em áreas protegidas.
\end{itemize}

\begin{atencao}[title={\faExclamationTriangle\quad Prioridade não é apenas adequabilidade}]
	Uma área altamente adequada pode não ser prioritária se apresentar alta incerteza, forte extrapolação climática ou ausência de viabilidade de conservação. A priorização deve integrar múltiplas dimensões.
\end{atencao}

\section{Preparação dos dados espaciais}

\rscriptfile{scripts/unidade17/01_estrutura_pastas.R}

\rscriptfile{scripts/unidade17/02_preparar_camadas_conservacao.R}

\section{Prioridade climática atual}

\rscriptfile{scripts/unidade17/03_prioridade_climatica_atual.R}

\begin{figure}[H]
	\centering
	\includegraphics[width=0.88\textwidth]{figuras/unidade17/unidade17_prioridade_climatica_atual.png}
	\caption{Prioridade climática atual para \textit{Dinizia excelsa}, baseada em adequabilidade atual e incerteza integrada.}
	\label{fig:unidade17_prioridade_atual}
\end{figure}

\section{Estabilidade futura e conservação}

\rscriptfile{scripts/unidade17/04_estabilidade_futura_conservacao.R}

\begin{figure}[H]
	\centering
	\includegraphics[width=0.96\textwidth]{figuras/unidade17/unidade17_estabilidade_futura_conservacao.png}
	\caption{Estabilidade futura de adequabilidade para \textit{Dinizia excelsa} sob cenários climáticos.}
	\label{fig:unidade17_estabilidade_futura}
\end{figure}

\section{Refúgios climáticos e estabilidade histórica}

\rscriptfile{scripts/unidade17/05_refugios_climaticos_conservacao.R}

\begin{figure}[H]
	\centering
	\includegraphics[width=0.96\textwidth]{figuras/unidade17/unidade17_refugios_climaticos.png}
	\caption{Áreas candidatas a refúgios climáticos para \textit{Dinizia excelsa}, integrando estabilidade atual e paleoclimática.}
	\label{fig:unidade17_refugios}
\end{figure}

\section{Lacunas de proteção}

\rscriptfile{scripts/unidade17/06_lacunas_protecao.R}

\begin{figure}[H]
	\centering
	\includegraphics[width=0.90\textwidth]{figuras/unidade17/unidade17_lacunas_protecao.png}
	\caption{Lacunas potenciais de proteção para \textit{Dinizia excelsa}. Áreas prioritárias fora de unidades de conservação indicam possíveis regiões estratégicas para conservação.}
	\label{fig:unidade17_lacunas}
\end{figure}

\section{Mapa integrado de prioridade}

\rscriptfile{scripts/unidade17/07_prioridade_integrada_conservacao.R}

\begin{figure}[H]
	\centering
	\includegraphics[width=0.90\textwidth]{figuras/unidade17/unidade17_prioridade_integrada.png}
	\caption{Mapa integrado de prioridade para conservação de \textit{Dinizia excelsa}.}
	\label{fig:unidade17_prioridade_integrada}
\end{figure}

\section{Síntese quantitativa}

\rscriptfile{scripts/unidade17/08_sintese_conservacao.R}

\begin{figure}[H]
	\centering
	\includegraphics[width=0.90\textwidth]{figuras/unidade17/unidade17_sintese_conservacao.png}
	\caption{Síntese quantitativa das classes de prioridade para conservação de \textit{Dinizia excelsa}.}
	\label{fig:unidade17_sintese}
\end{figure}

\section{Pipeline completo}

\rscriptfile{scripts/unidade17/09_pipeline_conservacao.R}

\section{Síntese da unidade}

A aplicação de SDM à conservação exige integração entre modelagem ecológica, mudanças climáticas, incerteza, estabilidade histórica e instrumentos territoriais de proteção. Para \textit{Dinizia excelsa}, a combinação desses elementos permite identificar áreas com maior potencial para conservação, monitoramento, restauração e planejamento estratégico.

\begin{conceito}[title={\faBook\quad Mensagem central}]
	A melhor aplicação conservacionista de SDM não é simplesmente mapear áreas adequadas, mas identificar áreas adequadas, estáveis, pouco incertas e estrategicamente relevantes para proteção da biodiversidade.
\end{conceito}

\section{Exercícios de fixação}

\begin{enumerate}
	\item Explique como SDM pode apoiar conservação de espécies arbóreas amazônicas.
	\item Por que áreas de alta adequabilidade e baixa incerteza são importantes?
	\item O que são lacunas de proteção?
	\item Como projeções futuras podem alterar prioridades de conservação?
	\item Por que estabilidade paleoclimática pode indicar possíveis refúgios?
	\item Como integrar SDM com áreas protegidas?
	\item Quais cuidados devem ser tomados antes de usar SDM em decisões públicas?
\end{enumerate}

\section{Síntese final da apostila}

Esta apostila encerra com uma mensagem central: Modelagem de Distribuição de Espécies não é apenas ajustar algoritmos. É um processo científico completo que conecta teoria ecológica, qualidade dos dados, hipóteses biogeográficas, variáveis ambientais, avaliação estatística, interpretação espacial, incerteza e conservação.

\begin{conceito}[title={\faBook\quad Mensagem final}]
	Um bom SDM não é aquele que apenas produz alto AUC. É aquele que possui dados confiáveis, hipóteses explícitas, área de calibração justificada, variáveis ecologicamente coerentes, modelos avaliados, incerteza comunicada e interpretação útil para ciência e conservação.
\end{conceito}

\section{Exercícios finais}

\begin{enumerate}
	\item Explique o fluxo completo de um projeto de SDM desde ocorrências até conservação.
	\item Por que a área acessível \(M\) influencia todo o pipeline?
	\item Compare os papéis de GLM, GAM, Random Forest e BRT no ensemble.
	\item Explique por que mapas de incerteza devem acompanhar mapas finais.
	\item Como projeções futuras podem apoiar decisões de conservação?
	\item Quais produtos finais devem ser entregues em um projeto reprodutível de SDM?
	\item Elabore um projeto de SDM para uma espécie amazônica de interesse ecológico ou conservacionista.
\end{enumerate}

% ============================================================
% UNIDADE 18
% ============================================================

%================================================
\chapter{Boas práticas, erros comuns e diretrizes internacionais para SDM}
%================================================

\markboth{Boas práticas em SDM}{Boas práticas em SDM}

\section{Introdução}

A Modelagem de Distribuição de Espécies (Species Distribution Modeling -- SDM)
evoluiu rapidamente nas últimas duas décadas, tornando-se uma das principais
ferramentas utilizadas em ecologia, biogeografia, conservação da biodiversidade,
planejamento ambiental e avaliação dos impactos das mudanças climáticas.

Paralelamente ao desenvolvimento de novos algoritmos e ao crescimento da
disponibilidade de dados ambientais e registros de ocorrência, surgiu uma
demanda crescente por maior rigor metodológico, transparência analítica e
reprodutibilidade científica. Atualmente, estudos de SDM não são avaliados
apenas pela sofisticação dos modelos utilizados, mas também pela qualidade dos
dados, clareza metodológica, robustez das inferências e disponibilidade dos
códigos e dados utilizados.

Diversos trabalhos demonstram que muitos dos erros observados em SDM não estão
relacionados aos algoritmos em si, mas às decisões tomadas durante a seleção de
variáveis ambientais, definição da área acessível ($M$), tratamento do viés
amostral, avaliação da extrapolação e interpretação inadequada das métricas de
desempenho \citep{guisan2000predictive, elith2006novel,
	barve2011crucial, guisan2017habitat, zurell2020standard}.

Ao mesmo tempo, periódicos científicos de alto impacto e organizações
internacionais vêm enfatizando a importância da ciência aberta, da
reprodutibilidade computacional e da adoção dos princípios FAIR para garantir
que os resultados produzidos possam ser verificados, reutilizados e integrados
em futuras sínteses ecológicas.

Esta unidade apresenta um conjunto de recomendações práticas destinadas a
auxiliar pesquisadores e estudantes na construção de modelos mais robustos,
transparentes e ecologicamente consistentes. O objetivo não é apenas evitar
erros comuns, mas promover uma visão crítica da modelagem ecológica como um
processo científico integrado que combina teoria ecológica, conhecimento
naturalístico, estatística, ciência de dados e tomada de decisão para a
conservação da biodiversidade.

\section{Por que modelos falham?}

Uma percepção equivocada frequentemente observada entre iniciantes em SDM é a
ideia de que modelos falham principalmente devido à escolha inadequada do
algoritmo. Embora diferentes algoritmos possuam características distintas,
estudos comparativos demonstram que fatores relacionados à qualidade dos dados,
à definição do problema ecológico e às decisões metodológicas exercem impacto
igual ou superior ao da escolha do algoritmo propriamente dito
\citep{elith2006novel, araujo2007ensemble, guisan2017habitat}.

Modelos podem apresentar elevados valores de AUC e TSS e, ainda assim,
produzir interpretações ecologicamente incorretas. Da mesma forma, modelos
estatisticamente simples podem fornecer inferências robustas quando baseados em
dados adequados, hipóteses bem definidas e procedimentos transparentes.

Nas seções seguintes são apresentados alguns dos erros mais frequentes
identificados na literatura especializada.

%------------------------------------------------
\section{Erro 1 - Utilizar todas as variáveis do WorldClim sem seleção prévia}

Um erro muito comum consiste em utilizar automaticamente todas as variáveis
bioclimáticas disponíveis, especialmente as 19 variáveis do WorldClim, sem
avaliar redundância estatística ou relevância ecológica. Embora essas variáveis
representem diferentes componentes do clima, muitas delas são fortemente
correlacionadas entre si.

A inclusão indiscriminada de variáveis ambientais pode aumentar a
multicolinearidade, dificultar a interpretação ecológica dos coeficientes,
favorecer o sobreajuste e reduzir a capacidade de transferência espacial e
temporal dos modelos \citep{dormann2013collinearity, guisan2017habitat}.

\begin{sdmwarningbox}
	Utilizar muitas variáveis ambientais sem seleção prévia pode gerar
	multicolinearidade, sobreajuste, instabilidade nas projeções futuras,
	baixa transferibilidade e interpretações ecológicas inconsistentes.
\end{sdmwarningbox}

Boas práticas incluem a avaliação da correlação entre variáveis, o cálculo do
Fator de Inflação da Variância (VIF), o uso criterioso de PCA quando
apropriado e, principalmente, a seleção de preditores com base em hipóteses
ecológicas relacionadas à espécie e à escala do estudo.

%------------------------------------------------
\section{Erro 2 - Ignorar a área acessível \texorpdfstring{($M$)}{M}}

A definição da área acessível ($M$) é uma das decisões mais importantes em
SDM. No contexto do \textit{BAM framework}, a distribuição observada de uma
espécie resulta da interação entre condições abióticas adequadas, interações
bióticas e áreas historicamente acessíveis por dispersão
\citep{soberon2005interpretation, soberon2009niches, barve2011crucial}.

Ignorar a área $M$ pode levar à seleção de background ou pseudo-ausências em
regiões que a espécie nunca poderia ter alcançado, confundindo ausência de
dispersão com inadequação ambiental.

\begin{sdmwarningbox}
	\textbf{Consequências principais.}
	
	A ausência de uma área $M$ ecologicamente justificada pode gerar
	pseudo-ausências irreais, background inadequado, viés na calibração do modelo
	e superestimação da distribuição potencial.
\end{sdmwarningbox}

A área $M$ deve ser definida com base no conhecimento biogeográfico da espécie,
barreiras geográficas, história de dispersão, limites de biomas, bacias
hidrográficas, ecorregiões ou outras unidades espacialmente coerentes.

%------------------------------------------------
\section{Erro 3 - Escolher o melhor modelo apenas pelo AUC}

O AUC é uma das métricas mais utilizadas em SDM por medir a capacidade
discriminatória do modelo. No entanto, seu uso isolado pode ser problemático,
pois altos valores de AUC não garantem que o modelo seja ecologicamente
plausível ou adequado para projeções espaciais e temporais
\citep{lobo2008auc, allouche2006assessing}.

Modelos com elevado AUC podem apresentar curvas de resposta biologicamente
incoerentes, forte dependência de viés amostral ou extrapolações inadequadas.
Assim, a avaliação do modelo deve combinar métricas estatísticas e interpretação
ecológica.

\begin{center}
	\fbox{
		\begin{minipage}{0.90\textwidth}
			
			\textbf{Consequências principais.}
			
			Escolher modelos apenas pelo maior AUC pode gerar falsa sensação de desempenho,
			supervalorização de pequenas diferenças estatísticas e interpretações
			ecológicas equivocadas.
			
		\end{minipage}
	}
\end{center}

Boas práticas incluem avaliar simultaneamente AUC, TSS, Índice de Boyce,
matrizes de confusão, curvas de resposta, mapas de adequabilidade e coerência
biogeográfica das predições.

%------------------------------------------------
\section{Erro 4 - Projetar para o futuro sem avaliar extrapolação}

Projeções futuras assumem que a relação estimada entre ocorrência e ambiente
no presente permanece válida sob novas condições climáticas. Entretanto,
cenários futuros podem conter combinações ambientais inexistentes no período de
calibração, resultando em extrapolação ambiental.

Quando a extrapolação não é avaliada, o modelo pode produzir previsões
aparentemente precisas, mas biologicamente frágeis, especialmente em regiões
com climas futuros muito diferentes dos climas atuais
\citep{elith2010functional, owens2013constraints, zurell2020standard}.

\begin{center}
	\fbox{
		\begin{minipage}{0.90\textwidth}
			
			\textbf{Consequências principais.}
			
			Projetar modelos para o futuro sem avaliar extrapolação pode gerar previsões
			biologicamente absurdas, mapas instáveis e interpretações equivocadas de perda
			ou ganho de áreas adequadas.
			
		\end{minipage}
	}
\end{center}

Boas práticas incluem calcular MESS, MOP, mapas de novidade ambiental e
identificar explicitamente áreas onde as projeções estão fora do domínio
ambiental observado durante a calibração.

%------------------------------------------------
\section{Erro 5 - Confundir adequabilidade ambiental com probabilidade de ocorrência}

Modelos de distribuição geralmente produzem índices de adequabilidade
ambiental, isto é, valores que expressam o grau de compatibilidade entre as
condições ambientais de um local e os registros utilizados na calibração do
modelo. Esses valores não devem ser interpretados automaticamente como
probabilidade real de ocorrência da espécie.

A ocorrência efetiva depende também de dispersão, interações bióticas,
história biogeográfica, dinâmica populacional, uso da paisagem e limitações de
amostragem \citep{soberon2007grinnellian, franklin2010mapping,
	peterson2011ecological}.

\begin{sdmwarningbox}
	\textbf{Consequência principal.}
	
	Confundir adequabilidade ambiental com probabilidade real de ocorrência pode
	levar à superestimação da distribuição da espécie e à interpretação incorreta
	de mapas de adequabilidade como mapas absolutos de presença e ausência.
\end{sdmwarningbox}

A interpretação recomendada é tratar os mapas contínuos como gradientes de
suporte ambiental ao modelo, evitando inferências absolutas sem validação
independente ou conhecimento ecológico adicional.

%------------------------------------------------
\section{Erro 6 - Ignorar o viés amostral}

Registros de ocorrência raramente representam amostras aleatórias da
distribuição das espécies. Em muitos casos, estão concentrados próximos a
estradas, rios navegáveis, cidades, centros de pesquisa, unidades de
conservação ou regiões historicamente mais inventariadas.

Quando esse viés não é tratado, o modelo pode aprender o padrão espacial do
esforço de coleta em vez da distribuição ecológica da espécie
\citep{phillips2009sample, kramer2013importance, fourcade2014mapping}.

\begin{center}
	\fbox{
		\begin{minipage}{0.90\textwidth}
			
			\textbf{Consequência principal.}
			
			Ignorar o viés amostral pode fazer com que o modelo reflita onde os
			pesquisadores coletaram mais dados, e não onde a espécie realmente encontra
			condições ambientais adequadas.
			
		\end{minipage}
	}
\end{center}

Boas práticas incluem rarefação espacial, filtros por distância mínima,
avaliação do esforço amostral, uso de background enviesado quando apropriado e
interpretação cautelosa de regiões subamostradas.

\section{Da qualidade metodológica à reprodutibilidade científica}

Evitar erros metodológicos representa apenas uma parte do processo de produção
de conhecimento científico confiável. Mesmo modelos ecologicamente robustos
podem tornar-se difíceis de verificar ou reutilizar quando os procedimentos
analíticos não são adequadamente documentados.

Por esse motivo, a comunidade científica internacional vem enfatizando a
necessidade de incorporar princípios de reprodutibilidade e ciência aberta ao
fluxo de trabalho em SDM. Essa abordagem busca garantir que os resultados
obtidos possam ser reproduzidos, auditados e reutilizados por outros
pesquisadores, fortalecendo a transparência e a credibilidade da pesquisa
ecológica.

%------------------------------------------------
\section{Reprodutibilidade científica em Modelagem de Distribuição de Espécies}
%------------------------------------------------

\subsection{Por que a reprodutibilidade é importante?}

A reprodutibilidade científica constitui um dos pilares fundamentais do método
científico moderno. Um estudo é considerado reprodutível quando outros
pesquisadores conseguem reproduzir os mesmos resultados utilizando os mesmos
dados, métodos analíticos e procedimentos computacionais descritos pelos autores.

Em Modelagem de Distribuição de Espécies (SDM), a reprodutibilidade possui
importância ainda maior devido à grande quantidade de etapas envolvidas no fluxo
de trabalho, incluindo aquisição de dados, filtragem espacial, seleção de
variáveis ambientais, definição da área acessível ($M$), ajuste de modelos,
avaliação estatística, projeções climáticas e análise de incertezas.

Pequenas diferenças nas versões dos softwares, bibliotecas ou algoritmos podem
produzir resultados distintos, dificultando a comparação entre estudos e a
verificação independente das conclusões. Por esse motivo, periódicos
internacionais e iniciativas de ciência aberta recomendam que todos os
procedimentos computacionais sejam documentados e disponibilizados de forma
transparente \citep{wilkinson2016fair, peng2011reproducible,
	sandve2013ten, zurell2020standard}.

%------------------------------------------------
\subsection{Controle do ambiente computacional}
%------------------------------------------------

Uma etapa fundamental para garantir reprodutibilidade consiste em registrar as
versões do software e dos pacotes utilizados durante a análise.

No ambiente R, a função \texttt{sessionInfo()} fornece informações detalhadas
sobre o sistema operacional, versão do R e bibliotecas carregadas durante a
sessão.

\begin{rcode}
	sessionInfo()
\end{rcode}

A saída dessa função deve ser arquivada juntamente com os scripts analíticos,
permitindo que outros pesquisadores reconstruam o ambiente computacional
original.

\begin{sdmpracticebox}
	Inclua sempre a saída de \texttt{sessionInfo()} nos materiais suplementares,
	repositórios GitHub ou documentos de reprodutibilidade do projeto.
\end{sdmpracticebox}

%------------------------------------------------
\subsection{Gerenciamento de dependências}
%------------------------------------------------

Ao longo do desenvolvimento de um projeto, versões de pacotes podem ser
atualizadas, modificadas ou removidas. Isso pode dificultar a reprodução futura
das análises.

O pacote \texttt{renv} permite criar ambientes isolados contendo exatamente as
versões utilizadas durante o desenvolvimento do estudo.

Após instalar os pacotes necessários, recomenda-se criar um snapshot do
ambiente:

\begin{rcode}
	renv::snapshot()
\end{rcode}

Esse comando gera um arquivo denominado \texttt{renv.lock}, que registra todas
as dependências do projeto.

Posteriormente, outro pesquisador poderá reconstruir exatamente o mesmo
ambiente utilizando:

\begin{rcode}
	renv::restore()
\end{rcode}

Essa abordagem reduz significativamente problemas associados à atualização de
bibliotecas e aumenta a confiabilidade da reprodução dos resultados.

%------------------------------------------------
\subsection{Controle de versão com Git}
%------------------------------------------------

O controle de versão permite registrar todas as modificações realizadas em um
projeto ao longo do tempo.

O sistema Git tornou-se o padrão internacional para gerenciamento de código
científico, oferecendo rastreabilidade, documentação histórica e colaboração
entre pesquisadores.

A criação de um repositório local pode ser iniciada com os seguintes comandos:

\begin{shellcode}
	git init
	git add .
	git commit -m "Primeira versão"
\end{shellcode}

Cada nova modificação relevante pode ser registrada por meio de novos
\textit{commits}, criando um histórico completo do desenvolvimento do projeto.

As principais vantagens incluem:

\begin{itemize}
	\item rastreamento de alterações;
	\item recuperação de versões anteriores;
	\item colaboração entre pesquisadores;
	\item documentação automática do fluxo de trabalho;
	\item integração com plataformas de ciência aberta.
\end{itemize}

%------------------------------------------------
\subsection{Disponibilização pública dos dados e scripts}
%------------------------------------------------

A transparência científica exige que dados, scripts e resultados sejam
disponibilizados sempre que possível.

Atualmente existem diversas plataformas especializadas para armazenamento,
compartilhamento e preservação de dados científicos.

\begin{table}[H]
	\centering
	\caption{Principais plataformas utilizadas para compartilhamento de dados e códigos científicos.}
	\label{tab:repositorios_ciencia_aberta}
	
	\begin{tabular}{
			p{0.20\textwidth}
			p{0.70\textwidth}
		}
		
		\toprule
		
		\textbf{Plataforma} &
		\textbf{Finalidade principal}
		
		\\
		\midrule
		
		GitHub &
		Versionamento de código, documentação e desenvolvimento colaborativo.
		
		\\
		
		Zenodo &
		Arquivamento permanente de dados e geração automática de DOI.
		
		\\
		
		Dryad &
		Repositório especializado para conjuntos de dados científicos.
		
		\\
		
		Figshare &
		Compartilhamento de dados, figuras, tabelas e materiais suplementares.
		
		\\
		
		\bottomrule
		
	\end{tabular}
	
\end{table}

Uma prática amplamente recomendada consiste em desenvolver o projeto em um
repositório GitHub e, posteriormente, arquivar uma versão estável no Zenodo,
obtendo um DOI permanente para citação.

%------------------------------------------------
\subsection{Princípios FAIR}
%------------------------------------------------

Em 2016 foi proposto um conjunto de diretrizes conhecido como FAIR Principles,
que se tornou referência internacional para gestão e compartilhamento de dados
científicos \citep{wilkinson2016fair}.

FAIR representa os seguintes princípios:

\begin{table}[H]
	\centering
	\caption{Princípios FAIR para gestão de dados científicos.}
	\label{tab:fair_principles}
	
	\begin{tabular}{
			p{0.15\textwidth}
			p{0.75\textwidth}
		}
		
		\toprule
		
		\textbf{Princípio} &
		\textbf{Descrição}
		
		\\
		\midrule
		
		Findable &
		Dados devem ser facilmente encontrados por humanos e sistemas computacionais.
		
		\\
		
		Accessible &
		Dados devem estar acessíveis por protocolos padronizados.
		
		\\
		
		Interoperable &
		Dados devem utilizar formatos compatíveis e padronizados.
		
		\\
		
		Reusable &
		Dados devem conter documentação suficiente para reutilização futura.
		
		\\
		
		\bottomrule
		
	\end{tabular}
	
\end{table}

A adoção dos princípios FAIR aumenta a transparência, a reutilização dos dados
e o impacto científico dos estudos.

%------------------------------------------------
\subsection{Checklist de reprodutibilidade}
%------------------------------------------------

Antes da publicação de um estudo de SDM, recomenda-se verificar os seguintes
itens:

\begin{itemize}
	\item Os scripts foram disponibilizados?
	\item As versões dos pacotes foram registradas?
	\item O ambiente computacional foi documentado?
	\item Os dados estão disponíveis em repositório público?
	\item O projeto possui DOI?
	\item Os procedimentos metodológicos podem ser reproduzidos integralmente?
	\item Os arquivos possuem metadados adequados?
	\item Os princípios FAIR foram considerados?
\end{itemize}

%------------------------------------------------
\subsection{Mensagem final}
%------------------------------------------------

A reprodutibilidade não deve ser vista como uma etapa adicional do projeto,
mas como parte integrante do processo científico. Em estudos de Modelagem de
Distribuição de Espécies, a documentação adequada dos dados, códigos,
dependências computacionais e decisões metodológicas aumenta a transparência,
fortalece a credibilidade dos resultados e permite que o conhecimento produzido
seja efetivamente reutilizado pela comunidade científica.

\begin{sdmkeybox}
	Um modelo de distribuição somente pode ser considerado plenamente
	científico quando seus resultados podem ser reproduzidos de forma
	independente por outros pesquisadores utilizando os mesmos dados,
	métodos e procedimentos descritos no estudo.
\end{sdmkeybox}

%------------------------------------------------
\section{Síntese dos principais erros e boas práticas}

\begin{table}[H]
	\centering
	\caption{Principais erros em SDM, consequências e boas práticas recomendadas.}
	\label{tab:erros_comuns_sdm}
	
	\scriptsize
	
	\begin{tabular}{
			p{0.24\textwidth}
			p{0.34\textwidth}
			p{0.34\textwidth}
		}
		
		\toprule
		\textbf{Erro comum} &
		\textbf{Consequências} &
		\textbf{Boas práticas} \\
		\midrule
		
		Usar todas as variáveis do WorldClim &
		Multicolinearidade, sobreajuste, projeções instáveis e baixa transferibilidade &
		Correlação, VIF, PCA e seleção ecológica de variáveis \\
		
		Ignorar a área acessível ($M$) &
		Pseudo-ausências irreais, background inadequado e viés de calibração &
		Delimitar $M$ com base em biogeografia, dispersão e barreiras históricas \\
		
		Escolher modelo apenas pelo AUC &
		Falsa sensação de desempenho e interpretações equivocadas &
		Combinar AUC, TSS, Boyce, curvas de resposta e plausibilidade ecológica \\
		
		Projetar para o futuro sem avaliar extrapolação &
		Previsões biologicamente frágeis e extrapolação climática não detectada &
		Usar MESS, MOP e mapas de novidade ambiental \\
		
		Confundir adequabilidade com probabilidade de ocorrência &
		Superestimação da ocupação real e interpretações incorretas &
		Interpretar mapas como gradientes de adequabilidade ambiental \\
		
		Ignorar viés amostral &
		Modelos refletem esforço de coleta, e não distribuição real &
		Rarefação espacial, filtros por distância e background adequado \\
		
		\bottomrule
	\end{tabular}
	
\end{table}

\section{Síntese da unidade}

Ao longo desta unidade foram discutidos alguns dos principais desafios
metodológicos enfrentados em estudos de Modelagem de Distribuição de Espécies.
Foi demonstrado que a qualidade de um modelo depende não apenas do algoritmo
utilizado, mas também da definição adequada da área acessível, da seleção de
variáveis ambientais, do controle do viés amostral, da avaliação da
extrapolação e da interpretação crítica das métricas de desempenho.

Também foram apresentados os fundamentos da reprodutibilidade científica,
controle de ambiente computacional, gerenciamento de dependências, controle de
versão, ciência aberta e princípios FAIR. Esses elementos constituem a base das
boas práticas atualmente recomendadas por periódicos internacionais,
organizações científicas e iniciativas globais de compartilhamento de dados.

Mais do que produzir mapas de adequabilidade ambiental, a modelagem moderna
busca gerar conhecimento transparente, verificável e reutilizável, capaz de
contribuir efetivamente para a compreensão da biodiversidade e para a tomada de
decisões em conservação sob cenários de rápidas mudanças ambientais.

%------------------------------------------------
\section{Mensagem final}

A maioria dos erros em SDM não está relacionada ao algoritmo utilizado, mas às
decisões tomadas antes e depois do ajuste do modelo. A preparação dos dados, a
definição da área acessível, a seleção de variáveis, o controle do viés
amostral, a avaliação da extrapolação e a interpretação ecológica dos
resultados são etapas tão importantes quanto o ajuste estatístico.

\begin{sdmkeybox}
	Um modelo simples, ecologicamente bem fundamentado e cuidadosamente avaliado
	pode ser mais confiável do que um modelo complexo ajustado com muitos
	preditores, área de calibração inadequada, viés amostral não tratado e
	interpretação exclusivamente baseada em métricas estatísticas.
\end{sdmkeybox}

%================================================
\chapter[Considerações finais]{Considerações Finais}
%================================================

\begin{objetivos}
	Ao final desta unidade, o estudante deverá ser capaz de reconhecer a Modelagem de Distribuição de Espécies como um campo integrador entre ecologia, biogeografia, estatística, ciência de dados, mudanças climáticas e conservação; compreender as principais decisões metodológicas discutidas ao longo da apostila; avaliar criticamente os limites e potencialidades dos SDMs; e estruturar interpretações ecológicas responsáveis a partir de modelos espaciais.
\end{objetivos}

\section{Síntese geral da apostila}

Esta apostila foi construída com o objetivo de apresentar a Modelagem de Distribuição de Espécies em R como um fluxo científico completo, crítico e reprodutível. Ao longo das unidades, partimos dos fundamentos conceituais do nicho ecológico e avançamos progressivamente para etapas aplicadas envolvendo dados reais de ocorrência, variáveis ambientais, seleção de preditores, algoritmos estatísticos, aprendizado de máquina, modelos presença-background, ensemble, avaliação, projeções futuras, paleoclima, extrapolação, incertezas e conservação.

O exemplo aplicado com \textit{Dinizia excelsa} permitiu demonstrar que um SDM não deve ser entendido apenas como uma técnica para produzir mapas de adequabilidade ambiental. Ao contrário, ele representa um processo analítico que conecta hipóteses ecológicas, dados espaciais, teoria do nicho, decisões metodológicas, avaliação estatística e interpretação biogeográfica. Nesse sentido, a modelagem de distribuição deve ser vista como uma ferramenta de inferência ecológica espacial, e não como uma simples rotina computacional.

\begin{conceito}[title={\faBook\quad Mensagem central da apostila}]
	Um bom SDM não é definido apenas pela beleza do mapa ou pelo valor de uma métrica. Ele depende da coerência entre pergunta ecológica, qualidade dos dados, definição da área acessível, seleção das variáveis, escolha dos algoritmos, avaliação independente, análise de incertezas e interpretação ecológica responsável.
\end{conceito}

\section{Do nicho ecológico ao espaço geográfico}

Um dos principais aprendizados desta apostila é que a distribuição geográfica de uma espécie não é determinada apenas pelo clima. A presença de uma espécie em determinado local resulta da interação entre condições abióticas, interações bióticas, acessibilidade histórica, dispersão, filtros ambientais, dinâmica populacional e história evolutiva. O arcabouço BAM, discutido nas unidades iniciais, reforça que a distribuição observada é apenas uma parte da distribuição potencial e que a área acessível \(M\) possui papel central na calibração de modelos.

A distinção entre espaço ambiental e espaço geográfico é fundamental. O modelo é ajustado a partir de relações no espaço ambiental, mas seus resultados são projetados no espaço geográfico. Essa transição exige cautela, porque padrões espaciais podem refletir tanto mecanismos ecológicos quanto estrutura amostral, autocorrelação, viés de coleta ou disponibilidade ambiental. Portanto, interpretar mapas de adequabilidade requer sempre retornar à pergunta ecológica original.

\section{Dados de ocorrência e responsabilidade analítica}

A qualidade dos dados de ocorrência é um dos pilares de qualquer SDM. Registros provenientes de inventários, herbários, bancos de dados públicos e plataformas como GBIF possuem enorme valor científico, mas também carregam limitações. Erros taxonômicos, coordenadas imprecisas, duplicatas, registros fora da área de estudo e viés espacial de amostragem podem influenciar diretamente os resultados.

Por isso, a limpeza dos dados não deve ser tratada como etapa burocrática. Ela é parte essencial da inferência ecológica. A rarefação espacial, a verificação taxonômica, a remoção de registros inconsistentes e a documentação das decisões tomadas são procedimentos fundamentais para aumentar a confiabilidade dos modelos.

\begin{atencao}[title={\faExclamationTriangle\quad Dados bons antes de modelos complexos}]
	Nenhum algoritmo compensa dados ruins, área acessível mal definida ou variáveis ambientais incoerentes. Em SDM, a qualidade da inferência começa antes do ajuste do modelo.
\end{atencao}

\section{Variáveis ambientais e interpretação ecológica}

As variáveis ambientais são a ponte entre a ecologia da espécie e a estrutura estatística do modelo. Entretanto, uma variável ambiental não é automaticamente um mecanismo ecológico. Uma variável pode apresentar alta importância no modelo porque resume um gradiente ambiental relevante, porque está correlacionada com outra variável causal, ou porque captura indiretamente processos não medidos.

Ao longo da apostila, enfatizamos a importância de selecionar variáveis com base em critérios ecológicos e estatísticos. A multicolinearidade, avaliada por correlação, VIF e PCA, deve ser tratada com cuidado, especialmente em modelos interpretativos. A seleção de variáveis deve buscar equilíbrio entre parcimônia, representatividade ambiental e plausibilidade ecológica.

\section{Algoritmos como diferentes lentes analíticas}

GLM, GAM, Random Forest, BRT e Maxent representam diferentes formas de aprender relações entre ocorrência e ambiente. GLM oferece maior simplicidade e interpretabilidade; GAM permite respostas suaves e não lineares; Random Forest captura padrões complexos e interações; BRT combina aprendizado sequencial e flexibilidade; Maxent é especialmente relevante para dados presença-background e presença-apenas.

Nenhum algoritmo é universalmente superior. Cada um possui pressupostos, vantagens e riscos. Modelos simples podem ser mais interpretáveis, mas menos flexíveis. Modelos complexos podem melhorar o desempenho preditivo, mas aumentar o risco de sobreajuste. Por isso, a comparação entre algoritmos deve considerar métricas, curvas de resposta, plausibilidade ecológica, estabilidade espacial e capacidade de transferência.

\begin{conceito}[title={\faBook\quad Algoritmos são hipóteses operacionais}]
	Cada algoritmo representa uma hipótese operacional diferente sobre como a espécie responde ao ambiente. Compará-los permite avaliar a robustez dos padrões e não apenas escolher automaticamente o maior valor de AUC.
\end{conceito}

\section{Avaliação, limiares e mapas binários}

A avaliação de modelos é uma etapa indispensável, mas nenhuma métrica isolada é suficiente. AUC, TSS, Kappa, sensibilidade, especificidade, acurácia e Boyce Index respondem a perguntas diferentes. O uso combinado dessas métricas permite uma avaliação mais completa, mas ainda assim deve ser interpretado com cautela.

A conversão de mapas contínuos em mapas binários pode ser útil para análises de área adequada, estabilidade, perda, ganho e conservação. No entanto, mapas binários dependem fortemente do limiar escolhido. Por isso, eles devem ser usados como instrumentos analíticos auxiliares, e não como representação definitiva da distribuição da espécie.

\section{Projeções futuras e mudanças climáticas}

As projeções futuras permitem avaliar como a adequabilidade ambiental pode mudar sob diferentes cenários climáticos. Nesta apostila, os cenários SSP foram apresentados como trajetórias alternativas, não como previsões absolutas. Isso é essencial para evitar interpretações determinísticas.

Modelos projetados para o futuro carregam incertezas associadas aos dados atuais, aos algoritmos, aos modelos climáticos, aos cenários de emissão, à resolução espacial e à extrapolação ambiental. Assim, mapas futuros devem ser interpretados como hipóteses condicionais: se a relação espécie-ambiente permanecer aproximadamente válida e se determinado cenário climático ocorrer, então certas regiões poderão ganhar, perder ou manter adequabilidade ambiental.

\section{Paleoclima e história biogeográfica}

A inclusão do paleoclima amplia a interpretação dos SDMs ao conectar a distribuição atual da espécie a condições ambientais passadas. Projeções para o Holoceno Médio, Último Máximo Glacial e Último Interglacial permitem formular hipóteses sobre estabilidade climática, refúgios históricos e persistência ambiental.

Essas análises são especialmente relevantes para espécies arbóreas amazônicas de longa vida, cujas distribuições contemporâneas podem refletir não apenas o clima atual, mas também legados históricos, dispersão limitada e estabilidade ambiental ao longo do tempo. Entretanto, projeções paleoclimáticas não demonstram ocorrência passada. Elas indicam áreas climaticamente compatíveis com o nicho estimado no presente.

\section{Transferência, extrapolação e incerteza}

A transferência espacial e temporal é uma das aplicações mais poderosas e, ao mesmo tempo, mais frágeis dos SDMs. Quando modelos são projetados para climas futuros ou passados, eles podem encontrar combinações ambientais não representadas na calibração. Nesses casos, as predições tornam-se extrapolativas.

Ferramentas como MESS e MOP ajudam a identificar áreas onde as projeções devem ser interpretadas com cautela. A integração de incerteza algorítmica, incerteza entre cenários, extrapolação e incerteza paleoclimática permite construir mapas mais transparentes e úteis para tomada de decisão.

\begin{atencao}[title={\faExclamationTriangle\quad Alta adequabilidade não significa alta confiança}]
	Uma área pode apresentar alta adequabilidade prevista e, ao mesmo tempo, alta incerteza ou forte extrapolação ambiental. Nessas situações, a interpretação deve ser cautelosa.
\end{atencao}

\section{Aplicações à conservação}

Em conservação, SDMs devem ser usados como camadas de evidência, não como respostas finais. A priorização espacial deve combinar adequabilidade atual, estabilidade futura, estabilidade paleoclimática, baixa incerteza, conectividade, proteção territorial, viabilidade ecológica e conhecimento de campo.

Para espécies arbóreas amazônicas como \textit{Dinizia excelsa}, essa abordagem pode apoiar inventários direcionados, planejamento de áreas prioritárias, restauração ecológica, monitoramento de populações, avaliação de lacunas de proteção e estratégias de conservação sob mudanças climáticas. O valor aplicado do SDM aumenta quando seus resultados são integrados a dados de campo, conhecimento local, políticas públicas e planejamento territorial.

\section{Boas práticas para estudos de SDM}

Com base no fluxo desenvolvido nesta apostila, recomenda-se que estudos de SDM sigam alguns princípios fundamentais:

\begin{enumerate}
	\item formular uma pergunta ecológica clara antes da modelagem;
	\item revisar a taxonomia e a qualidade dos registros de ocorrência;
	\item definir explicitamente a área acessível \(M\);
	\item escolher variáveis ambientais com justificativa ecológica;
	\item avaliar e reduzir multicolinearidade quando necessário;
	\item comparar algoritmos com diferentes níveis de complexidade;
	\item usar validação independente ou particionamento adequado dos dados;
	\item interpretar métricas junto com mapas e curvas de resposta;
	\item avaliar extrapolação em projeções espaciais ou temporais;
	\item quantificar incertezas;
	\item documentar todos os scripts e decisões metodológicas;
	\item comunicar limitações de forma transparente.
\end{enumerate}

\section{Limitações e perspectivas}

Apesar do avanço metodológico apresentado, SDMs continuam sendo simplificações da realidade ecológica. Muitos processos importantes podem não ser representados explicitamente, incluindo dispersão, demografia, interações bióticas, perturbações humanas, mudanças no uso da terra, plasticidade fenotípica, adaptação local e dinâmica populacional.

Perspectivas futuras incluem a integração de SDMs com modelos demográficos, dados genômicos, LiDAR, sensoriamento remoto, traços funcionais, inventários florestais permanentes, modelos de conectividade, cenários de desmatamento e planejamento sistemático da conservação. No contexto amazônico, essa integração é especialmente urgente diante das mudanças climáticas, fragmentação florestal e perda de biodiversidade.

\section{Mensagem final}

A Modelagem de Distribuição de Espécies é uma das ferramentas mais importantes da ecologia espacial contemporânea. Quando aplicada de forma crítica, ela permite transformar dados dispersos de ocorrência e ambiente em hipóteses espaciais úteis para ciência, ensino e conservação.

Contudo, a força de um SDM não está apenas no algoritmo, mas na coerência do processo científico. Um modelo bem construído deve ser ecologicamente fundamentado, estatisticamente avaliado, espacialmente explícito, reprodutível, transparente e interpretado com humildade diante das incertezas.

\begin{conceito}[title={\faLeaf\quad Consideração final}]
	Modelar a distribuição de uma espécie é, em última instância, tentar compreender como a vida ocupa o espaço sob as restrições do ambiente, da história e da mudança. Em um mundo em transformação, SDMs bem conduzidos podem ajudar a antecipar riscos, revelar oportunidades de conservação e orientar decisões mais responsáveis para a biodiversidade.
\end{conceito}

%================================================
\chapter*{Glossário}
\addcontentsline{toc}{chapter}{Glossário}
%================================================

\section*{Como utilizar este glossário}

Ao longo desta apostila foram apresentados conceitos provenientes de diversas
áreas do conhecimento, incluindo ecologia, biogeografia, climatologia,
estatística, ciência de dados espaciais e modelagem computacional.

Este glossário reúne os principais termos, siglas e conceitos utilizados em
Modelagem de Distribuição de Espécies (SDM), funcionando como uma referência
rápida para consulta durante a leitura, desenvolvimento de projetos e
interpretação dos resultados.

\begin{longtable}{
		p{0.18\textwidth}
		p{0.72\textwidth}
	}
	
	\caption{Glossário dos principais termos utilizados em Modelagem de Distribuição de Espécies.}
	\label{tab:glossario_sdm}
	\\
	
	\toprule
	\textbf{Termo} &
	\textbf{Definição}
	\\
	\midrule
	\endfirsthead
	
	\toprule
	\textbf{Termo} &
	\textbf{Definição}
	\\
	\midrule
	\endhead
	
	SDM &
	Species Distribution Model. Modelo utilizado para estimar a distribuição espacial de uma espécie em função de variáveis ambientais.
	\\
	
	ENM &
	Ecological Niche Model. Abordagem voltada à modelagem do nicho ecológico da espécie.
	\\
	
	HSM &
	Habitat Suitability Model. Modelo de adequabilidade de habitat.
	\\
	
	AUC &
	Area Under the Curve. Métrica baseada na curva ROC utilizada para avaliar a capacidade discriminatória do modelo.
	\\
	
	ROC &
	Receiver Operating Characteristic. Curva utilizada para avaliar o desempenho classificatório dos modelos.
	\\
	
	TSS &
	True Skill Statistic. Métrica que combina sensibilidade e especificidade.
	\\
	
	Kappa &
	Coeficiente de concordância utilizado para avaliar desempenho preditivo.
	\\
	
	Boyce Index &
	Métrica que avalia a concordância entre adequabilidade prevista e registros observados.
	\\
	
	VIF &
	Variance Inflation Factor. Indicador utilizado para detectar multicolinearidade entre variáveis preditoras.
	\\
	
	PCA &
	Principal Component Analysis. Técnica de redução de dimensionalidade baseada em componentes principais.
	\\
	
	GLM &
	Generalized Linear Model. Modelo linear generalizado amplamente utilizado em SDM.
	\\
	
	GAM &
	Generalized Additive Model. Extensão não linear dos modelos lineares generalizados.
	\\
	
	RF &
	Random Forest. Algoritmo baseado em múltiplas árvores de decisão.
	\\
	
	BRT &
	Boosted Regression Trees. Método baseado em árvores de regressão combinadas por boosting.
	\\
	
	Maxent &
	Maximum Entropy Modeling. Algoritmo amplamente utilizado para dados de presença apenas.
	\\
	
	Ensemble &
	Combinação de múltiplos modelos para produzir uma predição consensual.
	\\
	
	Background &
	Conjunto de condições ambientais utilizadas para representar a disponibilidade ambiental na área de calibração.
	\\
	
	Pseudo-ausência &
	Local sem registro da espécie utilizado para representar ausências em determinados algoritmos.
	\\
	
	Área M &
	Região acessível à espécie ao longo de sua história biogeográfica segundo o framework BAM.
	\\
	
	BAM &
	Framework conceitual que integra componentes Bióticos, Abióticos e de Movimento (dispersão).
	\\
	
	Nicho Fundamental &
	Conjunto de condições ambientais sob as quais a espécie pode sobreviver na ausência de restrições bióticas.
	\\
	
	Nicho Realizado &
	Porção do nicho fundamental efetivamente ocupada pela espécie.
	\\
	
	Adequabilidade Ambiental &
	Grau de compatibilidade entre as condições ambientais de um local e o nicho estimado da espécie.
	\\
	
	Transferibilidade &
	Capacidade do modelo manter desempenho quando aplicado a novas regiões ou períodos temporais.
	\\
	
	Extrapolação &
	Predição realizada sob condições ambientais não observadas durante a calibração.
	\\
	
	MESS &
	Multivariate Environmental Similarity Surface. Ferramenta para detectar extrapolação ambiental.
	\\
	
	MOP &
	Mobility-Oriented Parity. Método para identificar novidade ambiental em projeções espaciais ou temporais.
	\\
	
	CMIP6 &
	Coupled Model Intercomparison Project Phase 6. Conjunto internacional de modelos climáticos globais.
	\\
	
	GCM &
	Global Climate Model. Modelo climático global utilizado em projeções futuras.
	\\
	
	SSP &
	Shared Socioeconomic Pathway. Cenários socioeconômicos utilizados nas projeções climáticas.
	\\
	
	SSP126 &
	Cenário de baixas emissões e forte mitigação climática.
	\\
	
	SSP585 &
	Cenário de altas emissões e forte aquecimento global.
	\\
	
	LIG &
	Last Interglacial. Último Interglacial, aproximadamente 130 mil anos antes do presente.
	\\
	
	LGM &
	Last Glacial Maximum. Último Máximo Glacial, aproximadamente 21 mil anos antes do presente.
	\\
	
	Mid-Holocene &
	Holoceno Médio, aproximadamente 6 mil anos antes do presente.
	\\
	
	WorldClim &
	Base global de variáveis climáticas amplamente utilizada em SDM.
	\\
	
	CHELSA &
	Conjunto de variáveis climáticas de alta resolução derivadas de modelos atmosféricos.
	\\
	
	FAIR &
	Findable, Accessible, Interoperable and Reusable. Princípios internacionais para gestão de dados científicos.
	\\
	
	DOI &
	Digital Object Identifier. Identificador digital permanente para publicações e conjuntos de dados.
	\\
	
	Git &
	Sistema de controle de versão amplamente utilizado em ciência computacional.
	\\
	
	GitHub &
	Plataforma de hospedagem de código e colaboração científica baseada em Git.
	\\
	
	Zenodo &
	Repositório científico que permite arquivamento permanente e geração de DOI.
	\\
	
	Dryad &
	Repositório internacional para armazenamento e compartilhamento de dados científicos.
	\\
	
	Figshare &
	Plataforma de compartilhamento de dados, figuras e materiais suplementares.
	\\
	
	\bottomrule
	
\end{longtable}

\begin{conceito}
	\textbf{Observação.}
	
	A terminologia utilizada em SDM evolui continuamente à medida que novos
	algoritmos, bases de dados e abordagens metodológicas são desenvolvidos.
	Por esse motivo, recomenda-se que estudantes e pesquisadores consultem
	regularmente a literatura especializada.
\end{conceito}

%================================================
\chapter*{Leituras Essenciais}
\addcontentsline{toc}{chapter}{Leituras Essenciais}
%================================================

\section*{Como utilizar esta bibliografia}

A literatura sobre Modelagem de Distribuição de Espécies (Species Distribution
Modeling -- SDM) cresceu exponencialmente nas últimas décadas, acompanhando o
avanço da ecologia quantitativa, da biogeografia, da ciência de dados
ambientais e das mudanças climáticas.

Embora milhares de artigos tenham sido publicados sobre o tema, algumas obras
se destacam por estabelecer os fundamentos conceituais, metodológicos e
computacionais que estruturam a área atualmente.

Esta seleção reúne livros, artigos clássicos e trabalhos de referência que
servem como ponto de partida para estudantes, pesquisadores e profissionais
interessados em aprofundar seus conhecimentos em modelagem de nicho ecológico,
distribuição de espécies, mudanças climáticas e conservação da biodiversidade.

\section{Livros fundamentais}

\begin{table}[H]
	\centering
	\caption{Principais livros recomendados para aprofundamento em SDM.}
	\label{tab:livros_sdm}
	
	\small
	
	\begin{tabular}{
			p{0.30\textwidth}
			p{0.25\textwidth}
			p{0.35\textwidth}
		}
		
		\toprule
		
		\textbf{Livro} &
		\textbf{Autores} &
		\textbf{Contribuição}
		
		\\
		\midrule
		
		\textit{Mapping Species Distributions} &
		Franklin (2010) &
		Referência clássica sobre fundamentos ecológicos, estatísticos e espaciais da SDM.
		
		\\
		
		\textit{Habitat Suitability and Distribution Models} &
		Guisan, Thuiller \& Zimmermann (2017) &
		Obra moderna e abrangente sobre teoria, algoritmos e aplicações.
		
		\\
		
		\textit{Ecological Niches and Geographic Distributions} &
		Peterson et al. (2011) &
		Fundamentos conceituais do nicho ecológico e modelagem biogeográfica.
		
		\\
		
		\textit{Species Distribution Modeling with R} &
		Hijmans \& Elith &
		Aplicações práticas em ambiente R.
		
		\\
		
		\textit{Biogeography} &
		Lomolino et al. (2017) &
		Fundamentos biogeográficos necessários para interpretação dos modelos.
		
		\\
		
		\bottomrule
		
	\end{tabular}
	
\end{table}

\section{Artigos clássicos e fundamentais}

Os trabalhos apresentados a seguir são amplamente reconhecidos como marcos
conceituais e metodológicos da Modelagem de Distribuição de Espécies.

\begin{table}[H]
	\centering
	\caption{Artigos fundamentais para compreensão da evolução da SDM.}
	\label{tab:artigos_classicos_sdm}
	
	\small
	
	\begin{tabular}{
			p{0.30\textwidth}
			p{0.18\textwidth}
			p{0.40\textwidth}
		}
		
		\toprule
		
		\textbf{Referência} &
		\textbf{Ano} &
		\textbf{Contribuição}
		
		\\
		\midrule
		
		Guisan \& Zimmermann &
		2000 &
		Síntese pioneira sobre modelos preditivos de distribuição de espécies.
		
		\\
		
		Phillips, Anderson \& Schapire &
		2006 &
		Introdução formal do algoritmo Maxent.
		
		\\
		
		Elith et al. &
		2006 &
		Comparação abrangente de métodos modernos de SDM.
		
		\\
		
		Araújo \& New &
		2007 &
		Fundamentos da modelagem ensemble.
		
		\\
		
		Soberón \& Peterson &
		2005 &
		Integração entre nicho ecológico e distribuição geográfica.
		
		\\
		
		Barve et al. &
		2011 &
		Importância da definição da área acessível ($M$).
		
		\\
		
		Elith et al. &
		2010 &
		Avaliação da extrapolação ambiental.
		
		\\
		
		Zurell et al. &
		2020 &
		Protocolo internacional para relatórios e reprodutibilidade em SDM.
		
		\\
		
		\bottomrule
		
	\end{tabular}
	
\end{table}

\section{Leituras sobre mudanças climáticas e biodiversidade}

A integração entre SDM e mudanças climáticas tornou-se uma das principais áreas
de aplicação da modelagem ecológica. As obras abaixo são particularmente
relevantes para estudos de projeções futuras e conservação.

\begin{itemize}
	
	\item Thuiller et al. (2005) --- Climate change threats to plant diversity in Europe.
	
	\item Araújo et al. (2019) --- Standards for distribution models in climate change studies.
	
	\item IPCC (2021) --- Climate Change 2021: The Physical Science Basis.
	
	\item IPBES (2019) --- Global Assessment Report on Biodiversity and Ecosystem Services.
	
	\item Pecl et al. (2017) --- Biodiversity redistribution under climate change.
	
\end{itemize}

\section{Leituras sobre paleoclima e estabilidade climática}

A reconstrução de distribuições passadas representa uma importante aplicação da
SDM para investigação de refúgios climáticos, estabilidade ambiental e
história biogeográfica das espécies.

\begin{itemize}
	
	\item Svenning et al. (2011).
	
	\item Carnaval \& Moritz (2008).
	
	\item Graham et al. (2010).
	
	\item Peterson et al. (2011).
	
	\item Nogués-Bravo (2009).
	
\end{itemize}

\section{Leituras sobre reprodutibilidade científica}

A transparência e a reprodutibilidade tornaram-se componentes centrais da
pesquisa ecológica moderna. As referências abaixo fornecem excelentes pontos de
partida para implementação de boas práticas em projetos de SDM.

\begin{itemize}
	
	\item Peng (2011) --- Reproducible Research in Computational Science.
	
	\item Sandve et al. (2013) --- Ten Simple Rules for Reproducible Computational Research.
	
	\item Wilkinson et al. (2016) --- FAIR Guiding Principles.
	
	\item Zurell et al. (2020) --- Standard protocol for reporting species distribution models.
	
\end{itemize}

\section{Sequência recomendada de estudo}

Para leitores que desejam aprofundar seus conhecimentos em SDM, recomenda-se a
seguinte sequência de leitura:

\begin{enumerate}
	
	\item Guisan \& Zimmermann (2000);
	
	\item Franklin (2010);
	
	\item Peterson et al. (2011);
	
	\item Elith et al. (2006);
	
	\item Barve et al. (2011);
	
	\item Guisan, Thuiller \& Zimmermann (2017);
	
	\item Zurell et al. (2020).
	
\end{enumerate}

Essa sequência conduz o leitor desde os fundamentos ecológicos e
biogeográficos da modelagem até as práticas modernas de validação,
reprodutibilidade e aplicação em conservação sob mudanças climáticas.

\begin{dica}
	\textbf{Mensagem final ao leitor.}
	
	A Modelagem de Distribuição de Espécies é uma área em rápida evolução.
	Nenhum livro ou curso consegue esgotar completamente os conceitos,
	algoritmos e aplicações disponíveis.
	
	O desenvolvimento contínuo depende da leitura crítica da literatura
	científica, da prática com dados reais e da integração entre teoria
	ecológica, análise quantitativa e conservação da biodiversidade.
\end{dica}

\backmatter
\printbibliography[title={Referências}]

\end{document}
